\documentclass[11pt]{article}
\usepackage{amsmath,amssymb,amsthm,mathtools}
\usepackage[margin=1in]{geometry}
\usepackage{booktabs,longtable,array}
\usepackage{xcolor}
\usepackage[hidelinks]{hyperref}
\hypersetup{pdftitle={The singular spectrum of a sparse Mertens matrix},pdfauthor={Jeffery Kline},pdfsubject={Version 0.1.0, doi:10.5281/zenodo.23004980}}
\usepackage[expansion=false]{microtype}
\numberwithin{equation}{section}
\newtheorem{theorem}{Theorem}[section]
\newtheorem{proposition}[theorem]{Proposition}
\newtheorem{lemma}[theorem]{Lemma}
\newtheorem{corollary}[theorem]{Corollary}
\theoremstyle{definition}
\newtheorem{definition}[theorem]{Definition}
\theoremstyle{remark}
\newtheorem{remark}[theorem]{Remark}
\newcommand{\norm}[1]{\left\lVert#1\right\rVert}
\newcommand{\one}{\mathbf{1}}
\newcommand{\muvec}{\boldsymbol{\mu}}
\newcommand{\sfq}{\mathrm{sf}}
\DeclareMathOperator{\rank}{rank}
\DeclareMathOperator{\diag}{diag}
\DeclareMathOperator{\tr}{tr}
\DeclareMathOperator{\adj}{adj}
\DeclareMathOperator{\spanop}{span}
\DeclareMathOperator{\dist}{dist}
\title{The singular spectrum of a sparse Mertens matrix}
\author{Jeffery Kline}
\date{September 27, 2026\\[3pt]\small Version 0.1.0, \href{https://doi.org/10.5281/zenodo.23004980}{doi:10.5281/zenodo.23004980}}
\begin{document}
\maketitle
\begin{abstract}
We study the $n\times n$ matrix $B_n$ formed by placing ones on the diagonal and at $(i,i/P^+(i))$ for squarefree $i>1$, then replacing the first row by ones; all other entries are zero. Here $P^+(i)$ is the largest prime factor of $i$. Introduced by Kline (2019), this matrix satisfies $\det B_n=M(n)=\sum_{k\le n}\mu(k)$, where $\mu$ is the M\"obius function. We determine the multiplicity of the singular value one, compare upper singular values with square roots of parent degrees, and derive fixed-index lower limits from an explicit compact operator, including at singular matrices. The smallest singular value is exceptional: whenever $M(n)\ne0$,
\[
 \sigma_n(B_n)\sim\frac{|M(n)|}{\sqrt{Q(n)}\,W_n},\qquad
 W_n=n\exp\bigl(-(1+o(1))\sqrt{\log n\log\log n}\bigr),
\]
where $Q(n)$ counts the squarefree integers up to $n$ and $W_n$ is the norm of a vector of restricted M\"obius sums. This holds unconditionally. An earlier release by the author bounded $\sigma_n(B_n)$ between $|M(n)|n^{-3/2+o(1)}$ and $|M(n)|n^{-4/3+o(1)}$, and proved that it attains the lower end under a hypothesis implied by the Riemann hypothesis. We show that this hypothesis always holds. Classical estimates also give fixed-order moments and support-constrained sums of the restricted M\"obius sums, and products of singular values. Component projections give a geometric residual and a Dickman--Buchstab limit profile. The compact-Gram method and the restricted sums have prior arithmetic uses, and arithmetic priority for the moment formulations is not established. No stronger estimate for Mertens sums or prime counting is obtained: the formula for $\sigma_n(B_n)$ contains $M(n)$ and does not bound it.
\end{abstract}
\tableofcontents
\section{Introduction}

The determinant of a matrix can encode an arithmetic sum while leaving the geometry of most of its singular directions unexplained. The sparse Mertens matrix provides a particularly concrete example. Its nonzero entries come from a rooted tree on the squarefree integers, together with isolated nonsquarefree coordinates and a row of ones. Its determinant is the sum of the M\"obius function. This paper describes substantially more of its singular spectrum and separates the one direction carrying that sum from the remaining directions.

\paragraph{A six-integer example.}
An integer is squarefree if no prime factor occurs twice. Among $1,2,\ldots,6$, only $4=2^2$ fails this test. The primes $2,3,5$ each have parent $1$. The parent of $6=2\cdot3$ is $6/3=2$: we remove its largest prime factor, not an arbitrary prime factor. Draw arrows from parent to child:
\begin{center}
\setlength{\unitlength}{1pt}
\begin{picture}(240,78)
\put(20,36){\makebox(0,0){$1$}}
\put(30,41){\vector(2,1){42}}
\put(30,36){\vector(1,0){42}}
\put(30,31){\vector(2,-1){42}}
\put(82,62){\makebox(0,0){$2$}}
\put(82,36){\makebox(0,0){$3$}}
\put(82,10){\makebox(0,0){$5$}}
\put(92,62){\vector(1,0){42}}
\put(144,62){\makebox(0,0){$6$}}
\put(205,36){\makebox(0,0){$4$}}
\put(205,18){\makebox(0,0){\small isolated}}
\end{picture}
\end{center}
Rows and columns are ordered by their integer labels. Start with a one on every diagonal entry. For each arrow $j\to i$, put another one in row $i$, column $j$; leave all other entries zero. This gives $A_n$. Replace its first row by ones to obtain $B_n$. For example, the arrow $2\to6$ gives the entry in row six, column two. The isolated vertex $4$ keeps its diagonal one and is also included in the new first row. Thus the arrows describe the original parent links, not every nonzero entry of $B_n$.

\paragraph{Comparison with Redheffer's matrix.}
In the conventional orientation, Redheffer's matrix has entry one at $(i,j)$ when $j=1$ or $i\mid j$, and zero otherwise. Its transpose therefore has entry one when $i=1$ or $j\mid i$. Figure~\ref{fig:redheffer} uses this transpose so that both matrices have their row of ones at the top. Transposition preserves the determinant and singular values. Below that row, Redheffer's transpose includes every divisor of the row index; $B_n$ retains the diagonal and, for each squarefree index greater than one, just its parent. Both determinants equal $M(n)$. At $n=120$, the two matrices have respectively $721$ and $313$ nonzero entries.

\begin{figure}[t]
\centering
% Generated by build_comparison.py; exact entries, no raster image.
\begingroup\setlength{\unitlength}{1pt}
\definecolor{cmpRow}{HTML}{C27627}
\definecolor{cmpDiag}{HTML}{838B91}
\definecolor{cmpParent}{HTML}{007F7A}
\definecolor{cmpExtra}{HTML}{779CC3}
\begin{picture}(468,264)
\put(123.000,255.000){\makebox(0,0)[c]{\fontsize{10}{12}\selectfont Redheffer, transposed: $R_{120}^T$}}
\put(123.000,240.000){\makebox(0,0)[c]{\fontsize{9}{11}\selectfont 721 nonzero entries}}
\put(26.050,225.433){\color{cmpRow}\rule{1.517pt}{1.517pt}}
\put(27.667,225.433){\color{cmpRow}\rule{1.517pt}{1.517pt}}
\put(29.283,225.433){\color{cmpRow}\rule{1.517pt}{1.517pt}}
\put(30.900,225.433){\color{cmpRow}\rule{1.517pt}{1.517pt}}
\put(32.517,225.433){\color{cmpRow}\rule{1.517pt}{1.517pt}}
\put(34.133,225.433){\color{cmpRow}\rule{1.517pt}{1.517pt}}
\put(35.750,225.433){\color{cmpRow}\rule{1.517pt}{1.517pt}}
\put(37.367,225.433){\color{cmpRow}\rule{1.517pt}{1.517pt}}
\put(38.983,225.433){\color{cmpRow}\rule{1.517pt}{1.517pt}}
\put(40.600,225.433){\color{cmpRow}\rule{1.517pt}{1.517pt}}
\put(42.217,225.433){\color{cmpRow}\rule{1.517pt}{1.517pt}}
\put(43.833,225.433){\color{cmpRow}\rule{1.517pt}{1.517pt}}
\put(45.450,225.433){\color{cmpRow}\rule{1.517pt}{1.517pt}}
\put(47.067,225.433){\color{cmpRow}\rule{1.517pt}{1.517pt}}
\put(48.683,225.433){\color{cmpRow}\rule{1.517pt}{1.517pt}}
\put(50.300,225.433){\color{cmpRow}\rule{1.517pt}{1.517pt}}
\put(51.917,225.433){\color{cmpRow}\rule{1.517pt}{1.517pt}}
\put(53.533,225.433){\color{cmpRow}\rule{1.517pt}{1.517pt}}
\put(55.150,225.433){\color{cmpRow}\rule{1.517pt}{1.517pt}}
\put(56.767,225.433){\color{cmpRow}\rule{1.517pt}{1.517pt}}
\put(58.383,225.433){\color{cmpRow}\rule{1.517pt}{1.517pt}}
\put(60.000,225.433){\color{cmpRow}\rule{1.517pt}{1.517pt}}
\put(61.617,225.433){\color{cmpRow}\rule{1.517pt}{1.517pt}}
\put(63.233,225.433){\color{cmpRow}\rule{1.517pt}{1.517pt}}
\put(64.850,225.433){\color{cmpRow}\rule{1.517pt}{1.517pt}}
\put(66.467,225.433){\color{cmpRow}\rule{1.517pt}{1.517pt}}
\put(68.083,225.433){\color{cmpRow}\rule{1.517pt}{1.517pt}}
\put(69.700,225.433){\color{cmpRow}\rule{1.517pt}{1.517pt}}
\put(71.317,225.433){\color{cmpRow}\rule{1.517pt}{1.517pt}}
\put(72.933,225.433){\color{cmpRow}\rule{1.517pt}{1.517pt}}
\put(74.550,225.433){\color{cmpRow}\rule{1.517pt}{1.517pt}}
\put(76.167,225.433){\color{cmpRow}\rule{1.517pt}{1.517pt}}
\put(77.783,225.433){\color{cmpRow}\rule{1.517pt}{1.517pt}}
\put(79.400,225.433){\color{cmpRow}\rule{1.517pt}{1.517pt}}
\put(81.017,225.433){\color{cmpRow}\rule{1.517pt}{1.517pt}}
\put(82.633,225.433){\color{cmpRow}\rule{1.517pt}{1.517pt}}
\put(84.250,225.433){\color{cmpRow}\rule{1.517pt}{1.517pt}}
\put(85.867,225.433){\color{cmpRow}\rule{1.517pt}{1.517pt}}
\put(87.483,225.433){\color{cmpRow}\rule{1.517pt}{1.517pt}}
\put(89.100,225.433){\color{cmpRow}\rule{1.517pt}{1.517pt}}
\put(90.717,225.433){\color{cmpRow}\rule{1.517pt}{1.517pt}}
\put(92.333,225.433){\color{cmpRow}\rule{1.517pt}{1.517pt}}
\put(93.950,225.433){\color{cmpRow}\rule{1.517pt}{1.517pt}}
\put(95.567,225.433){\color{cmpRow}\rule{1.517pt}{1.517pt}}
\put(97.183,225.433){\color{cmpRow}\rule{1.517pt}{1.517pt}}
\put(98.800,225.433){\color{cmpRow}\rule{1.517pt}{1.517pt}}
\put(100.417,225.433){\color{cmpRow}\rule{1.517pt}{1.517pt}}
\put(102.033,225.433){\color{cmpRow}\rule{1.517pt}{1.517pt}}
\put(103.650,225.433){\color{cmpRow}\rule{1.517pt}{1.517pt}}
\put(105.267,225.433){\color{cmpRow}\rule{1.517pt}{1.517pt}}
\put(106.883,225.433){\color{cmpRow}\rule{1.517pt}{1.517pt}}
\put(108.500,225.433){\color{cmpRow}\rule{1.517pt}{1.517pt}}
\put(110.117,225.433){\color{cmpRow}\rule{1.517pt}{1.517pt}}
\put(111.733,225.433){\color{cmpRow}\rule{1.517pt}{1.517pt}}
\put(113.350,225.433){\color{cmpRow}\rule{1.517pt}{1.517pt}}
\put(114.967,225.433){\color{cmpRow}\rule{1.517pt}{1.517pt}}
\put(116.583,225.433){\color{cmpRow}\rule{1.517pt}{1.517pt}}
\put(118.200,225.433){\color{cmpRow}\rule{1.517pt}{1.517pt}}
\put(119.817,225.433){\color{cmpRow}\rule{1.517pt}{1.517pt}}
\put(121.433,225.433){\color{cmpRow}\rule{1.517pt}{1.517pt}}
\put(123.050,225.433){\color{cmpRow}\rule{1.517pt}{1.517pt}}
\put(124.667,225.433){\color{cmpRow}\rule{1.517pt}{1.517pt}}
\put(126.283,225.433){\color{cmpRow}\rule{1.517pt}{1.517pt}}
\put(127.900,225.433){\color{cmpRow}\rule{1.517pt}{1.517pt}}
\put(129.517,225.433){\color{cmpRow}\rule{1.517pt}{1.517pt}}
\put(131.133,225.433){\color{cmpRow}\rule{1.517pt}{1.517pt}}
\put(132.750,225.433){\color{cmpRow}\rule{1.517pt}{1.517pt}}
\put(134.367,225.433){\color{cmpRow}\rule{1.517pt}{1.517pt}}
\put(135.983,225.433){\color{cmpRow}\rule{1.517pt}{1.517pt}}
\put(137.600,225.433){\color{cmpRow}\rule{1.517pt}{1.517pt}}
\put(139.217,225.433){\color{cmpRow}\rule{1.517pt}{1.517pt}}
\put(140.833,225.433){\color{cmpRow}\rule{1.517pt}{1.517pt}}
\put(142.450,225.433){\color{cmpRow}\rule{1.517pt}{1.517pt}}
\put(144.067,225.433){\color{cmpRow}\rule{1.517pt}{1.517pt}}
\put(145.683,225.433){\color{cmpRow}\rule{1.517pt}{1.517pt}}
\put(147.300,225.433){\color{cmpRow}\rule{1.517pt}{1.517pt}}
\put(148.917,225.433){\color{cmpRow}\rule{1.517pt}{1.517pt}}
\put(150.533,225.433){\color{cmpRow}\rule{1.517pt}{1.517pt}}
\put(152.150,225.433){\color{cmpRow}\rule{1.517pt}{1.517pt}}
\put(153.767,225.433){\color{cmpRow}\rule{1.517pt}{1.517pt}}
\put(155.383,225.433){\color{cmpRow}\rule{1.517pt}{1.517pt}}
\put(157.000,225.433){\color{cmpRow}\rule{1.517pt}{1.517pt}}
\put(158.617,225.433){\color{cmpRow}\rule{1.517pt}{1.517pt}}
\put(160.233,225.433){\color{cmpRow}\rule{1.517pt}{1.517pt}}
\put(161.850,225.433){\color{cmpRow}\rule{1.517pt}{1.517pt}}
\put(163.467,225.433){\color{cmpRow}\rule{1.517pt}{1.517pt}}
\put(165.083,225.433){\color{cmpRow}\rule{1.517pt}{1.517pt}}
\put(166.700,225.433){\color{cmpRow}\rule{1.517pt}{1.517pt}}
\put(168.317,225.433){\color{cmpRow}\rule{1.517pt}{1.517pt}}
\put(169.933,225.433){\color{cmpRow}\rule{1.517pt}{1.517pt}}
\put(171.550,225.433){\color{cmpRow}\rule{1.517pt}{1.517pt}}
\put(173.167,225.433){\color{cmpRow}\rule{1.517pt}{1.517pt}}
\put(174.783,225.433){\color{cmpRow}\rule{1.517pt}{1.517pt}}
\put(176.400,225.433){\color{cmpRow}\rule{1.517pt}{1.517pt}}
\put(178.017,225.433){\color{cmpRow}\rule{1.517pt}{1.517pt}}
\put(179.633,225.433){\color{cmpRow}\rule{1.517pt}{1.517pt}}
\put(181.250,225.433){\color{cmpRow}\rule{1.517pt}{1.517pt}}
\put(182.867,225.433){\color{cmpRow}\rule{1.517pt}{1.517pt}}
\put(184.483,225.433){\color{cmpRow}\rule{1.517pt}{1.517pt}}
\put(186.100,225.433){\color{cmpRow}\rule{1.517pt}{1.517pt}}
\put(187.717,225.433){\color{cmpRow}\rule{1.517pt}{1.517pt}}
\put(189.333,225.433){\color{cmpRow}\rule{1.517pt}{1.517pt}}
\put(190.950,225.433){\color{cmpRow}\rule{1.517pt}{1.517pt}}
\put(192.567,225.433){\color{cmpRow}\rule{1.517pt}{1.517pt}}
\put(194.183,225.433){\color{cmpRow}\rule{1.517pt}{1.517pt}}
\put(195.800,225.433){\color{cmpRow}\rule{1.517pt}{1.517pt}}
\put(197.417,225.433){\color{cmpRow}\rule{1.517pt}{1.517pt}}
\put(199.033,225.433){\color{cmpRow}\rule{1.517pt}{1.517pt}}
\put(200.650,225.433){\color{cmpRow}\rule{1.517pt}{1.517pt}}
\put(202.267,225.433){\color{cmpRow}\rule{1.517pt}{1.517pt}}
\put(203.883,225.433){\color{cmpRow}\rule{1.517pt}{1.517pt}}
\put(205.500,225.433){\color{cmpRow}\rule{1.517pt}{1.517pt}}
\put(207.117,225.433){\color{cmpRow}\rule{1.517pt}{1.517pt}}
\put(208.733,225.433){\color{cmpRow}\rule{1.517pt}{1.517pt}}
\put(210.350,225.433){\color{cmpRow}\rule{1.517pt}{1.517pt}}
\put(211.967,225.433){\color{cmpRow}\rule{1.517pt}{1.517pt}}
\put(213.583,225.433){\color{cmpRow}\rule{1.517pt}{1.517pt}}
\put(215.200,225.433){\color{cmpRow}\rule{1.517pt}{1.517pt}}
\put(216.817,225.433){\color{cmpRow}\rule{1.517pt}{1.517pt}}
\put(218.433,225.433){\color{cmpRow}\rule{1.517pt}{1.517pt}}
\put(26.050,223.817){\color{cmpParent}\rule{1.517pt}{1.517pt}}
\put(27.667,223.817){\color{cmpDiag}\rule{1.517pt}{1.517pt}}
\put(26.050,222.200){\color{cmpParent}\rule{1.517pt}{1.517pt}}
\put(29.283,222.200){\color{cmpDiag}\rule{1.517pt}{1.517pt}}
\put(26.050,220.583){\color{cmpExtra}\rule{1.517pt}{1.517pt}}
\put(27.667,220.583){\color{cmpExtra}\rule{1.517pt}{1.517pt}}
\put(30.900,220.583){\color{cmpDiag}\rule{1.517pt}{1.517pt}}
\put(26.050,218.967){\color{cmpParent}\rule{1.517pt}{1.517pt}}
\put(32.517,218.967){\color{cmpDiag}\rule{1.517pt}{1.517pt}}
\put(26.050,217.350){\color{cmpExtra}\rule{1.517pt}{1.517pt}}
\put(27.667,217.350){\color{cmpParent}\rule{1.517pt}{1.517pt}}
\put(29.283,217.350){\color{cmpExtra}\rule{1.517pt}{1.517pt}}
\put(34.133,217.350){\color{cmpDiag}\rule{1.517pt}{1.517pt}}
\put(26.050,215.733){\color{cmpParent}\rule{1.517pt}{1.517pt}}
\put(35.750,215.733){\color{cmpDiag}\rule{1.517pt}{1.517pt}}
\put(26.050,214.117){\color{cmpExtra}\rule{1.517pt}{1.517pt}}
\put(27.667,214.117){\color{cmpExtra}\rule{1.517pt}{1.517pt}}
\put(30.900,214.117){\color{cmpExtra}\rule{1.517pt}{1.517pt}}
\put(37.367,214.117){\color{cmpDiag}\rule{1.517pt}{1.517pt}}
\put(26.050,212.500){\color{cmpExtra}\rule{1.517pt}{1.517pt}}
\put(29.283,212.500){\color{cmpExtra}\rule{1.517pt}{1.517pt}}
\put(38.983,212.500){\color{cmpDiag}\rule{1.517pt}{1.517pt}}
\put(26.050,210.883){\color{cmpExtra}\rule{1.517pt}{1.517pt}}
\put(27.667,210.883){\color{cmpParent}\rule{1.517pt}{1.517pt}}
\put(32.517,210.883){\color{cmpExtra}\rule{1.517pt}{1.517pt}}
\put(40.600,210.883){\color{cmpDiag}\rule{1.517pt}{1.517pt}}
\put(26.050,209.267){\color{cmpParent}\rule{1.517pt}{1.517pt}}
\put(42.217,209.267){\color{cmpDiag}\rule{1.517pt}{1.517pt}}
\put(26.050,207.650){\color{cmpExtra}\rule{1.517pt}{1.517pt}}
\put(27.667,207.650){\color{cmpExtra}\rule{1.517pt}{1.517pt}}
\put(29.283,207.650){\color{cmpExtra}\rule{1.517pt}{1.517pt}}
\put(30.900,207.650){\color{cmpExtra}\rule{1.517pt}{1.517pt}}
\put(34.133,207.650){\color{cmpExtra}\rule{1.517pt}{1.517pt}}
\put(43.833,207.650){\color{cmpDiag}\rule{1.517pt}{1.517pt}}
\put(26.050,206.033){\color{cmpParent}\rule{1.517pt}{1.517pt}}
\put(45.450,206.033){\color{cmpDiag}\rule{1.517pt}{1.517pt}}
\put(26.050,204.417){\color{cmpExtra}\rule{1.517pt}{1.517pt}}
\put(27.667,204.417){\color{cmpParent}\rule{1.517pt}{1.517pt}}
\put(35.750,204.417){\color{cmpExtra}\rule{1.517pt}{1.517pt}}
\put(47.067,204.417){\color{cmpDiag}\rule{1.517pt}{1.517pt}}
\put(26.050,202.800){\color{cmpExtra}\rule{1.517pt}{1.517pt}}
\put(29.283,202.800){\color{cmpParent}\rule{1.517pt}{1.517pt}}
\put(32.517,202.800){\color{cmpExtra}\rule{1.517pt}{1.517pt}}
\put(48.683,202.800){\color{cmpDiag}\rule{1.517pt}{1.517pt}}
\put(26.050,201.183){\color{cmpExtra}\rule{1.517pt}{1.517pt}}
\put(27.667,201.183){\color{cmpExtra}\rule{1.517pt}{1.517pt}}
\put(30.900,201.183){\color{cmpExtra}\rule{1.517pt}{1.517pt}}
\put(37.367,201.183){\color{cmpExtra}\rule{1.517pt}{1.517pt}}
\put(50.300,201.183){\color{cmpDiag}\rule{1.517pt}{1.517pt}}
\put(26.050,199.567){\color{cmpParent}\rule{1.517pt}{1.517pt}}
\put(51.917,199.567){\color{cmpDiag}\rule{1.517pt}{1.517pt}}
\put(26.050,197.950){\color{cmpExtra}\rule{1.517pt}{1.517pt}}
\put(27.667,197.950){\color{cmpExtra}\rule{1.517pt}{1.517pt}}
\put(29.283,197.950){\color{cmpExtra}\rule{1.517pt}{1.517pt}}
\put(34.133,197.950){\color{cmpExtra}\rule{1.517pt}{1.517pt}}
\put(38.983,197.950){\color{cmpExtra}\rule{1.517pt}{1.517pt}}
\put(53.533,197.950){\color{cmpDiag}\rule{1.517pt}{1.517pt}}
\put(26.050,196.333){\color{cmpParent}\rule{1.517pt}{1.517pt}}
\put(55.150,196.333){\color{cmpDiag}\rule{1.517pt}{1.517pt}}
\put(26.050,194.717){\color{cmpExtra}\rule{1.517pt}{1.517pt}}
\put(27.667,194.717){\color{cmpExtra}\rule{1.517pt}{1.517pt}}
\put(30.900,194.717){\color{cmpExtra}\rule{1.517pt}{1.517pt}}
\put(32.517,194.717){\color{cmpExtra}\rule{1.517pt}{1.517pt}}
\put(40.600,194.717){\color{cmpExtra}\rule{1.517pt}{1.517pt}}
\put(56.767,194.717){\color{cmpDiag}\rule{1.517pt}{1.517pt}}
\put(26.050,193.100){\color{cmpExtra}\rule{1.517pt}{1.517pt}}
\put(29.283,193.100){\color{cmpParent}\rule{1.517pt}{1.517pt}}
\put(35.750,193.100){\color{cmpExtra}\rule{1.517pt}{1.517pt}}
\put(58.383,193.100){\color{cmpDiag}\rule{1.517pt}{1.517pt}}
\put(26.050,191.483){\color{cmpExtra}\rule{1.517pt}{1.517pt}}
\put(27.667,191.483){\color{cmpParent}\rule{1.517pt}{1.517pt}}
\put(42.217,191.483){\color{cmpExtra}\rule{1.517pt}{1.517pt}}
\put(60.000,191.483){\color{cmpDiag}\rule{1.517pt}{1.517pt}}
\put(26.050,189.867){\color{cmpParent}\rule{1.517pt}{1.517pt}}
\put(61.617,189.867){\color{cmpDiag}\rule{1.517pt}{1.517pt}}
\put(26.050,188.250){\color{cmpExtra}\rule{1.517pt}{1.517pt}}
\put(27.667,188.250){\color{cmpExtra}\rule{1.517pt}{1.517pt}}
\put(29.283,188.250){\color{cmpExtra}\rule{1.517pt}{1.517pt}}
\put(30.900,188.250){\color{cmpExtra}\rule{1.517pt}{1.517pt}}
\put(34.133,188.250){\color{cmpExtra}\rule{1.517pt}{1.517pt}}
\put(37.367,188.250){\color{cmpExtra}\rule{1.517pt}{1.517pt}}
\put(43.833,188.250){\color{cmpExtra}\rule{1.517pt}{1.517pt}}
\put(63.233,188.250){\color{cmpDiag}\rule{1.517pt}{1.517pt}}
\put(26.050,186.633){\color{cmpExtra}\rule{1.517pt}{1.517pt}}
\put(32.517,186.633){\color{cmpExtra}\rule{1.517pt}{1.517pt}}
\put(64.850,186.633){\color{cmpDiag}\rule{1.517pt}{1.517pt}}
\put(26.050,185.017){\color{cmpExtra}\rule{1.517pt}{1.517pt}}
\put(27.667,185.017){\color{cmpParent}\rule{1.517pt}{1.517pt}}
\put(45.450,185.017){\color{cmpExtra}\rule{1.517pt}{1.517pt}}
\put(66.467,185.017){\color{cmpDiag}\rule{1.517pt}{1.517pt}}
\put(26.050,183.400){\color{cmpExtra}\rule{1.517pt}{1.517pt}}
\put(29.283,183.400){\color{cmpExtra}\rule{1.517pt}{1.517pt}}
\put(38.983,183.400){\color{cmpExtra}\rule{1.517pt}{1.517pt}}
\put(68.083,183.400){\color{cmpDiag}\rule{1.517pt}{1.517pt}}
\put(26.050,181.783){\color{cmpExtra}\rule{1.517pt}{1.517pt}}
\put(27.667,181.783){\color{cmpExtra}\rule{1.517pt}{1.517pt}}
\put(30.900,181.783){\color{cmpExtra}\rule{1.517pt}{1.517pt}}
\put(35.750,181.783){\color{cmpExtra}\rule{1.517pt}{1.517pt}}
\put(47.067,181.783){\color{cmpExtra}\rule{1.517pt}{1.517pt}}
\put(69.700,181.783){\color{cmpDiag}\rule{1.517pt}{1.517pt}}
\put(26.050,180.167){\color{cmpParent}\rule{1.517pt}{1.517pt}}
\put(71.317,180.167){\color{cmpDiag}\rule{1.517pt}{1.517pt}}
\put(26.050,178.550){\color{cmpExtra}\rule{1.517pt}{1.517pt}}
\put(27.667,178.550){\color{cmpExtra}\rule{1.517pt}{1.517pt}}
\put(29.283,178.550){\color{cmpExtra}\rule{1.517pt}{1.517pt}}
\put(32.517,178.550){\color{cmpExtra}\rule{1.517pt}{1.517pt}}
\put(34.133,178.550){\color{cmpParent}\rule{1.517pt}{1.517pt}}
\put(40.600,178.550){\color{cmpExtra}\rule{1.517pt}{1.517pt}}
\put(48.683,178.550){\color{cmpExtra}\rule{1.517pt}{1.517pt}}
\put(72.933,178.550){\color{cmpDiag}\rule{1.517pt}{1.517pt}}
\put(26.050,176.933){\color{cmpParent}\rule{1.517pt}{1.517pt}}
\put(74.550,176.933){\color{cmpDiag}\rule{1.517pt}{1.517pt}}
\put(26.050,175.317){\color{cmpExtra}\rule{1.517pt}{1.517pt}}
\put(27.667,175.317){\color{cmpExtra}\rule{1.517pt}{1.517pt}}
\put(30.900,175.317){\color{cmpExtra}\rule{1.517pt}{1.517pt}}
\put(37.367,175.317){\color{cmpExtra}\rule{1.517pt}{1.517pt}}
\put(50.300,175.317){\color{cmpExtra}\rule{1.517pt}{1.517pt}}
\put(76.167,175.317){\color{cmpDiag}\rule{1.517pt}{1.517pt}}
\put(26.050,173.700){\color{cmpExtra}\rule{1.517pt}{1.517pt}}
\put(29.283,173.700){\color{cmpParent}\rule{1.517pt}{1.517pt}}
\put(42.217,173.700){\color{cmpExtra}\rule{1.517pt}{1.517pt}}
\put(77.783,173.700){\color{cmpDiag}\rule{1.517pt}{1.517pt}}
\put(26.050,172.083){\color{cmpExtra}\rule{1.517pt}{1.517pt}}
\put(27.667,172.083){\color{cmpParent}\rule{1.517pt}{1.517pt}}
\put(51.917,172.083){\color{cmpExtra}\rule{1.517pt}{1.517pt}}
\put(79.400,172.083){\color{cmpDiag}\rule{1.517pt}{1.517pt}}
\put(26.050,170.467){\color{cmpExtra}\rule{1.517pt}{1.517pt}}
\put(32.517,170.467){\color{cmpParent}\rule{1.517pt}{1.517pt}}
\put(35.750,170.467){\color{cmpExtra}\rule{1.517pt}{1.517pt}}
\put(81.017,170.467){\color{cmpDiag}\rule{1.517pt}{1.517pt}}
\put(26.050,168.850){\color{cmpExtra}\rule{1.517pt}{1.517pt}}
\put(27.667,168.850){\color{cmpExtra}\rule{1.517pt}{1.517pt}}
\put(29.283,168.850){\color{cmpExtra}\rule{1.517pt}{1.517pt}}
\put(30.900,168.850){\color{cmpExtra}\rule{1.517pt}{1.517pt}}
\put(34.133,168.850){\color{cmpExtra}\rule{1.517pt}{1.517pt}}
\put(38.983,168.850){\color{cmpExtra}\rule{1.517pt}{1.517pt}}
\put(43.833,168.850){\color{cmpExtra}\rule{1.517pt}{1.517pt}}
\put(53.533,168.850){\color{cmpExtra}\rule{1.517pt}{1.517pt}}
\put(82.633,168.850){\color{cmpDiag}\rule{1.517pt}{1.517pt}}
\put(26.050,167.233){\color{cmpParent}\rule{1.517pt}{1.517pt}}
\put(84.250,167.233){\color{cmpDiag}\rule{1.517pt}{1.517pt}}
\put(26.050,165.617){\color{cmpExtra}\rule{1.517pt}{1.517pt}}
\put(27.667,165.617){\color{cmpParent}\rule{1.517pt}{1.517pt}}
\put(55.150,165.617){\color{cmpExtra}\rule{1.517pt}{1.517pt}}
\put(85.867,165.617){\color{cmpDiag}\rule{1.517pt}{1.517pt}}
\put(26.050,164.000){\color{cmpExtra}\rule{1.517pt}{1.517pt}}
\put(29.283,164.000){\color{cmpParent}\rule{1.517pt}{1.517pt}}
\put(45.450,164.000){\color{cmpExtra}\rule{1.517pt}{1.517pt}}
\put(87.483,164.000){\color{cmpDiag}\rule{1.517pt}{1.517pt}}
\put(26.050,162.383){\color{cmpExtra}\rule{1.517pt}{1.517pt}}
\put(27.667,162.383){\color{cmpExtra}\rule{1.517pt}{1.517pt}}
\put(30.900,162.383){\color{cmpExtra}\rule{1.517pt}{1.517pt}}
\put(32.517,162.383){\color{cmpExtra}\rule{1.517pt}{1.517pt}}
\put(37.367,162.383){\color{cmpExtra}\rule{1.517pt}{1.517pt}}
\put(40.600,162.383){\color{cmpExtra}\rule{1.517pt}{1.517pt}}
\put(56.767,162.383){\color{cmpExtra}\rule{1.517pt}{1.517pt}}
\put(89.100,162.383){\color{cmpDiag}\rule{1.517pt}{1.517pt}}
\put(26.050,160.767){\color{cmpParent}\rule{1.517pt}{1.517pt}}
\put(90.717,160.767){\color{cmpDiag}\rule{1.517pt}{1.517pt}}
\put(26.050,159.150){\color{cmpExtra}\rule{1.517pt}{1.517pt}}
\put(27.667,159.150){\color{cmpExtra}\rule{1.517pt}{1.517pt}}
\put(29.283,159.150){\color{cmpExtra}\rule{1.517pt}{1.517pt}}
\put(34.133,159.150){\color{cmpParent}\rule{1.517pt}{1.517pt}}
\put(35.750,159.150){\color{cmpExtra}\rule{1.517pt}{1.517pt}}
\put(47.067,159.150){\color{cmpExtra}\rule{1.517pt}{1.517pt}}
\put(58.383,159.150){\color{cmpExtra}\rule{1.517pt}{1.517pt}}
\put(92.333,159.150){\color{cmpDiag}\rule{1.517pt}{1.517pt}}
\put(26.050,157.533){\color{cmpParent}\rule{1.517pt}{1.517pt}}
\put(93.950,157.533){\color{cmpDiag}\rule{1.517pt}{1.517pt}}
\put(26.050,155.917){\color{cmpExtra}\rule{1.517pt}{1.517pt}}
\put(27.667,155.917){\color{cmpExtra}\rule{1.517pt}{1.517pt}}
\put(30.900,155.917){\color{cmpExtra}\rule{1.517pt}{1.517pt}}
\put(42.217,155.917){\color{cmpExtra}\rule{1.517pt}{1.517pt}}
\put(60.000,155.917){\color{cmpExtra}\rule{1.517pt}{1.517pt}}
\put(95.567,155.917){\color{cmpDiag}\rule{1.517pt}{1.517pt}}
\put(26.050,154.300){\color{cmpExtra}\rule{1.517pt}{1.517pt}}
\put(29.283,154.300){\color{cmpExtra}\rule{1.517pt}{1.517pt}}
\put(32.517,154.300){\color{cmpExtra}\rule{1.517pt}{1.517pt}}
\put(38.983,154.300){\color{cmpExtra}\rule{1.517pt}{1.517pt}}
\put(48.683,154.300){\color{cmpExtra}\rule{1.517pt}{1.517pt}}
\put(97.183,154.300){\color{cmpDiag}\rule{1.517pt}{1.517pt}}
\put(26.050,152.683){\color{cmpExtra}\rule{1.517pt}{1.517pt}}
\put(27.667,152.683){\color{cmpParent}\rule{1.517pt}{1.517pt}}
\put(61.617,152.683){\color{cmpExtra}\rule{1.517pt}{1.517pt}}
\put(98.800,152.683){\color{cmpDiag}\rule{1.517pt}{1.517pt}}
\put(26.050,151.067){\color{cmpParent}\rule{1.517pt}{1.517pt}}
\put(100.417,151.067){\color{cmpDiag}\rule{1.517pt}{1.517pt}}
\put(26.050,149.450){\color{cmpExtra}\rule{1.517pt}{1.517pt}}
\put(27.667,149.450){\color{cmpExtra}\rule{1.517pt}{1.517pt}}
\put(29.283,149.450){\color{cmpExtra}\rule{1.517pt}{1.517pt}}
\put(30.900,149.450){\color{cmpExtra}\rule{1.517pt}{1.517pt}}
\put(34.133,149.450){\color{cmpExtra}\rule{1.517pt}{1.517pt}}
\put(37.367,149.450){\color{cmpExtra}\rule{1.517pt}{1.517pt}}
\put(43.833,149.450){\color{cmpExtra}\rule{1.517pt}{1.517pt}}
\put(50.300,149.450){\color{cmpExtra}\rule{1.517pt}{1.517pt}}
\put(63.233,149.450){\color{cmpExtra}\rule{1.517pt}{1.517pt}}
\put(102.033,149.450){\color{cmpDiag}\rule{1.517pt}{1.517pt}}
\put(26.050,147.833){\color{cmpExtra}\rule{1.517pt}{1.517pt}}
\put(35.750,147.833){\color{cmpExtra}\rule{1.517pt}{1.517pt}}
\put(103.650,147.833){\color{cmpDiag}\rule{1.517pt}{1.517pt}}
\put(26.050,146.217){\color{cmpExtra}\rule{1.517pt}{1.517pt}}
\put(27.667,146.217){\color{cmpExtra}\rule{1.517pt}{1.517pt}}
\put(32.517,146.217){\color{cmpExtra}\rule{1.517pt}{1.517pt}}
\put(40.600,146.217){\color{cmpExtra}\rule{1.517pt}{1.517pt}}
\put(64.850,146.217){\color{cmpExtra}\rule{1.517pt}{1.517pt}}
\put(105.267,146.217){\color{cmpDiag}\rule{1.517pt}{1.517pt}}
\put(26.050,144.600){\color{cmpExtra}\rule{1.517pt}{1.517pt}}
\put(29.283,144.600){\color{cmpParent}\rule{1.517pt}{1.517pt}}
\put(51.917,144.600){\color{cmpExtra}\rule{1.517pt}{1.517pt}}
\put(106.883,144.600){\color{cmpDiag}\rule{1.517pt}{1.517pt}}
\put(26.050,142.983){\color{cmpExtra}\rule{1.517pt}{1.517pt}}
\put(27.667,142.983){\color{cmpExtra}\rule{1.517pt}{1.517pt}}
\put(30.900,142.983){\color{cmpExtra}\rule{1.517pt}{1.517pt}}
\put(45.450,142.983){\color{cmpExtra}\rule{1.517pt}{1.517pt}}
\put(66.467,142.983){\color{cmpExtra}\rule{1.517pt}{1.517pt}}
\put(108.500,142.983){\color{cmpDiag}\rule{1.517pt}{1.517pt}}
\put(26.050,141.367){\color{cmpParent}\rule{1.517pt}{1.517pt}}
\put(110.117,141.367){\color{cmpDiag}\rule{1.517pt}{1.517pt}}
\put(26.050,139.750){\color{cmpExtra}\rule{1.517pt}{1.517pt}}
\put(27.667,139.750){\color{cmpExtra}\rule{1.517pt}{1.517pt}}
\put(29.283,139.750){\color{cmpExtra}\rule{1.517pt}{1.517pt}}
\put(34.133,139.750){\color{cmpExtra}\rule{1.517pt}{1.517pt}}
\put(38.983,139.750){\color{cmpExtra}\rule{1.517pt}{1.517pt}}
\put(53.533,139.750){\color{cmpExtra}\rule{1.517pt}{1.517pt}}
\put(68.083,139.750){\color{cmpExtra}\rule{1.517pt}{1.517pt}}
\put(111.733,139.750){\color{cmpDiag}\rule{1.517pt}{1.517pt}}
\put(26.050,138.133){\color{cmpExtra}\rule{1.517pt}{1.517pt}}
\put(32.517,138.133){\color{cmpParent}\rule{1.517pt}{1.517pt}}
\put(42.217,138.133){\color{cmpExtra}\rule{1.517pt}{1.517pt}}
\put(113.350,138.133){\color{cmpDiag}\rule{1.517pt}{1.517pt}}
\put(26.050,136.517){\color{cmpExtra}\rule{1.517pt}{1.517pt}}
\put(27.667,136.517){\color{cmpExtra}\rule{1.517pt}{1.517pt}}
\put(30.900,136.517){\color{cmpExtra}\rule{1.517pt}{1.517pt}}
\put(35.750,136.517){\color{cmpExtra}\rule{1.517pt}{1.517pt}}
\put(37.367,136.517){\color{cmpExtra}\rule{1.517pt}{1.517pt}}
\put(47.067,136.517){\color{cmpExtra}\rule{1.517pt}{1.517pt}}
\put(69.700,136.517){\color{cmpExtra}\rule{1.517pt}{1.517pt}}
\put(114.967,136.517){\color{cmpDiag}\rule{1.517pt}{1.517pt}}
\put(26.050,134.900){\color{cmpExtra}\rule{1.517pt}{1.517pt}}
\put(29.283,134.900){\color{cmpParent}\rule{1.517pt}{1.517pt}}
\put(55.150,134.900){\color{cmpExtra}\rule{1.517pt}{1.517pt}}
\put(116.583,134.900){\color{cmpDiag}\rule{1.517pt}{1.517pt}}
\put(26.050,133.283){\color{cmpExtra}\rule{1.517pt}{1.517pt}}
\put(27.667,133.283){\color{cmpParent}\rule{1.517pt}{1.517pt}}
\put(71.317,133.283){\color{cmpExtra}\rule{1.517pt}{1.517pt}}
\put(118.200,133.283){\color{cmpDiag}\rule{1.517pt}{1.517pt}}
\put(26.050,131.667){\color{cmpParent}\rule{1.517pt}{1.517pt}}
\put(119.817,131.667){\color{cmpDiag}\rule{1.517pt}{1.517pt}}
\put(26.050,130.050){\color{cmpExtra}\rule{1.517pt}{1.517pt}}
\put(27.667,130.050){\color{cmpExtra}\rule{1.517pt}{1.517pt}}
\put(29.283,130.050){\color{cmpExtra}\rule{1.517pt}{1.517pt}}
\put(30.900,130.050){\color{cmpExtra}\rule{1.517pt}{1.517pt}}
\put(32.517,130.050){\color{cmpExtra}\rule{1.517pt}{1.517pt}}
\put(34.133,130.050){\color{cmpExtra}\rule{1.517pt}{1.517pt}}
\put(40.600,130.050){\color{cmpExtra}\rule{1.517pt}{1.517pt}}
\put(43.833,130.050){\color{cmpExtra}\rule{1.517pt}{1.517pt}}
\put(48.683,130.050){\color{cmpExtra}\rule{1.517pt}{1.517pt}}
\put(56.767,130.050){\color{cmpExtra}\rule{1.517pt}{1.517pt}}
\put(72.933,130.050){\color{cmpExtra}\rule{1.517pt}{1.517pt}}
\put(121.433,130.050){\color{cmpDiag}\rule{1.517pt}{1.517pt}}
\put(26.050,128.433){\color{cmpParent}\rule{1.517pt}{1.517pt}}
\put(123.050,128.433){\color{cmpDiag}\rule{1.517pt}{1.517pt}}
\put(26.050,126.817){\color{cmpExtra}\rule{1.517pt}{1.517pt}}
\put(27.667,126.817){\color{cmpParent}\rule{1.517pt}{1.517pt}}
\put(74.550,126.817){\color{cmpExtra}\rule{1.517pt}{1.517pt}}
\put(124.667,126.817){\color{cmpDiag}\rule{1.517pt}{1.517pt}}
\put(26.050,125.200){\color{cmpExtra}\rule{1.517pt}{1.517pt}}
\put(29.283,125.200){\color{cmpExtra}\rule{1.517pt}{1.517pt}}
\put(35.750,125.200){\color{cmpExtra}\rule{1.517pt}{1.517pt}}
\put(38.983,125.200){\color{cmpExtra}\rule{1.517pt}{1.517pt}}
\put(58.383,125.200){\color{cmpExtra}\rule{1.517pt}{1.517pt}}
\put(126.283,125.200){\color{cmpDiag}\rule{1.517pt}{1.517pt}}
\put(26.050,123.583){\color{cmpExtra}\rule{1.517pt}{1.517pt}}
\put(27.667,123.583){\color{cmpExtra}\rule{1.517pt}{1.517pt}}
\put(30.900,123.583){\color{cmpExtra}\rule{1.517pt}{1.517pt}}
\put(37.367,123.583){\color{cmpExtra}\rule{1.517pt}{1.517pt}}
\put(50.300,123.583){\color{cmpExtra}\rule{1.517pt}{1.517pt}}
\put(76.167,123.583){\color{cmpExtra}\rule{1.517pt}{1.517pt}}
\put(127.900,123.583){\color{cmpDiag}\rule{1.517pt}{1.517pt}}
\put(26.050,121.967){\color{cmpExtra}\rule{1.517pt}{1.517pt}}
\put(32.517,121.967){\color{cmpParent}\rule{1.517pt}{1.517pt}}
\put(45.450,121.967){\color{cmpExtra}\rule{1.517pt}{1.517pt}}
\put(129.517,121.967){\color{cmpDiag}\rule{1.517pt}{1.517pt}}
\put(26.050,120.350){\color{cmpExtra}\rule{1.517pt}{1.517pt}}
\put(27.667,120.350){\color{cmpExtra}\rule{1.517pt}{1.517pt}}
\put(29.283,120.350){\color{cmpExtra}\rule{1.517pt}{1.517pt}}
\put(34.133,120.350){\color{cmpParent}\rule{1.517pt}{1.517pt}}
\put(42.217,120.350){\color{cmpExtra}\rule{1.517pt}{1.517pt}}
\put(60.000,120.350){\color{cmpExtra}\rule{1.517pt}{1.517pt}}
\put(77.783,120.350){\color{cmpExtra}\rule{1.517pt}{1.517pt}}
\put(131.133,120.350){\color{cmpDiag}\rule{1.517pt}{1.517pt}}
\put(26.050,118.733){\color{cmpParent}\rule{1.517pt}{1.517pt}}
\put(132.750,118.733){\color{cmpDiag}\rule{1.517pt}{1.517pt}}
\put(26.050,117.117){\color{cmpExtra}\rule{1.517pt}{1.517pt}}
\put(27.667,117.117){\color{cmpExtra}\rule{1.517pt}{1.517pt}}
\put(30.900,117.117){\color{cmpExtra}\rule{1.517pt}{1.517pt}}
\put(51.917,117.117){\color{cmpExtra}\rule{1.517pt}{1.517pt}}
\put(79.400,117.117){\color{cmpExtra}\rule{1.517pt}{1.517pt}}
\put(134.367,117.117){\color{cmpDiag}\rule{1.517pt}{1.517pt}}
\put(26.050,115.500){\color{cmpExtra}\rule{1.517pt}{1.517pt}}
\put(29.283,115.500){\color{cmpParent}\rule{1.517pt}{1.517pt}}
\put(61.617,115.500){\color{cmpExtra}\rule{1.517pt}{1.517pt}}
\put(135.983,115.500){\color{cmpDiag}\rule{1.517pt}{1.517pt}}
\put(26.050,113.883){\color{cmpExtra}\rule{1.517pt}{1.517pt}}
\put(27.667,113.883){\color{cmpExtra}\rule{1.517pt}{1.517pt}}
\put(32.517,113.883){\color{cmpExtra}\rule{1.517pt}{1.517pt}}
\put(35.750,113.883){\color{cmpExtra}\rule{1.517pt}{1.517pt}}
\put(40.600,113.883){\color{cmpParent}\rule{1.517pt}{1.517pt}}
\put(47.067,113.883){\color{cmpExtra}\rule{1.517pt}{1.517pt}}
\put(81.017,113.883){\color{cmpExtra}\rule{1.517pt}{1.517pt}}
\put(137.600,113.883){\color{cmpDiag}\rule{1.517pt}{1.517pt}}
\put(26.050,112.267){\color{cmpParent}\rule{1.517pt}{1.517pt}}
\put(139.217,112.267){\color{cmpDiag}\rule{1.517pt}{1.517pt}}
\put(26.050,110.650){\color{cmpExtra}\rule{1.517pt}{1.517pt}}
\put(27.667,110.650){\color{cmpExtra}\rule{1.517pt}{1.517pt}}
\put(29.283,110.650){\color{cmpExtra}\rule{1.517pt}{1.517pt}}
\put(30.900,110.650){\color{cmpExtra}\rule{1.517pt}{1.517pt}}
\put(34.133,110.650){\color{cmpExtra}\rule{1.517pt}{1.517pt}}
\put(37.367,110.650){\color{cmpExtra}\rule{1.517pt}{1.517pt}}
\put(38.983,110.650){\color{cmpExtra}\rule{1.517pt}{1.517pt}}
\put(43.833,110.650){\color{cmpExtra}\rule{1.517pt}{1.517pt}}
\put(53.533,110.650){\color{cmpExtra}\rule{1.517pt}{1.517pt}}
\put(63.233,110.650){\color{cmpExtra}\rule{1.517pt}{1.517pt}}
\put(82.633,110.650){\color{cmpExtra}\rule{1.517pt}{1.517pt}}
\put(140.833,110.650){\color{cmpDiag}\rule{1.517pt}{1.517pt}}
\put(26.050,109.033){\color{cmpParent}\rule{1.517pt}{1.517pt}}
\put(142.450,109.033){\color{cmpDiag}\rule{1.517pt}{1.517pt}}
\put(26.050,107.417){\color{cmpExtra}\rule{1.517pt}{1.517pt}}
\put(27.667,107.417){\color{cmpParent}\rule{1.517pt}{1.517pt}}
\put(84.250,107.417){\color{cmpExtra}\rule{1.517pt}{1.517pt}}
\put(144.067,107.417){\color{cmpDiag}\rule{1.517pt}{1.517pt}}
\put(26.050,105.800){\color{cmpExtra}\rule{1.517pt}{1.517pt}}
\put(29.283,105.800){\color{cmpExtra}\rule{1.517pt}{1.517pt}}
\put(32.517,105.800){\color{cmpExtra}\rule{1.517pt}{1.517pt}}
\put(48.683,105.800){\color{cmpExtra}\rule{1.517pt}{1.517pt}}
\put(64.850,105.800){\color{cmpExtra}\rule{1.517pt}{1.517pt}}
\put(145.683,105.800){\color{cmpDiag}\rule{1.517pt}{1.517pt}}
\put(26.050,104.183){\color{cmpExtra}\rule{1.517pt}{1.517pt}}
\put(27.667,104.183){\color{cmpExtra}\rule{1.517pt}{1.517pt}}
\put(30.900,104.183){\color{cmpExtra}\rule{1.517pt}{1.517pt}}
\put(55.150,104.183){\color{cmpExtra}\rule{1.517pt}{1.517pt}}
\put(85.867,104.183){\color{cmpExtra}\rule{1.517pt}{1.517pt}}
\put(147.300,104.183){\color{cmpDiag}\rule{1.517pt}{1.517pt}}
\put(26.050,102.567){\color{cmpExtra}\rule{1.517pt}{1.517pt}}
\put(35.750,102.567){\color{cmpParent}\rule{1.517pt}{1.517pt}}
\put(42.217,102.567){\color{cmpExtra}\rule{1.517pt}{1.517pt}}
\put(148.917,102.567){\color{cmpDiag}\rule{1.517pt}{1.517pt}}
\put(26.050,100.950){\color{cmpExtra}\rule{1.517pt}{1.517pt}}
\put(27.667,100.950){\color{cmpExtra}\rule{1.517pt}{1.517pt}}
\put(29.283,100.950){\color{cmpExtra}\rule{1.517pt}{1.517pt}}
\put(34.133,100.950){\color{cmpParent}\rule{1.517pt}{1.517pt}}
\put(45.450,100.950){\color{cmpExtra}\rule{1.517pt}{1.517pt}}
\put(66.467,100.950){\color{cmpExtra}\rule{1.517pt}{1.517pt}}
\put(87.483,100.950){\color{cmpExtra}\rule{1.517pt}{1.517pt}}
\put(150.533,100.950){\color{cmpDiag}\rule{1.517pt}{1.517pt}}
\put(26.050,99.333){\color{cmpParent}\rule{1.517pt}{1.517pt}}
\put(152.150,99.333){\color{cmpDiag}\rule{1.517pt}{1.517pt}}
\put(26.050,97.717){\color{cmpExtra}\rule{1.517pt}{1.517pt}}
\put(27.667,97.717){\color{cmpExtra}\rule{1.517pt}{1.517pt}}
\put(30.900,97.717){\color{cmpExtra}\rule{1.517pt}{1.517pt}}
\put(32.517,97.717){\color{cmpExtra}\rule{1.517pt}{1.517pt}}
\put(37.367,97.717){\color{cmpExtra}\rule{1.517pt}{1.517pt}}
\put(40.600,97.717){\color{cmpExtra}\rule{1.517pt}{1.517pt}}
\put(50.300,97.717){\color{cmpExtra}\rule{1.517pt}{1.517pt}}
\put(56.767,97.717){\color{cmpExtra}\rule{1.517pt}{1.517pt}}
\put(89.100,97.717){\color{cmpExtra}\rule{1.517pt}{1.517pt}}
\put(153.767,97.717){\color{cmpDiag}\rule{1.517pt}{1.517pt}}
\put(26.050,96.100){\color{cmpExtra}\rule{1.517pt}{1.517pt}}
\put(29.283,96.100){\color{cmpExtra}\rule{1.517pt}{1.517pt}}
\put(38.983,96.100){\color{cmpExtra}\rule{1.517pt}{1.517pt}}
\put(68.083,96.100){\color{cmpExtra}\rule{1.517pt}{1.517pt}}
\put(155.383,96.100){\color{cmpDiag}\rule{1.517pt}{1.517pt}}
\put(26.050,94.483){\color{cmpExtra}\rule{1.517pt}{1.517pt}}
\put(27.667,94.483){\color{cmpParent}\rule{1.517pt}{1.517pt}}
\put(90.717,94.483){\color{cmpExtra}\rule{1.517pt}{1.517pt}}
\put(157.000,94.483){\color{cmpDiag}\rule{1.517pt}{1.517pt}}
\put(26.050,92.867){\color{cmpParent}\rule{1.517pt}{1.517pt}}
\put(158.617,92.867){\color{cmpDiag}\rule{1.517pt}{1.517pt}}
\put(26.050,91.250){\color{cmpExtra}\rule{1.517pt}{1.517pt}}
\put(27.667,91.250){\color{cmpExtra}\rule{1.517pt}{1.517pt}}
\put(29.283,91.250){\color{cmpExtra}\rule{1.517pt}{1.517pt}}
\put(30.900,91.250){\color{cmpExtra}\rule{1.517pt}{1.517pt}}
\put(34.133,91.250){\color{cmpExtra}\rule{1.517pt}{1.517pt}}
\put(35.750,91.250){\color{cmpExtra}\rule{1.517pt}{1.517pt}}
\put(43.833,91.250){\color{cmpExtra}\rule{1.517pt}{1.517pt}}
\put(47.067,91.250){\color{cmpExtra}\rule{1.517pt}{1.517pt}}
\put(58.383,91.250){\color{cmpExtra}\rule{1.517pt}{1.517pt}}
\put(69.700,91.250){\color{cmpExtra}\rule{1.517pt}{1.517pt}}
\put(92.333,91.250){\color{cmpExtra}\rule{1.517pt}{1.517pt}}
\put(160.233,91.250){\color{cmpDiag}\rule{1.517pt}{1.517pt}}
\put(26.050,89.633){\color{cmpExtra}\rule{1.517pt}{1.517pt}}
\put(32.517,89.633){\color{cmpParent}\rule{1.517pt}{1.517pt}}
\put(51.917,89.633){\color{cmpExtra}\rule{1.517pt}{1.517pt}}
\put(161.850,89.633){\color{cmpDiag}\rule{1.517pt}{1.517pt}}
\put(26.050,88.017){\color{cmpExtra}\rule{1.517pt}{1.517pt}}
\put(27.667,88.017){\color{cmpParent}\rule{1.517pt}{1.517pt}}
\put(93.950,88.017){\color{cmpExtra}\rule{1.517pt}{1.517pt}}
\put(163.467,88.017){\color{cmpDiag}\rule{1.517pt}{1.517pt}}
\put(26.050,86.400){\color{cmpExtra}\rule{1.517pt}{1.517pt}}
\put(29.283,86.400){\color{cmpParent}\rule{1.517pt}{1.517pt}}
\put(71.317,86.400){\color{cmpExtra}\rule{1.517pt}{1.517pt}}
\put(165.083,86.400){\color{cmpDiag}\rule{1.517pt}{1.517pt}}
\put(26.050,84.783){\color{cmpExtra}\rule{1.517pt}{1.517pt}}
\put(27.667,84.783){\color{cmpExtra}\rule{1.517pt}{1.517pt}}
\put(30.900,84.783){\color{cmpExtra}\rule{1.517pt}{1.517pt}}
\put(37.367,84.783){\color{cmpExtra}\rule{1.517pt}{1.517pt}}
\put(42.217,84.783){\color{cmpExtra}\rule{1.517pt}{1.517pt}}
\put(60.000,84.783){\color{cmpExtra}\rule{1.517pt}{1.517pt}}
\put(95.567,84.783){\color{cmpExtra}\rule{1.517pt}{1.517pt}}
\put(166.700,84.783){\color{cmpDiag}\rule{1.517pt}{1.517pt}}
\put(26.050,83.167){\color{cmpParent}\rule{1.517pt}{1.517pt}}
\put(168.317,83.167){\color{cmpDiag}\rule{1.517pt}{1.517pt}}
\put(26.050,81.550){\color{cmpExtra}\rule{1.517pt}{1.517pt}}
\put(27.667,81.550){\color{cmpExtra}\rule{1.517pt}{1.517pt}}
\put(29.283,81.550){\color{cmpExtra}\rule{1.517pt}{1.517pt}}
\put(32.517,81.550){\color{cmpExtra}\rule{1.517pt}{1.517pt}}
\put(34.133,81.550){\color{cmpExtra}\rule{1.517pt}{1.517pt}}
\put(38.983,81.550){\color{cmpExtra}\rule{1.517pt}{1.517pt}}
\put(40.600,81.550){\color{cmpExtra}\rule{1.517pt}{1.517pt}}
\put(48.683,81.550){\color{cmpExtra}\rule{1.517pt}{1.517pt}}
\put(53.533,81.550){\color{cmpExtra}\rule{1.517pt}{1.517pt}}
\put(72.933,81.550){\color{cmpExtra}\rule{1.517pt}{1.517pt}}
\put(97.183,81.550){\color{cmpExtra}\rule{1.517pt}{1.517pt}}
\put(169.933,81.550){\color{cmpDiag}\rule{1.517pt}{1.517pt}}
\put(26.050,79.933){\color{cmpExtra}\rule{1.517pt}{1.517pt}}
\put(35.750,79.933){\color{cmpParent}\rule{1.517pt}{1.517pt}}
\put(45.450,79.933){\color{cmpExtra}\rule{1.517pt}{1.517pt}}
\put(171.550,79.933){\color{cmpDiag}\rule{1.517pt}{1.517pt}}
\put(26.050,78.317){\color{cmpExtra}\rule{1.517pt}{1.517pt}}
\put(27.667,78.317){\color{cmpExtra}\rule{1.517pt}{1.517pt}}
\put(30.900,78.317){\color{cmpExtra}\rule{1.517pt}{1.517pt}}
\put(61.617,78.317){\color{cmpExtra}\rule{1.517pt}{1.517pt}}
\put(98.800,78.317){\color{cmpExtra}\rule{1.517pt}{1.517pt}}
\put(173.167,78.317){\color{cmpDiag}\rule{1.517pt}{1.517pt}}
\put(26.050,76.700){\color{cmpExtra}\rule{1.517pt}{1.517pt}}
\put(29.283,76.700){\color{cmpParent}\rule{1.517pt}{1.517pt}}
\put(74.550,76.700){\color{cmpExtra}\rule{1.517pt}{1.517pt}}
\put(174.783,76.700){\color{cmpDiag}\rule{1.517pt}{1.517pt}}
\put(26.050,75.083){\color{cmpExtra}\rule{1.517pt}{1.517pt}}
\put(27.667,75.083){\color{cmpParent}\rule{1.517pt}{1.517pt}}
\put(100.417,75.083){\color{cmpExtra}\rule{1.517pt}{1.517pt}}
\put(176.400,75.083){\color{cmpDiag}\rule{1.517pt}{1.517pt}}
\put(26.050,73.467){\color{cmpExtra}\rule{1.517pt}{1.517pt}}
\put(32.517,73.467){\color{cmpParent}\rule{1.517pt}{1.517pt}}
\put(55.150,73.467){\color{cmpExtra}\rule{1.517pt}{1.517pt}}
\put(178.017,73.467){\color{cmpDiag}\rule{1.517pt}{1.517pt}}
\put(26.050,71.850){\color{cmpExtra}\rule{1.517pt}{1.517pt}}
\put(27.667,71.850){\color{cmpExtra}\rule{1.517pt}{1.517pt}}
\put(29.283,71.850){\color{cmpExtra}\rule{1.517pt}{1.517pt}}
\put(30.900,71.850){\color{cmpExtra}\rule{1.517pt}{1.517pt}}
\put(34.133,71.850){\color{cmpExtra}\rule{1.517pt}{1.517pt}}
\put(37.367,71.850){\color{cmpExtra}\rule{1.517pt}{1.517pt}}
\put(43.833,71.850){\color{cmpExtra}\rule{1.517pt}{1.517pt}}
\put(50.300,71.850){\color{cmpExtra}\rule{1.517pt}{1.517pt}}
\put(63.233,71.850){\color{cmpExtra}\rule{1.517pt}{1.517pt}}
\put(76.167,71.850){\color{cmpExtra}\rule{1.517pt}{1.517pt}}
\put(102.033,71.850){\color{cmpExtra}\rule{1.517pt}{1.517pt}}
\put(179.633,71.850){\color{cmpDiag}\rule{1.517pt}{1.517pt}}
\put(26.050,70.233){\color{cmpParent}\rule{1.517pt}{1.517pt}}
\put(181.250,70.233){\color{cmpDiag}\rule{1.517pt}{1.517pt}}
\put(26.050,68.617){\color{cmpExtra}\rule{1.517pt}{1.517pt}}
\put(27.667,68.617){\color{cmpExtra}\rule{1.517pt}{1.517pt}}
\put(35.750,68.617){\color{cmpExtra}\rule{1.517pt}{1.517pt}}
\put(47.067,68.617){\color{cmpExtra}\rule{1.517pt}{1.517pt}}
\put(103.650,68.617){\color{cmpExtra}\rule{1.517pt}{1.517pt}}
\put(182.867,68.617){\color{cmpDiag}\rule{1.517pt}{1.517pt}}
\put(26.050,67.000){\color{cmpExtra}\rule{1.517pt}{1.517pt}}
\put(29.283,67.000){\color{cmpExtra}\rule{1.517pt}{1.517pt}}
\put(38.983,67.000){\color{cmpExtra}\rule{1.517pt}{1.517pt}}
\put(42.217,67.000){\color{cmpExtra}\rule{1.517pt}{1.517pt}}
\put(77.783,67.000){\color{cmpExtra}\rule{1.517pt}{1.517pt}}
\put(184.483,67.000){\color{cmpDiag}\rule{1.517pt}{1.517pt}}
\put(26.050,65.383){\color{cmpExtra}\rule{1.517pt}{1.517pt}}
\put(27.667,65.383){\color{cmpExtra}\rule{1.517pt}{1.517pt}}
\put(30.900,65.383){\color{cmpExtra}\rule{1.517pt}{1.517pt}}
\put(32.517,65.383){\color{cmpExtra}\rule{1.517pt}{1.517pt}}
\put(40.600,65.383){\color{cmpExtra}\rule{1.517pt}{1.517pt}}
\put(56.767,65.383){\color{cmpExtra}\rule{1.517pt}{1.517pt}}
\put(64.850,65.383){\color{cmpExtra}\rule{1.517pt}{1.517pt}}
\put(105.267,65.383){\color{cmpExtra}\rule{1.517pt}{1.517pt}}
\put(186.100,65.383){\color{cmpDiag}\rule{1.517pt}{1.517pt}}
\put(26.050,63.767){\color{cmpParent}\rule{1.517pt}{1.517pt}}
\put(187.717,63.767){\color{cmpDiag}\rule{1.517pt}{1.517pt}}
\put(26.050,62.150){\color{cmpExtra}\rule{1.517pt}{1.517pt}}
\put(27.667,62.150){\color{cmpExtra}\rule{1.517pt}{1.517pt}}
\put(29.283,62.150){\color{cmpExtra}\rule{1.517pt}{1.517pt}}
\put(34.133,62.150){\color{cmpParent}\rule{1.517pt}{1.517pt}}
\put(51.917,62.150){\color{cmpExtra}\rule{1.517pt}{1.517pt}}
\put(79.400,62.150){\color{cmpExtra}\rule{1.517pt}{1.517pt}}
\put(106.883,62.150){\color{cmpExtra}\rule{1.517pt}{1.517pt}}
\put(189.333,62.150){\color{cmpDiag}\rule{1.517pt}{1.517pt}}
\put(26.050,60.533){\color{cmpParent}\rule{1.517pt}{1.517pt}}
\put(190.950,60.533){\color{cmpDiag}\rule{1.517pt}{1.517pt}}
\put(26.050,58.917){\color{cmpExtra}\rule{1.517pt}{1.517pt}}
\put(27.667,58.917){\color{cmpExtra}\rule{1.517pt}{1.517pt}}
\put(30.900,58.917){\color{cmpExtra}\rule{1.517pt}{1.517pt}}
\put(37.367,58.917){\color{cmpExtra}\rule{1.517pt}{1.517pt}}
\put(45.450,58.917){\color{cmpExtra}\rule{1.517pt}{1.517pt}}
\put(66.467,58.917){\color{cmpExtra}\rule{1.517pt}{1.517pt}}
\put(108.500,58.917){\color{cmpExtra}\rule{1.517pt}{1.517pt}}
\put(192.567,58.917){\color{cmpDiag}\rule{1.517pt}{1.517pt}}
\put(26.050,57.300){\color{cmpExtra}\rule{1.517pt}{1.517pt}}
\put(29.283,57.300){\color{cmpExtra}\rule{1.517pt}{1.517pt}}
\put(32.517,57.300){\color{cmpExtra}\rule{1.517pt}{1.517pt}}
\put(35.750,57.300){\color{cmpExtra}\rule{1.517pt}{1.517pt}}
\put(48.683,57.300){\color{cmpParent}\rule{1.517pt}{1.517pt}}
\put(58.383,57.300){\color{cmpExtra}\rule{1.517pt}{1.517pt}}
\put(81.017,57.300){\color{cmpExtra}\rule{1.517pt}{1.517pt}}
\put(194.183,57.300){\color{cmpDiag}\rule{1.517pt}{1.517pt}}
\put(26.050,55.683){\color{cmpExtra}\rule{1.517pt}{1.517pt}}
\put(27.667,55.683){\color{cmpParent}\rule{1.517pt}{1.517pt}}
\put(110.117,55.683){\color{cmpExtra}\rule{1.517pt}{1.517pt}}
\put(195.800,55.683){\color{cmpDiag}\rule{1.517pt}{1.517pt}}
\put(26.050,54.067){\color{cmpParent}\rule{1.517pt}{1.517pt}}
\put(197.417,54.067){\color{cmpDiag}\rule{1.517pt}{1.517pt}}
\put(26.050,52.450){\color{cmpExtra}\rule{1.517pt}{1.517pt}}
\put(27.667,52.450){\color{cmpExtra}\rule{1.517pt}{1.517pt}}
\put(29.283,52.450){\color{cmpExtra}\rule{1.517pt}{1.517pt}}
\put(30.900,52.450){\color{cmpExtra}\rule{1.517pt}{1.517pt}}
\put(34.133,52.450){\color{cmpExtra}\rule{1.517pt}{1.517pt}}
\put(38.983,52.450){\color{cmpExtra}\rule{1.517pt}{1.517pt}}
\put(43.833,52.450){\color{cmpExtra}\rule{1.517pt}{1.517pt}}
\put(53.533,52.450){\color{cmpExtra}\rule{1.517pt}{1.517pt}}
\put(68.083,52.450){\color{cmpExtra}\rule{1.517pt}{1.517pt}}
\put(82.633,52.450){\color{cmpExtra}\rule{1.517pt}{1.517pt}}
\put(111.733,52.450){\color{cmpExtra}\rule{1.517pt}{1.517pt}}
\put(199.033,52.450){\color{cmpDiag}\rule{1.517pt}{1.517pt}}
\put(26.050,50.833){\color{cmpParent}\rule{1.517pt}{1.517pt}}
\put(200.650,50.833){\color{cmpDiag}\rule{1.517pt}{1.517pt}}
\put(26.050,49.217){\color{cmpExtra}\rule{1.517pt}{1.517pt}}
\put(27.667,49.217){\color{cmpExtra}\rule{1.517pt}{1.517pt}}
\put(32.517,49.217){\color{cmpExtra}\rule{1.517pt}{1.517pt}}
\put(40.600,49.217){\color{cmpParent}\rule{1.517pt}{1.517pt}}
\put(42.217,49.217){\color{cmpExtra}\rule{1.517pt}{1.517pt}}
\put(60.000,49.217){\color{cmpExtra}\rule{1.517pt}{1.517pt}}
\put(113.350,49.217){\color{cmpExtra}\rule{1.517pt}{1.517pt}}
\put(202.267,49.217){\color{cmpDiag}\rule{1.517pt}{1.517pt}}
\put(26.050,47.600){\color{cmpExtra}\rule{1.517pt}{1.517pt}}
\put(29.283,47.600){\color{cmpParent}\rule{1.517pt}{1.517pt}}
\put(84.250,47.600){\color{cmpExtra}\rule{1.517pt}{1.517pt}}
\put(203.883,47.600){\color{cmpDiag}\rule{1.517pt}{1.517pt}}
\put(26.050,45.983){\color{cmpExtra}\rule{1.517pt}{1.517pt}}
\put(27.667,45.983){\color{cmpExtra}\rule{1.517pt}{1.517pt}}
\put(30.900,45.983){\color{cmpExtra}\rule{1.517pt}{1.517pt}}
\put(35.750,45.983){\color{cmpExtra}\rule{1.517pt}{1.517pt}}
\put(37.367,45.983){\color{cmpExtra}\rule{1.517pt}{1.517pt}}
\put(47.067,45.983){\color{cmpExtra}\rule{1.517pt}{1.517pt}}
\put(50.300,45.983){\color{cmpExtra}\rule{1.517pt}{1.517pt}}
\put(69.700,45.983){\color{cmpExtra}\rule{1.517pt}{1.517pt}}
\put(114.967,45.983){\color{cmpExtra}\rule{1.517pt}{1.517pt}}
\put(205.500,45.983){\color{cmpDiag}\rule{1.517pt}{1.517pt}}
\put(26.050,44.367){\color{cmpParent}\rule{1.517pt}{1.517pt}}
\put(207.117,44.367){\color{cmpDiag}\rule{1.517pt}{1.517pt}}
\put(26.050,42.750){\color{cmpExtra}\rule{1.517pt}{1.517pt}}
\put(27.667,42.750){\color{cmpExtra}\rule{1.517pt}{1.517pt}}
\put(29.283,42.750){\color{cmpExtra}\rule{1.517pt}{1.517pt}}
\put(34.133,42.750){\color{cmpParent}\rule{1.517pt}{1.517pt}}
\put(55.150,42.750){\color{cmpExtra}\rule{1.517pt}{1.517pt}}
\put(85.867,42.750){\color{cmpExtra}\rule{1.517pt}{1.517pt}}
\put(116.583,42.750){\color{cmpExtra}\rule{1.517pt}{1.517pt}}
\put(208.733,42.750){\color{cmpDiag}\rule{1.517pt}{1.517pt}}
\put(26.050,41.133){\color{cmpExtra}\rule{1.517pt}{1.517pt}}
\put(32.517,41.133){\color{cmpParent}\rule{1.517pt}{1.517pt}}
\put(61.617,41.133){\color{cmpExtra}\rule{1.517pt}{1.517pt}}
\put(210.350,41.133){\color{cmpDiag}\rule{1.517pt}{1.517pt}}
\put(26.050,39.517){\color{cmpExtra}\rule{1.517pt}{1.517pt}}
\put(27.667,39.517){\color{cmpExtra}\rule{1.517pt}{1.517pt}}
\put(30.900,39.517){\color{cmpExtra}\rule{1.517pt}{1.517pt}}
\put(71.317,39.517){\color{cmpExtra}\rule{1.517pt}{1.517pt}}
\put(118.200,39.517){\color{cmpExtra}\rule{1.517pt}{1.517pt}}
\put(211.967,39.517){\color{cmpDiag}\rule{1.517pt}{1.517pt}}
\put(26.050,37.900){\color{cmpExtra}\rule{1.517pt}{1.517pt}}
\put(29.283,37.900){\color{cmpExtra}\rule{1.517pt}{1.517pt}}
\put(38.983,37.900){\color{cmpExtra}\rule{1.517pt}{1.517pt}}
\put(45.450,37.900){\color{cmpExtra}\rule{1.517pt}{1.517pt}}
\put(87.483,37.900){\color{cmpExtra}\rule{1.517pt}{1.517pt}}
\put(213.583,37.900){\color{cmpDiag}\rule{1.517pt}{1.517pt}}
\put(26.050,36.283){\color{cmpExtra}\rule{1.517pt}{1.517pt}}
\put(27.667,36.283){\color{cmpParent}\rule{1.517pt}{1.517pt}}
\put(119.817,36.283){\color{cmpExtra}\rule{1.517pt}{1.517pt}}
\put(215.200,36.283){\color{cmpDiag}\rule{1.517pt}{1.517pt}}
\put(26.050,34.667){\color{cmpExtra}\rule{1.517pt}{1.517pt}}
\put(35.750,34.667){\color{cmpParent}\rule{1.517pt}{1.517pt}}
\put(51.917,34.667){\color{cmpExtra}\rule{1.517pt}{1.517pt}}
\put(216.817,34.667){\color{cmpDiag}\rule{1.517pt}{1.517pt}}
\put(26.050,33.050){\color{cmpExtra}\rule{1.517pt}{1.517pt}}
\put(27.667,33.050){\color{cmpExtra}\rule{1.517pt}{1.517pt}}
\put(29.283,33.050){\color{cmpExtra}\rule{1.517pt}{1.517pt}}
\put(30.900,33.050){\color{cmpExtra}\rule{1.517pt}{1.517pt}}
\put(32.517,33.050){\color{cmpExtra}\rule{1.517pt}{1.517pt}}
\put(34.133,33.050){\color{cmpExtra}\rule{1.517pt}{1.517pt}}
\put(37.367,33.050){\color{cmpExtra}\rule{1.517pt}{1.517pt}}
\put(40.600,33.050){\color{cmpExtra}\rule{1.517pt}{1.517pt}}
\put(43.833,33.050){\color{cmpExtra}\rule{1.517pt}{1.517pt}}
\put(48.683,33.050){\color{cmpExtra}\rule{1.517pt}{1.517pt}}
\put(56.767,33.050){\color{cmpExtra}\rule{1.517pt}{1.517pt}}
\put(63.233,33.050){\color{cmpExtra}\rule{1.517pt}{1.517pt}}
\put(72.933,33.050){\color{cmpExtra}\rule{1.517pt}{1.517pt}}
\put(89.100,33.050){\color{cmpExtra}\rule{1.517pt}{1.517pt}}
\put(121.433,33.050){\color{cmpExtra}\rule{1.517pt}{1.517pt}}
\put(218.433,33.050){\color{cmpDiag}\rule{1.517pt}{1.517pt}}
\put(26.808,23.000){\makebox(0,0)[c]{\fontsize{8}{10}\selectfont 1}}
\put(21.000,226.192){\makebox(0,0)[r]{\fontsize{8}{10}\selectfont 1}}
\put(73.692,23.000){\makebox(0,0)[c]{\fontsize{8}{10}\selectfont 30}}
\put(21.000,179.308){\makebox(0,0)[r]{\fontsize{8}{10}\selectfont 30}}
\put(122.192,23.000){\makebox(0,0)[c]{\fontsize{8}{10}\selectfont 60}}
\put(21.000,130.808){\makebox(0,0)[r]{\fontsize{8}{10}\selectfont 60}}
\put(170.692,23.000){\makebox(0,0)[c]{\fontsize{8}{10}\selectfont 90}}
\put(21.000,82.308){\makebox(0,0)[r]{\fontsize{8}{10}\selectfont 90}}
\put(219.192,23.000){\makebox(0,0)[c]{\fontsize{8}{10}\selectfont 120}}
\put(21.000,33.808){\makebox(0,0)[r]{\fontsize{8}{10}\selectfont 120}}
\put(123.000,7.000){\makebox(0,0)[c]{\fontsize{9}{11}\selectfont column $j$}}
\put(23.000,237.000){\makebox(0,0)[r]{\fontsize{8}{10}\selectfont row $i$}}
\put(367.000,255.000){\makebox(0,0)[c]{\fontsize{10}{12}\selectfont Sparse parent matrix: $B_{120}$}}
\put(367.000,240.000){\makebox(0,0)[c]{\fontsize{9}{11}\selectfont 313 nonzero entries}}
\put(270.050,225.433){\color{cmpRow}\rule{1.517pt}{1.517pt}}
\put(271.667,225.433){\color{cmpRow}\rule{1.517pt}{1.517pt}}
\put(273.283,225.433){\color{cmpRow}\rule{1.517pt}{1.517pt}}
\put(274.900,225.433){\color{cmpRow}\rule{1.517pt}{1.517pt}}
\put(276.517,225.433){\color{cmpRow}\rule{1.517pt}{1.517pt}}
\put(278.133,225.433){\color{cmpRow}\rule{1.517pt}{1.517pt}}
\put(279.750,225.433){\color{cmpRow}\rule{1.517pt}{1.517pt}}
\put(281.367,225.433){\color{cmpRow}\rule{1.517pt}{1.517pt}}
\put(282.983,225.433){\color{cmpRow}\rule{1.517pt}{1.517pt}}
\put(284.600,225.433){\color{cmpRow}\rule{1.517pt}{1.517pt}}
\put(286.217,225.433){\color{cmpRow}\rule{1.517pt}{1.517pt}}
\put(287.833,225.433){\color{cmpRow}\rule{1.517pt}{1.517pt}}
\put(289.450,225.433){\color{cmpRow}\rule{1.517pt}{1.517pt}}
\put(291.067,225.433){\color{cmpRow}\rule{1.517pt}{1.517pt}}
\put(292.683,225.433){\color{cmpRow}\rule{1.517pt}{1.517pt}}
\put(294.300,225.433){\color{cmpRow}\rule{1.517pt}{1.517pt}}
\put(295.917,225.433){\color{cmpRow}\rule{1.517pt}{1.517pt}}
\put(297.533,225.433){\color{cmpRow}\rule{1.517pt}{1.517pt}}
\put(299.150,225.433){\color{cmpRow}\rule{1.517pt}{1.517pt}}
\put(300.767,225.433){\color{cmpRow}\rule{1.517pt}{1.517pt}}
\put(302.383,225.433){\color{cmpRow}\rule{1.517pt}{1.517pt}}
\put(304.000,225.433){\color{cmpRow}\rule{1.517pt}{1.517pt}}
\put(305.617,225.433){\color{cmpRow}\rule{1.517pt}{1.517pt}}
\put(307.233,225.433){\color{cmpRow}\rule{1.517pt}{1.517pt}}
\put(308.850,225.433){\color{cmpRow}\rule{1.517pt}{1.517pt}}
\put(310.467,225.433){\color{cmpRow}\rule{1.517pt}{1.517pt}}
\put(312.083,225.433){\color{cmpRow}\rule{1.517pt}{1.517pt}}
\put(313.700,225.433){\color{cmpRow}\rule{1.517pt}{1.517pt}}
\put(315.317,225.433){\color{cmpRow}\rule{1.517pt}{1.517pt}}
\put(316.933,225.433){\color{cmpRow}\rule{1.517pt}{1.517pt}}
\put(318.550,225.433){\color{cmpRow}\rule{1.517pt}{1.517pt}}
\put(320.167,225.433){\color{cmpRow}\rule{1.517pt}{1.517pt}}
\put(321.783,225.433){\color{cmpRow}\rule{1.517pt}{1.517pt}}
\put(323.400,225.433){\color{cmpRow}\rule{1.517pt}{1.517pt}}
\put(325.017,225.433){\color{cmpRow}\rule{1.517pt}{1.517pt}}
\put(326.633,225.433){\color{cmpRow}\rule{1.517pt}{1.517pt}}
\put(328.250,225.433){\color{cmpRow}\rule{1.517pt}{1.517pt}}
\put(329.867,225.433){\color{cmpRow}\rule{1.517pt}{1.517pt}}
\put(331.483,225.433){\color{cmpRow}\rule{1.517pt}{1.517pt}}
\put(333.100,225.433){\color{cmpRow}\rule{1.517pt}{1.517pt}}
\put(334.717,225.433){\color{cmpRow}\rule{1.517pt}{1.517pt}}
\put(336.333,225.433){\color{cmpRow}\rule{1.517pt}{1.517pt}}
\put(337.950,225.433){\color{cmpRow}\rule{1.517pt}{1.517pt}}
\put(339.567,225.433){\color{cmpRow}\rule{1.517pt}{1.517pt}}
\put(341.183,225.433){\color{cmpRow}\rule{1.517pt}{1.517pt}}
\put(342.800,225.433){\color{cmpRow}\rule{1.517pt}{1.517pt}}
\put(344.417,225.433){\color{cmpRow}\rule{1.517pt}{1.517pt}}
\put(346.033,225.433){\color{cmpRow}\rule{1.517pt}{1.517pt}}
\put(347.650,225.433){\color{cmpRow}\rule{1.517pt}{1.517pt}}
\put(349.267,225.433){\color{cmpRow}\rule{1.517pt}{1.517pt}}
\put(350.883,225.433){\color{cmpRow}\rule{1.517pt}{1.517pt}}
\put(352.500,225.433){\color{cmpRow}\rule{1.517pt}{1.517pt}}
\put(354.117,225.433){\color{cmpRow}\rule{1.517pt}{1.517pt}}
\put(355.733,225.433){\color{cmpRow}\rule{1.517pt}{1.517pt}}
\put(357.350,225.433){\color{cmpRow}\rule{1.517pt}{1.517pt}}
\put(358.967,225.433){\color{cmpRow}\rule{1.517pt}{1.517pt}}
\put(360.583,225.433){\color{cmpRow}\rule{1.517pt}{1.517pt}}
\put(362.200,225.433){\color{cmpRow}\rule{1.517pt}{1.517pt}}
\put(363.817,225.433){\color{cmpRow}\rule{1.517pt}{1.517pt}}
\put(365.433,225.433){\color{cmpRow}\rule{1.517pt}{1.517pt}}
\put(367.050,225.433){\color{cmpRow}\rule{1.517pt}{1.517pt}}
\put(368.667,225.433){\color{cmpRow}\rule{1.517pt}{1.517pt}}
\put(370.283,225.433){\color{cmpRow}\rule{1.517pt}{1.517pt}}
\put(371.900,225.433){\color{cmpRow}\rule{1.517pt}{1.517pt}}
\put(373.517,225.433){\color{cmpRow}\rule{1.517pt}{1.517pt}}
\put(375.133,225.433){\color{cmpRow}\rule{1.517pt}{1.517pt}}
\put(376.750,225.433){\color{cmpRow}\rule{1.517pt}{1.517pt}}
\put(378.367,225.433){\color{cmpRow}\rule{1.517pt}{1.517pt}}
\put(379.983,225.433){\color{cmpRow}\rule{1.517pt}{1.517pt}}
\put(381.600,225.433){\color{cmpRow}\rule{1.517pt}{1.517pt}}
\put(383.217,225.433){\color{cmpRow}\rule{1.517pt}{1.517pt}}
\put(384.833,225.433){\color{cmpRow}\rule{1.517pt}{1.517pt}}
\put(386.450,225.433){\color{cmpRow}\rule{1.517pt}{1.517pt}}
\put(388.067,225.433){\color{cmpRow}\rule{1.517pt}{1.517pt}}
\put(389.683,225.433){\color{cmpRow}\rule{1.517pt}{1.517pt}}
\put(391.300,225.433){\color{cmpRow}\rule{1.517pt}{1.517pt}}
\put(392.917,225.433){\color{cmpRow}\rule{1.517pt}{1.517pt}}
\put(394.533,225.433){\color{cmpRow}\rule{1.517pt}{1.517pt}}
\put(396.150,225.433){\color{cmpRow}\rule{1.517pt}{1.517pt}}
\put(397.767,225.433){\color{cmpRow}\rule{1.517pt}{1.517pt}}
\put(399.383,225.433){\color{cmpRow}\rule{1.517pt}{1.517pt}}
\put(401.000,225.433){\color{cmpRow}\rule{1.517pt}{1.517pt}}
\put(402.617,225.433){\color{cmpRow}\rule{1.517pt}{1.517pt}}
\put(404.233,225.433){\color{cmpRow}\rule{1.517pt}{1.517pt}}
\put(405.850,225.433){\color{cmpRow}\rule{1.517pt}{1.517pt}}
\put(407.467,225.433){\color{cmpRow}\rule{1.517pt}{1.517pt}}
\put(409.083,225.433){\color{cmpRow}\rule{1.517pt}{1.517pt}}
\put(410.700,225.433){\color{cmpRow}\rule{1.517pt}{1.517pt}}
\put(412.317,225.433){\color{cmpRow}\rule{1.517pt}{1.517pt}}
\put(413.933,225.433){\color{cmpRow}\rule{1.517pt}{1.517pt}}
\put(415.550,225.433){\color{cmpRow}\rule{1.517pt}{1.517pt}}
\put(417.167,225.433){\color{cmpRow}\rule{1.517pt}{1.517pt}}
\put(418.783,225.433){\color{cmpRow}\rule{1.517pt}{1.517pt}}
\put(420.400,225.433){\color{cmpRow}\rule{1.517pt}{1.517pt}}
\put(422.017,225.433){\color{cmpRow}\rule{1.517pt}{1.517pt}}
\put(423.633,225.433){\color{cmpRow}\rule{1.517pt}{1.517pt}}
\put(425.250,225.433){\color{cmpRow}\rule{1.517pt}{1.517pt}}
\put(426.867,225.433){\color{cmpRow}\rule{1.517pt}{1.517pt}}
\put(428.483,225.433){\color{cmpRow}\rule{1.517pt}{1.517pt}}
\put(430.100,225.433){\color{cmpRow}\rule{1.517pt}{1.517pt}}
\put(431.717,225.433){\color{cmpRow}\rule{1.517pt}{1.517pt}}
\put(433.333,225.433){\color{cmpRow}\rule{1.517pt}{1.517pt}}
\put(434.950,225.433){\color{cmpRow}\rule{1.517pt}{1.517pt}}
\put(436.567,225.433){\color{cmpRow}\rule{1.517pt}{1.517pt}}
\put(438.183,225.433){\color{cmpRow}\rule{1.517pt}{1.517pt}}
\put(439.800,225.433){\color{cmpRow}\rule{1.517pt}{1.517pt}}
\put(441.417,225.433){\color{cmpRow}\rule{1.517pt}{1.517pt}}
\put(443.033,225.433){\color{cmpRow}\rule{1.517pt}{1.517pt}}
\put(444.650,225.433){\color{cmpRow}\rule{1.517pt}{1.517pt}}
\put(446.267,225.433){\color{cmpRow}\rule{1.517pt}{1.517pt}}
\put(447.883,225.433){\color{cmpRow}\rule{1.517pt}{1.517pt}}
\put(449.500,225.433){\color{cmpRow}\rule{1.517pt}{1.517pt}}
\put(451.117,225.433){\color{cmpRow}\rule{1.517pt}{1.517pt}}
\put(452.733,225.433){\color{cmpRow}\rule{1.517pt}{1.517pt}}
\put(454.350,225.433){\color{cmpRow}\rule{1.517pt}{1.517pt}}
\put(455.967,225.433){\color{cmpRow}\rule{1.517pt}{1.517pt}}
\put(457.583,225.433){\color{cmpRow}\rule{1.517pt}{1.517pt}}
\put(459.200,225.433){\color{cmpRow}\rule{1.517pt}{1.517pt}}
\put(460.817,225.433){\color{cmpRow}\rule{1.517pt}{1.517pt}}
\put(462.433,225.433){\color{cmpRow}\rule{1.517pt}{1.517pt}}
\put(270.050,223.817){\color{cmpParent}\rule{1.517pt}{1.517pt}}
\put(271.667,223.817){\color{cmpDiag}\rule{1.517pt}{1.517pt}}
\put(270.050,222.200){\color{cmpParent}\rule{1.517pt}{1.517pt}}
\put(273.283,222.200){\color{cmpDiag}\rule{1.517pt}{1.517pt}}
\put(274.900,220.583){\color{cmpDiag}\rule{1.517pt}{1.517pt}}
\put(270.050,218.967){\color{cmpParent}\rule{1.517pt}{1.517pt}}
\put(276.517,218.967){\color{cmpDiag}\rule{1.517pt}{1.517pt}}
\put(271.667,217.350){\color{cmpParent}\rule{1.517pt}{1.517pt}}
\put(278.133,217.350){\color{cmpDiag}\rule{1.517pt}{1.517pt}}
\put(270.050,215.733){\color{cmpParent}\rule{1.517pt}{1.517pt}}
\put(279.750,215.733){\color{cmpDiag}\rule{1.517pt}{1.517pt}}
\put(281.367,214.117){\color{cmpDiag}\rule{1.517pt}{1.517pt}}
\put(282.983,212.500){\color{cmpDiag}\rule{1.517pt}{1.517pt}}
\put(271.667,210.883){\color{cmpParent}\rule{1.517pt}{1.517pt}}
\put(284.600,210.883){\color{cmpDiag}\rule{1.517pt}{1.517pt}}
\put(270.050,209.267){\color{cmpParent}\rule{1.517pt}{1.517pt}}
\put(286.217,209.267){\color{cmpDiag}\rule{1.517pt}{1.517pt}}
\put(287.833,207.650){\color{cmpDiag}\rule{1.517pt}{1.517pt}}
\put(270.050,206.033){\color{cmpParent}\rule{1.517pt}{1.517pt}}
\put(289.450,206.033){\color{cmpDiag}\rule{1.517pt}{1.517pt}}
\put(271.667,204.417){\color{cmpParent}\rule{1.517pt}{1.517pt}}
\put(291.067,204.417){\color{cmpDiag}\rule{1.517pt}{1.517pt}}
\put(273.283,202.800){\color{cmpParent}\rule{1.517pt}{1.517pt}}
\put(292.683,202.800){\color{cmpDiag}\rule{1.517pt}{1.517pt}}
\put(294.300,201.183){\color{cmpDiag}\rule{1.517pt}{1.517pt}}
\put(270.050,199.567){\color{cmpParent}\rule{1.517pt}{1.517pt}}
\put(295.917,199.567){\color{cmpDiag}\rule{1.517pt}{1.517pt}}
\put(297.533,197.950){\color{cmpDiag}\rule{1.517pt}{1.517pt}}
\put(270.050,196.333){\color{cmpParent}\rule{1.517pt}{1.517pt}}
\put(299.150,196.333){\color{cmpDiag}\rule{1.517pt}{1.517pt}}
\put(300.767,194.717){\color{cmpDiag}\rule{1.517pt}{1.517pt}}
\put(273.283,193.100){\color{cmpParent}\rule{1.517pt}{1.517pt}}
\put(302.383,193.100){\color{cmpDiag}\rule{1.517pt}{1.517pt}}
\put(271.667,191.483){\color{cmpParent}\rule{1.517pt}{1.517pt}}
\put(304.000,191.483){\color{cmpDiag}\rule{1.517pt}{1.517pt}}
\put(270.050,189.867){\color{cmpParent}\rule{1.517pt}{1.517pt}}
\put(305.617,189.867){\color{cmpDiag}\rule{1.517pt}{1.517pt}}
\put(307.233,188.250){\color{cmpDiag}\rule{1.517pt}{1.517pt}}
\put(308.850,186.633){\color{cmpDiag}\rule{1.517pt}{1.517pt}}
\put(271.667,185.017){\color{cmpParent}\rule{1.517pt}{1.517pt}}
\put(310.467,185.017){\color{cmpDiag}\rule{1.517pt}{1.517pt}}
\put(312.083,183.400){\color{cmpDiag}\rule{1.517pt}{1.517pt}}
\put(313.700,181.783){\color{cmpDiag}\rule{1.517pt}{1.517pt}}
\put(270.050,180.167){\color{cmpParent}\rule{1.517pt}{1.517pt}}
\put(315.317,180.167){\color{cmpDiag}\rule{1.517pt}{1.517pt}}
\put(278.133,178.550){\color{cmpParent}\rule{1.517pt}{1.517pt}}
\put(316.933,178.550){\color{cmpDiag}\rule{1.517pt}{1.517pt}}
\put(270.050,176.933){\color{cmpParent}\rule{1.517pt}{1.517pt}}
\put(318.550,176.933){\color{cmpDiag}\rule{1.517pt}{1.517pt}}
\put(320.167,175.317){\color{cmpDiag}\rule{1.517pt}{1.517pt}}
\put(273.283,173.700){\color{cmpParent}\rule{1.517pt}{1.517pt}}
\put(321.783,173.700){\color{cmpDiag}\rule{1.517pt}{1.517pt}}
\put(271.667,172.083){\color{cmpParent}\rule{1.517pt}{1.517pt}}
\put(323.400,172.083){\color{cmpDiag}\rule{1.517pt}{1.517pt}}
\put(276.517,170.467){\color{cmpParent}\rule{1.517pt}{1.517pt}}
\put(325.017,170.467){\color{cmpDiag}\rule{1.517pt}{1.517pt}}
\put(326.633,168.850){\color{cmpDiag}\rule{1.517pt}{1.517pt}}
\put(270.050,167.233){\color{cmpParent}\rule{1.517pt}{1.517pt}}
\put(328.250,167.233){\color{cmpDiag}\rule{1.517pt}{1.517pt}}
\put(271.667,165.617){\color{cmpParent}\rule{1.517pt}{1.517pt}}
\put(329.867,165.617){\color{cmpDiag}\rule{1.517pt}{1.517pt}}
\put(273.283,164.000){\color{cmpParent}\rule{1.517pt}{1.517pt}}
\put(331.483,164.000){\color{cmpDiag}\rule{1.517pt}{1.517pt}}
\put(333.100,162.383){\color{cmpDiag}\rule{1.517pt}{1.517pt}}
\put(270.050,160.767){\color{cmpParent}\rule{1.517pt}{1.517pt}}
\put(334.717,160.767){\color{cmpDiag}\rule{1.517pt}{1.517pt}}
\put(278.133,159.150){\color{cmpParent}\rule{1.517pt}{1.517pt}}
\put(336.333,159.150){\color{cmpDiag}\rule{1.517pt}{1.517pt}}
\put(270.050,157.533){\color{cmpParent}\rule{1.517pt}{1.517pt}}
\put(337.950,157.533){\color{cmpDiag}\rule{1.517pt}{1.517pt}}
\put(339.567,155.917){\color{cmpDiag}\rule{1.517pt}{1.517pt}}
\put(341.183,154.300){\color{cmpDiag}\rule{1.517pt}{1.517pt}}
\put(271.667,152.683){\color{cmpParent}\rule{1.517pt}{1.517pt}}
\put(342.800,152.683){\color{cmpDiag}\rule{1.517pt}{1.517pt}}
\put(270.050,151.067){\color{cmpParent}\rule{1.517pt}{1.517pt}}
\put(344.417,151.067){\color{cmpDiag}\rule{1.517pt}{1.517pt}}
\put(346.033,149.450){\color{cmpDiag}\rule{1.517pt}{1.517pt}}
\put(347.650,147.833){\color{cmpDiag}\rule{1.517pt}{1.517pt}}
\put(349.267,146.217){\color{cmpDiag}\rule{1.517pt}{1.517pt}}
\put(273.283,144.600){\color{cmpParent}\rule{1.517pt}{1.517pt}}
\put(350.883,144.600){\color{cmpDiag}\rule{1.517pt}{1.517pt}}
\put(352.500,142.983){\color{cmpDiag}\rule{1.517pt}{1.517pt}}
\put(270.050,141.367){\color{cmpParent}\rule{1.517pt}{1.517pt}}
\put(354.117,141.367){\color{cmpDiag}\rule{1.517pt}{1.517pt}}
\put(355.733,139.750){\color{cmpDiag}\rule{1.517pt}{1.517pt}}
\put(276.517,138.133){\color{cmpParent}\rule{1.517pt}{1.517pt}}
\put(357.350,138.133){\color{cmpDiag}\rule{1.517pt}{1.517pt}}
\put(358.967,136.517){\color{cmpDiag}\rule{1.517pt}{1.517pt}}
\put(273.283,134.900){\color{cmpParent}\rule{1.517pt}{1.517pt}}
\put(360.583,134.900){\color{cmpDiag}\rule{1.517pt}{1.517pt}}
\put(271.667,133.283){\color{cmpParent}\rule{1.517pt}{1.517pt}}
\put(362.200,133.283){\color{cmpDiag}\rule{1.517pt}{1.517pt}}
\put(270.050,131.667){\color{cmpParent}\rule{1.517pt}{1.517pt}}
\put(363.817,131.667){\color{cmpDiag}\rule{1.517pt}{1.517pt}}
\put(365.433,130.050){\color{cmpDiag}\rule{1.517pt}{1.517pt}}
\put(270.050,128.433){\color{cmpParent}\rule{1.517pt}{1.517pt}}
\put(367.050,128.433){\color{cmpDiag}\rule{1.517pt}{1.517pt}}
\put(271.667,126.817){\color{cmpParent}\rule{1.517pt}{1.517pt}}
\put(368.667,126.817){\color{cmpDiag}\rule{1.517pt}{1.517pt}}
\put(370.283,125.200){\color{cmpDiag}\rule{1.517pt}{1.517pt}}
\put(371.900,123.583){\color{cmpDiag}\rule{1.517pt}{1.517pt}}
\put(276.517,121.967){\color{cmpParent}\rule{1.517pt}{1.517pt}}
\put(373.517,121.967){\color{cmpDiag}\rule{1.517pt}{1.517pt}}
\put(278.133,120.350){\color{cmpParent}\rule{1.517pt}{1.517pt}}
\put(375.133,120.350){\color{cmpDiag}\rule{1.517pt}{1.517pt}}
\put(270.050,118.733){\color{cmpParent}\rule{1.517pt}{1.517pt}}
\put(376.750,118.733){\color{cmpDiag}\rule{1.517pt}{1.517pt}}
\put(378.367,117.117){\color{cmpDiag}\rule{1.517pt}{1.517pt}}
\put(273.283,115.500){\color{cmpParent}\rule{1.517pt}{1.517pt}}
\put(379.983,115.500){\color{cmpDiag}\rule{1.517pt}{1.517pt}}
\put(284.600,113.883){\color{cmpParent}\rule{1.517pt}{1.517pt}}
\put(381.600,113.883){\color{cmpDiag}\rule{1.517pt}{1.517pt}}
\put(270.050,112.267){\color{cmpParent}\rule{1.517pt}{1.517pt}}
\put(383.217,112.267){\color{cmpDiag}\rule{1.517pt}{1.517pt}}
\put(384.833,110.650){\color{cmpDiag}\rule{1.517pt}{1.517pt}}
\put(270.050,109.033){\color{cmpParent}\rule{1.517pt}{1.517pt}}
\put(386.450,109.033){\color{cmpDiag}\rule{1.517pt}{1.517pt}}
\put(271.667,107.417){\color{cmpParent}\rule{1.517pt}{1.517pt}}
\put(388.067,107.417){\color{cmpDiag}\rule{1.517pt}{1.517pt}}
\put(389.683,105.800){\color{cmpDiag}\rule{1.517pt}{1.517pt}}
\put(391.300,104.183){\color{cmpDiag}\rule{1.517pt}{1.517pt}}
\put(279.750,102.567){\color{cmpParent}\rule{1.517pt}{1.517pt}}
\put(392.917,102.567){\color{cmpDiag}\rule{1.517pt}{1.517pt}}
\put(278.133,100.950){\color{cmpParent}\rule{1.517pt}{1.517pt}}
\put(394.533,100.950){\color{cmpDiag}\rule{1.517pt}{1.517pt}}
\put(270.050,99.333){\color{cmpParent}\rule{1.517pt}{1.517pt}}
\put(396.150,99.333){\color{cmpDiag}\rule{1.517pt}{1.517pt}}
\put(397.767,97.717){\color{cmpDiag}\rule{1.517pt}{1.517pt}}
\put(399.383,96.100){\color{cmpDiag}\rule{1.517pt}{1.517pt}}
\put(271.667,94.483){\color{cmpParent}\rule{1.517pt}{1.517pt}}
\put(401.000,94.483){\color{cmpDiag}\rule{1.517pt}{1.517pt}}
\put(270.050,92.867){\color{cmpParent}\rule{1.517pt}{1.517pt}}
\put(402.617,92.867){\color{cmpDiag}\rule{1.517pt}{1.517pt}}
\put(404.233,91.250){\color{cmpDiag}\rule{1.517pt}{1.517pt}}
\put(276.517,89.633){\color{cmpParent}\rule{1.517pt}{1.517pt}}
\put(405.850,89.633){\color{cmpDiag}\rule{1.517pt}{1.517pt}}
\put(271.667,88.017){\color{cmpParent}\rule{1.517pt}{1.517pt}}
\put(407.467,88.017){\color{cmpDiag}\rule{1.517pt}{1.517pt}}
\put(273.283,86.400){\color{cmpParent}\rule{1.517pt}{1.517pt}}
\put(409.083,86.400){\color{cmpDiag}\rule{1.517pt}{1.517pt}}
\put(410.700,84.783){\color{cmpDiag}\rule{1.517pt}{1.517pt}}
\put(270.050,83.167){\color{cmpParent}\rule{1.517pt}{1.517pt}}
\put(412.317,83.167){\color{cmpDiag}\rule{1.517pt}{1.517pt}}
\put(413.933,81.550){\color{cmpDiag}\rule{1.517pt}{1.517pt}}
\put(279.750,79.933){\color{cmpParent}\rule{1.517pt}{1.517pt}}
\put(415.550,79.933){\color{cmpDiag}\rule{1.517pt}{1.517pt}}
\put(417.167,78.317){\color{cmpDiag}\rule{1.517pt}{1.517pt}}
\put(273.283,76.700){\color{cmpParent}\rule{1.517pt}{1.517pt}}
\put(418.783,76.700){\color{cmpDiag}\rule{1.517pt}{1.517pt}}
\put(271.667,75.083){\color{cmpParent}\rule{1.517pt}{1.517pt}}
\put(420.400,75.083){\color{cmpDiag}\rule{1.517pt}{1.517pt}}
\put(276.517,73.467){\color{cmpParent}\rule{1.517pt}{1.517pt}}
\put(422.017,73.467){\color{cmpDiag}\rule{1.517pt}{1.517pt}}
\put(423.633,71.850){\color{cmpDiag}\rule{1.517pt}{1.517pt}}
\put(270.050,70.233){\color{cmpParent}\rule{1.517pt}{1.517pt}}
\put(425.250,70.233){\color{cmpDiag}\rule{1.517pt}{1.517pt}}
\put(426.867,68.617){\color{cmpDiag}\rule{1.517pt}{1.517pt}}
\put(428.483,67.000){\color{cmpDiag}\rule{1.517pt}{1.517pt}}
\put(430.100,65.383){\color{cmpDiag}\rule{1.517pt}{1.517pt}}
\put(270.050,63.767){\color{cmpParent}\rule{1.517pt}{1.517pt}}
\put(431.717,63.767){\color{cmpDiag}\rule{1.517pt}{1.517pt}}
\put(278.133,62.150){\color{cmpParent}\rule{1.517pt}{1.517pt}}
\put(433.333,62.150){\color{cmpDiag}\rule{1.517pt}{1.517pt}}
\put(270.050,60.533){\color{cmpParent}\rule{1.517pt}{1.517pt}}
\put(434.950,60.533){\color{cmpDiag}\rule{1.517pt}{1.517pt}}
\put(436.567,58.917){\color{cmpDiag}\rule{1.517pt}{1.517pt}}
\put(292.683,57.300){\color{cmpParent}\rule{1.517pt}{1.517pt}}
\put(438.183,57.300){\color{cmpDiag}\rule{1.517pt}{1.517pt}}
\put(271.667,55.683){\color{cmpParent}\rule{1.517pt}{1.517pt}}
\put(439.800,55.683){\color{cmpDiag}\rule{1.517pt}{1.517pt}}
\put(270.050,54.067){\color{cmpParent}\rule{1.517pt}{1.517pt}}
\put(441.417,54.067){\color{cmpDiag}\rule{1.517pt}{1.517pt}}
\put(443.033,52.450){\color{cmpDiag}\rule{1.517pt}{1.517pt}}
\put(270.050,50.833){\color{cmpParent}\rule{1.517pt}{1.517pt}}
\put(444.650,50.833){\color{cmpDiag}\rule{1.517pt}{1.517pt}}
\put(284.600,49.217){\color{cmpParent}\rule{1.517pt}{1.517pt}}
\put(446.267,49.217){\color{cmpDiag}\rule{1.517pt}{1.517pt}}
\put(273.283,47.600){\color{cmpParent}\rule{1.517pt}{1.517pt}}
\put(447.883,47.600){\color{cmpDiag}\rule{1.517pt}{1.517pt}}
\put(449.500,45.983){\color{cmpDiag}\rule{1.517pt}{1.517pt}}
\put(270.050,44.367){\color{cmpParent}\rule{1.517pt}{1.517pt}}
\put(451.117,44.367){\color{cmpDiag}\rule{1.517pt}{1.517pt}}
\put(278.133,42.750){\color{cmpParent}\rule{1.517pt}{1.517pt}}
\put(452.733,42.750){\color{cmpDiag}\rule{1.517pt}{1.517pt}}
\put(276.517,41.133){\color{cmpParent}\rule{1.517pt}{1.517pt}}
\put(454.350,41.133){\color{cmpDiag}\rule{1.517pt}{1.517pt}}
\put(455.967,39.517){\color{cmpDiag}\rule{1.517pt}{1.517pt}}
\put(457.583,37.900){\color{cmpDiag}\rule{1.517pt}{1.517pt}}
\put(271.667,36.283){\color{cmpParent}\rule{1.517pt}{1.517pt}}
\put(459.200,36.283){\color{cmpDiag}\rule{1.517pt}{1.517pt}}
\put(279.750,34.667){\color{cmpParent}\rule{1.517pt}{1.517pt}}
\put(460.817,34.667){\color{cmpDiag}\rule{1.517pt}{1.517pt}}
\put(462.433,33.050){\color{cmpDiag}\rule{1.517pt}{1.517pt}}
\put(270.808,23.000){\makebox(0,0)[c]{\fontsize{8}{10}\selectfont 1}}
\put(265.000,226.192){\makebox(0,0)[r]{\fontsize{8}{10}\selectfont 1}}
\put(317.692,23.000){\makebox(0,0)[c]{\fontsize{8}{10}\selectfont 30}}
\put(265.000,179.308){\makebox(0,0)[r]{\fontsize{8}{10}\selectfont 30}}
\put(366.192,23.000){\makebox(0,0)[c]{\fontsize{8}{10}\selectfont 60}}
\put(265.000,130.808){\makebox(0,0)[r]{\fontsize{8}{10}\selectfont 60}}
\put(414.692,23.000){\makebox(0,0)[c]{\fontsize{8}{10}\selectfont 90}}
\put(265.000,82.308){\makebox(0,0)[r]{\fontsize{8}{10}\selectfont 90}}
\put(463.192,23.000){\makebox(0,0)[c]{\fontsize{8}{10}\selectfont 120}}
\put(265.000,33.808){\makebox(0,0)[r]{\fontsize{8}{10}\selectfont 120}}
\put(367.000,7.000){\makebox(0,0)[c]{\fontsize{9}{11}\selectfont column $j$}}
\put(267.000,237.000){\makebox(0,0)[r]{\fontsize{8}{10}\selectfont row $i$}}
\end{picture}
\endgroup

\caption{Nonzero entries at $n=120$, with identical row and column scales. Blank entries are zero. Orange marks the first row, gray the remaining diagonal, green the parent links retained in $B_n$, and blue the other divisor entries present only in $R_n^T$. For example, row $30$ of $R_n^T$ has ones in columns $1,2,3,5,6,10,15,30$, whereas row $30$ of $B_n$ has ones only in columns $6=30/5$ and $30$.}
\label{fig:redheffer}
\end{figure}

All constructions are defined formally in Section~\ref{sec:definitions}. The parent-incidence matrix is denoted by $S_n$, the triangular matrix $I_n+S_n$ by $A_n$, and the matrix obtained by replacing its first row by ones by $B_n$. Singular values are always listed in decreasing order. The integer $k_n$ counts vertices that have at least one child.

The principal spectral conclusions are these. For $n\ge4$, exactly $n-2k_n-1$ singular values of $B_n$ equal one. Exactly $k_n+1$ are larger than one and $k_n$ are smaller than one, with a possible zero included in the latter count. Moreover $k_n=O_H(n/(\log n)^H)$ for every fixed positive $H$. Apart from the largest singular value, which is asymptotic to $\sqrt n$, the fixed upper singular values are asymptotic to
\[
 \sqrt{\frac{n}{a_r\log n}},\qquad r=1,2,\ldots,
\]
where $a_r$ is the $r$th squarefree positive integer. The underlying comparison has a uniform additive error less than three and therefore also applies to growing indices whenever the corresponding degree tends to infinity.

At the lower end, the inverse Gram matrices of $A_n$, divided by $n$, converge in Hilbert--Schmidt norm to an explicit positive compact operator. Removing a specified rank-one part yields an operator $L$ whose positive eigenvalues determine every fixed nonexceptional lower singular value of $B_n$. The limit remains valid at integers for which the Mertens sum vanishes. The proof at those integers uses an exact Moore--Penrose inverse rather than an invalid division by the determinant.

The exceptional direction requires a different ingredient. We define the restricted signed count
\[
 R(x,y)=\sum_{\substack{d\le x\\P^-(d)>y}}\mu(d),
\]
including the term $d=1$. Here $P^-(d)$ is the smallest prime factor and $\mu$ is the M\"obius function. The relevant inverse vector $w_n$ has these signed counts as its squarefree coordinates. We calculate the leading exponential scale of its Euclidean norm $W_n$ by first summing $R(n/p,p)^2$ over primes, then controlling the composite indices:
\[
 W_n=n\exp\bigl(-(1+o(1))\sqrt{\log n\log\log n}\bigr).
\]
Combined with the exact inverse, this gives the unconditional equivalent
\[
 \sigma_n(B_n)\sim\frac{|M(n)|}{\sqrt{Q(n)}\,W_n}
 =|M(n)|\,n^{-3/2}\exp\bigl((1+o(1))\sqrt{\log n\log\log n}\bigr)
 \qquad(M(n)\ne0),
\]
where $Q(n)$ counts the squarefree integers up to $n$. Consequently $\sigma_n(B_n)\ll_\varepsilon n^{-1+\varepsilon}$ for every $\varepsilon>0$ if and only if $M(n)\ll_\varepsilon n^{1/2+\varepsilon}$ for every $\varepsilon>0$, which is equivalent to the Riemann hypothesis. The formula contains $M(n)$, so it restates the Mertens sum rather than bounding it. We also give all fixed moments of order greater than one and a sharp leading exponent under a prescribed support bound for least-prime weights. These analytic results use existing cancellation estimates, not a new source of ordinary M\"obius cancellation.

\subsection{Relation to earlier work and the scope of novelty}

The sparse matrix $B_n$ was introduced in Kline~\cite{kline2019}, and its dominant eigenvalues were determined in~\cite{kline2020}. It is the matrix denoted by $\mathcal R_n$ in that work. We use the letter $B$ here to distinguish it from both Redheffer's matrix and the restricted sum $R(x,y)$; this is a change of notation, not a new matrix construction. The base-volume identity $\det(H_nH_n^T)=Q(n)$ in Section~\ref{sec:geometry} is a case of the bordered-matrix identity of~\cite{kline2020b}.

The closest predecessor is the author's earlier release~\cite{klineSmallest} on the smallest singular value of the same matrix. It proved the triangular inverse, the expression of $w_n$ by restricted M\"obius sums, the rank-one inverse formula, $\sigma_1(B_n)=\sqrt n\,(1+O(1/n))$, $\norm{A_n^{-1}}=n^{1/2+o(1)}$, and the unconditional bracket
\[
 |M(n)|\,n^{-3/2+o(1)}\le\sigma_n(B_n)\le|M(n)|\,n^{-4/3+o(1)},
\]
together with the unconditional equivalence of the Riemann hypothesis with $\sigma_n(B_n)\ll_\varepsilon n^{-1+\varepsilon}$. It proved $\sigma_n(B_n)=|M(n)|n^{-3/2+o(1)}$ under a rank-one dominance condition implied by the Riemann hypothesis, and it posed the shape $W_n=n\exp(-(1+o(1))\sqrt{\log n\log\log n})$ as an open problem. We rederive the shared identities so that no earlier manuscript is needed to read this one. Our estimate $\sigma_1(B_n)=\sqrt n+O(1)$ is weaker than the earlier one; it is included because the same argument gives the rest of the upper spectrum.

The new relations to that release are these. Theorem~\ref{arith:energy} proves the norm shape posed there. Theorem~\ref{spec:minimum} then shows that the dominance condition holds unconditionally, so the bracket becomes an asymptotic equivalent without any hypothesis. Theorem~\ref{spec:compact} replaces $n^{1/2+o(1)}$ by the exact limit $\norm{A_n^{-1}}/\sqrt n\to\norm{K}^{1/2}$. The unit spectrum, the degree comparison, the fixed lower spectrum, the products of singular values, and the geometric results do not appear in the earlier release. Unpublished notes of the author also contain a two-sided bound of order $\sqrt n$ for $\norm{A_n^{-1}}$.

Another release by the author~\cite{klineExtremal} studies row-modified weighted divisor matrices with entries $(j/i)^s$ when $j$ divides $i$. Its dominant inverse term and extreme singular-value analysis are related in method, but those divisor matrices differ from the parent forest here.

Hilberdink~\cite{hilberdink}, in particular Theorems~2.2 and~3.2 and Section~4(b), develops compact Gram limits for multiplicative Toeplitz matrices and derives arithmetic singular-value asymptotics. Thus neither the general method nor the idea of a fixed lower-spectrum limit is claimed as new. A multiplicative Toeplitz entry has the form $f(i/j)$, with a prescribed support on quotients. The forest inverse here does not have that form: its entries at $(2,1)$ and $(6,3)$ are respectively $-1$ and $0$, although both quotients equal two. Our candidate contribution is the explicit ancestor-density kernel, its uniform summable tail, the deflation for the bordered matrix, and the exact unit-spectrum and degree comparisons.

Alladi~\cite{alladi1982} already studies restricted M\"obius sums and the unrestricted bounded least-prime weighted extremum. McNew~\cite{mcnew} studies the analogous smooth-number maximum and supplies an accessible account of the smooth-count expansions used here. The saddle scale itself is classical. The arithmetic moment and support-constrained formulations below are derived refinements; we do not claim that a bounded search has established their publication priority. The full theorem pages of Alladi's 1982 paper were not accessible in the present survey, so a decisive theorem-by-theorem comparison remains a release obligation. The zero value of the unweighted thin-prime reciprocal series is already included in the density theorem of Kural, McDonald, and Sah~\cite{kms}; the quantitative weighted tail is the additional formulation considered here.

These comparisons cover a bounded literature search, recorded with the release; they do not establish worldwide novelty. The geometric component functional is included as an explicit interpretation and consequence of classical smooth and rough counting, rather than as a separate priority headline. No result improves a known estimate for the Mertens function, proves the prime number theorem by a new geometric argument, or advances the Riemann hypothesis.

\subsection{Logical dependencies}

The finite forest identities, exact unit spectrum, degree comparison, and compact inverse-Gram limit use elementary counting and linear algebra. The explicit upper constants use the prime number theorem. The restricted-sum moment scale uses a classical Mertens bound and a two-term smooth-count expansion, both stated with their needed ranges. Those estimates separate the exceptional singular direction. The component profile uses fixed-parameter smooth and rough asymptotics. An optional appendix applies a general Hal\'asz theorem to the component residual; it is not needed for the spectral theorems.

The proofs are self-contained relative to these explicitly stated classical inputs. No numerical experiment is required to prove any assertion in the manuscript. Numerical approximations to the limiting eigenvalues are omitted here because finite sections alone do not certify their limiting decimal values.

\section{Definitions and exact finite identities}\label{sec:definitions}

\subsection{Arithmetic notation}

An integer is \emph{squarefree} if no square of a prime divides it. The M\"obius function is $\mu(1)=1$, $\mu(m)=(-1)^r$ when $m$ is a product of $r$ distinct primes, and $\mu(m)=0$ otherwise. Let
\[
 M(x)=\sum_{m\le x}\mu(m),\qquad
 Q(n)=\sum_{m\le n}\mu(m)^2,\qquad c_0=\frac1{\zeta(2)}=\frac6{\pi^2}.
\]
Here and below a sum over integers bounded by a real number uses its floor. The function $\zeta(s)=\sum_{m\ge1}m^{-s}$ for $\Re s>1$ is the Riemann zeta function. Elementary squarefree counting gives $Q(n)=c_0n+O(\sqrt n)$: use $\mu(m)^2=\sum_{d^2\mid m}\mu(d)$, interchange sums, and bound the tail of $\sum d^{-2}$.

Write $P^+(m)$ and $P^-(m)$ for the largest and smallest prime factors of $m>1$, with $P^+(1)=1$ and $P^-(1)=\infty$. The function $\pi(x)$ counts primes at most $x$; the unadorned constant $\pi$ retains its usual geometric meaning. An integer is $y$-smooth if $P^+(m)\le y$ and $y$-rough if $P^-(m)>y$. The inequalities are inclusive for smoothness and strict for roughness. For $x\ge1$ and $y\ge1$, define
\begin{align}
 \Psi(x,y)&=\#\{m\le x:P^+(m)\le y\},\notag\\
 \Phi(x,y)&=\#\{m\le x:P^-(m)>y\},\notag\\
 \Phi_{\sfq}(x,y)&=\sum_{\substack{m\le x\\P^-(m)>y}}\mu(m)^2,\label{def:counts}\\
 R(x,y)&=\sum_{\substack{m\le x\\P^-(m)>y}}\mu(m).\notag
\end{align}
All four quantities include the integer one. In particular $\Phi_{\sfq}(x,y)\ge1$. The signed sum $R$ is not a count of primes and need not be positive. For arithmetic functions $f,g$, their Dirichlet convolution is $(f*g)(m)=\sum_{d\mid m}f(d)g(m/d)$. An indicator $\mathbf1_{\mathcal E}$ equals one when the condition $\mathcal E$ holds and zero otherwise.

For sufficiently large $x$ put $s(x)=\sqrt{\log x\log\log x}$; all logarithms are natural. For a prime $p\le n$ set $b_p(n)=R(n/p,p)$ and define the two nonnegative sums
\begin{equation}\label{def:energies}
 E_P(n)=\sum_{p\le n}|b_p(n)|^2,\qquad
 E(n)=\sum_{\substack{2\le j\le n\\\mu(j)^2=1}}
          R(n/j,P^+(j))^2.
\end{equation}
The word \emph{energy}, when used for these quantities, means precisely this sum of squared signed counts. It carries no additional normalization. The later component quantity $F(n,y)$ has different denominators and will be defined separately.

For positive functions, $f\sim g$ means $f/g\to1$. The notation $f\ll g$ or $f=O(g)$ means $|f|\le Cg$ eventually; subscripts indicate permitted dependence of $C$. The expression $o(1)$ tends to zero in the stated limit. An equality such as $E_P=n^2\exp(-(2+o(1))s(n))$ specifies a leading logarithmic scale, not a multiplicative equivalent with an explicit prefactor.

The notation $f\asymp g$ means that both $f\ll g$ and $g\ll f$ hold. Unless another limit is specified, asymptotic statements concern $n\to\infty$ through positive integers.

\subsection{The tree and its matrices}

For each squarefree $i>1$ draw an edge from its parent $i/P^+(i)$ to $i$. Deleting the largest prime repeatedly reaches one, so the squarefree vertices at most $n$ form a tree. Each nonsquarefree vertex is isolated. We say $j\preceq i$ if $j$ and $i$ are squarefree and $j$ occurs on the parent chain from $i$ to one, including $j=i$. Thus $j$ consists of an initial segment of the prime factors of $i$, listed in increasing order. A descendant of $j$ is a vertex $i$ with $j\preceq i$. For squarefree $j$, its number of descendants is
\[
 D_j(n)=\Phi_{\sfq}(n/j,P^+(j)).
\]
Its number of children is
\begin{equation}\label{def:degree}
 d_j(n)=\max\{0,\pi(n/j)-\pi(P^+(j))\}.
\end{equation}
Set $d_j(n)=0$ for nonsquarefree $j$, and let $k_n=\#\{j\le n:d_j(n)>0\}$. Let $d_{(1)}(n)\ge\cdots\ge d_{(n)}(n)$ be these degrees in decreasing order, including zeros. In particular $d_1(n)=\pi(n)$.

In $\mathbb R^n$, $e_j$ is the $j$th coordinate vector, $\one_n=(1,\ldots,1)^T$, and $I_n$ is the identity. Define
\begin{equation}\label{def:matrices}
 (S_n)_{ij}=\begin{cases}1,&\mu(i)^2=1,\ i>1,\ j=i/P^+(i),\\0,&\text{otherwise},\end{cases}
 \quad A_n=I_n+S_n,\quad
 B_n=A_n+e_1u_n^T,\quad u_n=\one_n-e_1.
\end{equation}
The first row of $A_n$ is $e_1^T$ and the first row of $B_n$ is $\one_n^T$. Their other rows agree. The matrix $S_n$ is strictly lower triangular; hence it is nilpotent and $A_n$ is invertible with determinant one. Set
\begin{equation}\label{def:vectors}
 C_n=A_n^{-1},\qquad
 \muvec_n=(\mu(1),\ldots,\mu(n))^T,\qquad
 w_n=C_n^Tu_n,\qquad W_n=\norm{w_n}_2.
\end{equation}
Thus $w_n$ is a specifically defined vector, rather than a standard arithmetic function. Its coordinate description below explains why restricted M\"obius sums enter.

For $n=6$, the nonzero off-diagonal entries of $S_n$ are at $(2,1),(3,1),(5,1),(6,2)$; the index four is isolated. In this case
\[
 \muvec_6=(1,-1,-1,0,-1,1)^T,\qquad
 w_6=(-2,0,1,1,1,1)^T.
\]
This example also distinguishes the ancestor relation from divisibility: three divides six, but is not its ancestor.

\subsection{Norms, singular values, and projections}

The real Hilbert space $\ell^2(\mathbb N)$ consists of sequences $x=(x_j)_{j\ge1}$ for which $\sum_jx_j^2<\infty$, with inner product $\langle x,y\rangle=\sum_jx_jy_j$. A vector in $\mathbb R^n$ is viewed there by adjoining zero coordinates; a finite matrix acts as zero outside its first $n$ coordinates whenever an infinite-space limit is discussed.

The Euclidean inner product is $x^Ty$ and its norm is $\norm{x}_2$. The operator norm of a matrix $T$ is $\norm{T}_2=\sup_{\norm{x}_2=1}\norm{Tx}_2$; its Frobenius norm is $\norm{T}_F=(\sum_{ij}|T_{ij}|^2)^{1/2}$. Its singular values are the square roots of the eigenvalues of $T^TT$, counted with multiplicity. For the square matrices here we write
\[
 \sigma_1(T)\ge\cdots\ge\sigma_n(T)\ge0.
\]
When the matrix is omitted, $\sigma_i$ means $\sigma_i(B_n)$. For a nonzero vector $z$, the orthogonal projection onto its perpendicular subspace is explicitly
\begin{equation}\label{def:projection}
 \mathsf P_z=I-\frac{zz^T}{z^Tz},\qquad
 \mathsf P_z x=x-z\frac{z^Tx}{z^Tz}.
\end{equation}
For a unit vector the denominator is one. This projection removes the component parallel to $z$; it is not the projection onto its span. All later projections specify their ambient space. A subspace \emph{reduces} a matrix if it and its orthogonal complement are invariant under both the matrix and its transpose.

\begin{proposition}[Exact inverse and determinant]\label{def:inverse}
On squarefree indices,
\[
 (C_n)_{ij}=\begin{cases}\mu(i/j),&j\preceq i,\\0,&\text{otherwise}.
 \end{cases}
\]
Each nonsquarefree row and column of $C_n$ has just its diagonal entry one. Moreover
\begin{equation}\label{def:rankone}
 C_ne_1=\muvec_n,\qquad \det B_n=M(n),\qquad
 B_n^{-1}=C_n-\frac{\muvec_nw_n^T}{M(n)}\quad(M(n)\ne0).
\end{equation}
For all $n$, including $M(n)=0$,
\begin{equation}\label{def:adjugate}
 \adj(B_n)=M(n)C_n-\muvec_nw_n^T.
\end{equation}
\end{proposition}
\begin{proof}
The finite geometric series $C_n=\sum_{r=0}^{n-1}(-S_n)^r$ counts the unique parent path from an ancestor $j$ to a descendant $i$. A path of length $r$ adds $r$ distinct primes and therefore has sign $(-1)^r=\mu(i/j)$. Isolated vertices give the stated diagonal entries. This also proves the first column formula.

The determinant identity for a rank-one update gives
\[
 \det(A_n+e_1u_n^T)=1+u_n^TC_ne_1
 =1+\sum_{j=2}^n\mu(j)=M(n).
\]
Multiplication verifies the inverse formula when the scalar is nonzero. To obtain the adjugate identity without a nonsingularity restriction, apply the same calculation to $A_n+t e_1u_n^T$. For all but at most one real $t$ its adjugate is
\[
 (1+t u_n^TC_ne_1)C_n-t\muvec_nw_n^T.
\]
Both sides are matrix polynomials in $t$, so the identity holds for every $t$, in particular at $t=1$.
\end{proof}

\begin{proposition}[Coordinates and squared norm of $w_n$]\label{def:wnorm}
For $n\ge2$,
\[
 (w_n)_j=\begin{cases}
 M(n)-1,&j=1,\\
 R(n/j,P^+(j)),&j\ge2\text{ squarefree},\\
 1,&j\text{ nonsquarefree}.
 \end{cases}
\]
Consequently
\begin{equation}\label{def:wnormidentity}
 W_n^2=E(n)+(M(n)-1)^2+n-Q(n),\qquad
 \norm{\muvec_n}_2^2=Q(n).
\end{equation}
\end{proposition}
\begin{proof}
The $j$th coordinate of $C_n^Tu_n$ is the sum of column $j$ of $C_n$ over rows two through $n$. For $j=1$ it is the sum of $\mu(i)$ with its first term removed. For squarefree $j\ge2$, descendants are exactly $jd$ with $P^-(d)>P^+(j)$ and $d$ squarefree; the corresponding inverse entries are $\mu(d)$. For an isolated nonsquarefree column only its diagonal one contributes. Squaring and summing proves the identities.
\end{proof}

Two further counts will be useful. Distinct nonzero columns of $S_n$ have disjoint supports, so
\begin{equation}\label{def:orthocolumns}
 S_n^TS_n=\diag(d_1(n),\ldots,d_n(n)),\qquad
 \sum_jd_j(n)=Q(n)-1.
\end{equation}
Since all nonzero entries of $A_n$ and $B_n$ equal one,
\begin{equation}\label{def:frobenius}
 \norm{A_n}_F^2=n+Q(n)-1,\qquad
 \norm{B_n}_F^2=2n+Q(n)-2,\qquad \tr B_n=n.
\end{equation}

\section{Restricted M\"obius sums and their moments}\label{sec:arithmetic}

This section develops the arithmetic tools needed for the exceptional singular
value. All analytic cancellation comes from classical Mertens and smooth-number
estimates. The moment and support-constrained statements are consequences of
those inputs; no arithmetic priority claim is attached to their formulation.
Throughout, $n\to\infty$, and we write
\[
 L=\log n,\qquad s=s(n)=\sqrt{L\log L},\qquad
 b_p(n)=R(n/p,p),\qquad E_P(n)=\sum_{p\le n}b_p(n)^2.
\]
All sums indexed by $p$ or $q$ as primes below are explicitly so indicated when
there could be ambiguity. Arguments of counting functions may be real, with the
usual floor convention. Constants may depend on fixed compact intervals and on
fixed moment parameters, but never on a prime varying in the stated interval.

\subsection{Classical inputs and their ranges}

Let $\rho$ be the Dickman function, defined by $\rho(v)=1$ for $0\le v\le1$,
$\rho(v)=0$ for $v<0$, and
\[
 v\rho'(v)=-\rho(v-1)\qquad(v>1).
\]
We use the following classical estimates, recording the needed specialization
rather than assuming an asymptotic for $R$ itself. The ordinary Mertens bound is
\begin{equation}\label{arith:mertens}
 |M(z)|\le C z\exp\left\{-c\frac{(\log z)^{3/5}}
 {(\log\log z)^{1/5}}\right\}
 \le C z e^{-c(\log z)^{7/12}}\qquad(z\ge z_0).
\end{equation}
The second inequality follows after changing $c,C$, since
$(\log z)^{1/60}/(\log\log z)^{1/5}\to\infty$.
For a modern explicit statement implying the first inequality, including a
logarithmic prefactor absorbable into the exponential, see
Lee--Leong~\cite[Theorem~1.1, (12)]{leeleong}.

Here is the precise smooth-number input. Fix $0<c_-<c_+<\infty$ and $D>0$.
Suppose
\[
 c_-s\le t=\log y\le c_+s,\qquad
 |\log x-L|\le Ds,\qquad
 \log x-L^{9/10}\le\log z\le\log x.
\]
Uniformly over these choices, with $v=\log z/t$,
\begin{equation}\label{arith:smooth-input}
 \Psi(z,y)=z\left\{\rho(v)+(\gamma-1)\frac{\rho'(v)}t
 +O\left(\rho(v)\left(\frac{\log(v+1)}t\right)^2\right)\right\},
\end{equation}
where $\gamma$ is Euler's constant. We also use, uniformly as $v\to\infty$,
\begin{equation}\label{arith:dickman-input}
 -\rho'(v)\asymp\rho(v)\log(v+1),\qquad
 |\rho''(v)|\ll\rho(v)\log^2(v+1),\qquad
 \log\rho(v)=-(1+o(1))v\log v.
\end{equation}
Only arguments tending uniformly to infinity occur in the derivative estimates,
so the low-argument changes of differentiability are irrelevant.

For clarity, these are specializations of established smooth-number theory,
not additional hypotheses. McNew~\cite[(7), (13), (14), (17)]{mcnew}
records the Hildebrand first-order estimate, the de Bruijn--Saias expansion,
and the derivative estimates; see also Saias~\cite{saias} and the survey
of Hildebrand--Tenenbaum~\cite{hildebrandtenenbaum}.
The first-order statement alone would not suffice here.
In the notation of that expansion, take two terms (the derivative order
parameter $k=1$). The coefficient of $\rho'(v)/t$ is $\gamma-1$,
and the relative remainder after the two terms is of order
$(\log(v+1)/t)^2$ in the present range.
The source range contains
\[
 y>\exp\{(\log\log z)^{5/3+\varepsilon}\},\qquad
 y>(\log z)^{1+\varepsilon},\qquad y<z,
\]
with any fixed sufficiently small $\varepsilon>0$.
The extra boundary inequalities for $k=1$ in (13) are
\[
 \frac v2\ge\frac{\log t}{t},\qquad
 v-1\ge\frac{\log t}{t}.
\]
Indeed our substitutions give $\log z=L(1-o(1))$,
$t\asymp\sqrt{L\log L}$, and
$v\asymp\sqrt{L/\log L}$. All these inequalities therefore hold uniformly
with a diverging margin. The exponentially small relative error in the
Saias estimate is $o((\log(v+1)/t)^2)$ here. Since
$\rho'(v)/(t\rho(v))=O(\log(v+1)/t)=o(1)$, the multiplicative
remainder in the source yields exactly the additive remainder in
\eqref{arith:smooth-input}. In particular, no exclusion of integer $v$ is
needed in this range. The derivative and ratio estimates recorded there,
together with the defining Dickman equation, give
\eqref{arith:dickman-input}.

We also use the prime number theorem and Chebyshev's upper bound
$\pi(z)\ll z/\log z$. The series application additionally uses
Mertens' product bound
\[
 \prod_{p\le y}(1-1/p)^{-1}\ll\log(2y).
\]
The weaker estimate
\begin{equation}\label{arith:harmonic-product}
 \prod_{p\le y}(1-1/p)^{-1}
 =\exp\{O(\log\log(3y))\}
\end{equation}
already suffices for the moment proof. Partial summation of Chebyshev gives
\[
 \sum_{p\le y}p^{-1}=O(\log\log(3y)),
\]
and higher prime powers have a bounded total contribution to the logarithm
of the product.

\subsection{Two exact convolutions and a uniform upper bound}

Let $S_y(a)$ be the indicator of $P^+(a)\le y$, including $a=1$.
The coefficient identity
$(\mu*S_y)(d)=\mu(d)\one_{P^-(d)>y}$ follows prime by prime:
the local factor is $1$ at primes at most $y$ and $1-X$ at larger primes.
Summing its coefficients in either order gives
\begin{equation}\label{arith:convolutions}
 R(x,y)=\sum_{\substack{a\le x\\P^+(a)\le y}}M(x/a)
       =\sum_{k\le x}\mu(k)\Psi(x/k,y).
\end{equation}
These are finite identities and include the coefficient at $1$.
Replacing $M$ by $|M|$ in the first identity gives a positive majorant
which increases with $y$.

\begin{lemma}\label{arith:upper-local}
Fix $0<c_-<c_+<\infty$ and $D>0$. Uniformly for
$|\log x-L|\le Ds$ and $\log y=cs$ with $c_-\le c\le c_+$,
\[
 |R(x,y)|\le x\exp\left\{-\left(\frac1{2c}+o(1)\right)s\right\}.
\]
The same bound holds for the positive majorant just described.
\end{lemma}
\begin{proof}
Put $h=L^{9/10}$ and split the first sum in
\eqref{arith:convolutions} at $a=xe^{-h}$. By
\eqref{arith:mertens} with exponent $7/12$, the part below this point is
at most
\[
 Cxe^{-ch^{7/12}}\sum_{P^+(a)\le y}\frac1a
 \le x\exp\{-cL^{21/40}+O(\log L)\}=x e^{-\omega(s)}.
\]
Here and below $e^{-\omega(s)}$ means a quantity bounded by
$e^{-As}$ for every fixed $A>0$ once $n$ is sufficiently large.
For the remaining part use $|M(x/a)|\le x/a$ and Rankin's inequality.
With $u=\log x/t$, $t=\log y$, and $\alpha=\log u/t$, it is at most
\begin{equation}\label{arith:rankin}
 xe^{-\alpha(\log x-h)}
 \prod_{p\le y}(1-p^{-1+\alpha})^{-1}.
\end{equation}
Uniformly in this range, $u\to\infty$ and $0<\alpha\le1/4$ eventually.
Higher prime-power terms in the logarithm of the product are $O(1)$,
and Chebyshev partial summation bounds the prime terms by
\[
 O(1)+O(e^{\alpha t}/t)
 +O\left(\int_{\log2}^t\frac{e^{\alpha v}}v\,dv\right)
 =O(\log t+e^{\alpha t})=O(\log L+u)=o(s).
\]
To see the integral bound, split at $1/\alpha$. Below this point use a
bounded exponential; above it substitute $w=\alpha v$ and use
$\int_1^T e^w/w\,dw\ll e^T$. Finally $\alpha h=o(s)$ and
\[
 \alpha\log x=u\log u=\left(\frac1{2c}+o(1)\right)s
\]
uniformly on the fixed compact interval. Substitution in
\eqref{arith:rankin} completes the proof.
\end{proof}

\subsection{A relative asymptotic at the saddle scale}

\begin{lemma}\label{arith:bridge}
Under the ranges of Lemma~\ref{arith:upper-local}, put
$t=\log y$ and $u=\log x/t$. Then
\begin{equation}\label{arith:rough-asymptotic}
 R(x,y)=(1+o(1))\frac{x\rho'(u)}t
\end{equation}
uniformly. In particular it is negative for all sufficiently large $n$
throughout this range.
\end{lemma}
\begin{proof}
Set $a_0=7/12$, $K=\lfloor\exp((\log L)^2)\rfloor$,
$H=\exp(L^{9/10})$, and $\beta=\log(u+1)/t$.
Here $\beta\asymp\sqrt{\log L/L}$. For $0\le w\le L^{9/10}$,
all arguments between $u-w/t$ and $u$ are $(1-o(1))u$.
Integrating the logarithmic derivative bound in
\eqref{arith:dickman-input} gives
\begin{equation}\label{arith:dickman-shift}
 \rho(u-w/t)\le\rho(u)e^{C\beta w}.
\end{equation}

First truncate the second identity in \eqref{arith:convolutions} without
losing M\"obius cancellation. Interchanging its tail sums gives exactly
\begin{equation}\label{arith:tail-exact}
 T_K:=\sum_{K<k\le x}\mu(k)\Psi(x/k,y)
 =\sum_{\substack{a\le x/K\\P^+(a)\le y}}
       \{M(x/a)-M(K)\}.
\end{equation}
An endpoint with $x/a=K$ contributes zero, so this formula is unchanged by
floors. The contribution from $M(K)$ is bounded using
\eqref{arith:smooth-input} and \eqref{arith:dickman-shift} by
\[
 |M(K)|\Psi(x/K,y)
 \ll x\rho(u)\exp\{-c(\log K)^{a_0}+C\beta\log K\}.
\]
For the terms with $K<x/a<H$, split $x/a$ into dyadic intervals
$[q,2q)$ starting at $K$, cutting the last one at $H$ if necessary.
There are at most $\Psi(x/q,y)$ eligible integers $a$ in one interval.
The Mertens bound and the same two smooth estimates therefore bound its
contribution by
\[
 Cx\rho(u)\exp\{-c'(\log q)^{a_0}+C\beta\log q\}.
\]
Uniformly for $\log K\le v\le\log H$,
\[
 \beta v^{1-a_0}\le
 O\left(L^{-1/2+(9/10)(1-a_0)}\sqrt{\log L}\right)=o(1).
\]
Thus the positive term in the exponent is absorbed. Even counting all
$O(L)$ possible dyadic intervals, their sum is at most
$x\rho(u)\exp\{-c''(\log K)^{a_0}\}$ after reducing $c''$,
since $(\log K)^{a_0}\asymp(\log L)^{7/6}\gg\log L$.

For $x/a\ge H$, use \eqref{arith:harmonic-product} to obtain
\[
 \sum_{\substack{P^+(a)\le y\\x/a\ge H}}|M(x/a)|
 \ll x e^{-c(\log H)^{a_0}}
        \prod_{p\le y}(1-1/p)^{-1}
 \le x\exp\{-cL^{21/40}+O(\log L)\}.
\]
By \eqref{arith:dickman-input}, $|\log\rho(u)|=O(s)$; since
$21/40>1/2$, this last bound is $o(x\rho(u)\beta)$.
The previous bounds have the same property because
$(\log L)^{7/6}\gg|\log\beta|$. We have proved
\begin{equation}\label{arith:tail-small}
 T_K=o(x\rho(u)\beta)
\end{equation}
uniformly.

It remains to expand only the short initial sum $k\le K$.
Since $\beta\log K=o(1)$, Taylor's theorem and
\eqref{arith:dickman-input}--\eqref{arith:dickman-shift} give
\begin{align*}
 \rho(u-\log k/t)&=\rho(u)-\frac{\log k}{t}\rho'(u)
       +O\bigl(\rho(u)\beta^2(\log k)^2\bigr),\\
 \frac{\rho'(u-\log k/t)}t&=\frac{\rho'(u)}t
       +O\bigl(\rho(u)\beta^2\log k\bigr).
\end{align*}
Let
$A_0(K)=\sum_{k\le K}\mu(k)/k$ and
$A_1(K)=\sum_{k\le K}\mu(k)\log k/k$.
Applying \eqref{arith:smooth-input} and summing its errors absolutely yields
\begin{align}\label{arith:short-expansion}
 \sum_{k\le K}\mu(k)\Psi(x/k,y)
 &=x\rho(u)A_0(K)-\frac{x\rho'(u)}t A_1(K)\notag\\
 &\quad +(\gamma-1)\frac{x\rho'(u)}t A_0(K)
 +O\bigl(x\rho(u)\beta^2(1+\log K)^3\bigr).
\end{align}
The cubic factor bounds $\sum_{k\le K}(1+\log k)^2/k$.

For completeness, the two M\"obius moments here satisfy
\begin{equation}\label{arith:mobius-moments}
 A_0(K)=O(e^{-c_1(\log K)^{a_0}}),\qquad
 A_1(K)=-1+O(e^{-c_1(\log K)^{a_0}}).
\end{equation}
Indeed partial summation with \eqref{arith:mertens} proves convergence
of both infinite series and bounds their tails by a polynomial in
$\log K$ times $e^{-c(\log K)^{a_0}}$.
The same integrable bounds justify letting $\sigma\downarrow1$ in
$\sum\mu(k)k^{-\sigma}=1/\zeta(\sigma)$ and its derivative.
The simple pole of $\zeta$ at $1$ then gives values $0$ and $-1$,
respectively. Absorb the polynomial into a smaller $c_1$ to obtain
\eqref{arith:mobius-moments}.
Now $\beta(1+\log K)^3=o(1)$ and
$|\rho'(u)|/t\asymp\rho(u)\beta$.
Equations \eqref{arith:tail-small}--\eqref{arith:mobius-moments}
leave precisely $x\rho'(u)/t$ as the main term, proving the lemma.
\end{proof}

The two-term input is essential to this argument. A relative error of order
$\beta$ in the first-order smooth estimate would be as large as the surviving
term after $A_0(\infty)=0$ cancels the density term.

\subsection{Moments, the maximum, and prime concentration}

\begin{theorem}\label{arith:moments}
For every fixed real $q>1$,
\begin{equation}\label{arith:moment-law}
 \sum_{p\le n}|b_p(n)|^q
 =n^q\exp\{-(\sqrt{2q(q-1)}+o(1))s\}.
\end{equation}
Moreover
\begin{equation}\label{arith:maximum}
 \max_{p\le n}|b_p(n)|
 =n\exp\{-(\sqrt2+o(1))s\}.
\end{equation}
For every fixed $0<\delta<1$, there is $\eta_\delta>0$ such that
\begin{equation}\label{arith:concentration}
 \sum_{\substack{p\le n\\|\log p/s-1|>\delta}}b_p(n)^2
 \le n^2e^{-(2+\eta_\delta)s}
\end{equation}
for all sufficiently large $n$.
\end{theorem}
\begin{proof}
For a fixed compact interval $0<c_-\le\log p/s\le c_+$,
Lemma~\ref{arith:upper-local} applied to $x=n/p$, $y=p$ gives
\[
 |b_p(n)|\le\frac np
 \exp\left\{-\left(\frac{s}{2\log p}+o(1)\right)s\right\}.
\]
In a bin $e^{as}<p\le e^{bs}$ within this interval, it follows that
\begin{equation}\label{arith:bin}
 \sum_{e^{as}<p\le e^{bs}}|b_p(n)|^q
 \ll_q n^q\exp\left\{-\left((q-1)a+\frac q{2b}-o(1)\right)s\right\},
\end{equation}
since $\sum_{m>e^{as}}m^{-q}\ll_q e^{-(q-1)as}$.
Cover the compact interval by finitely many bins of mesh at most $\varepsilon$.
As the mesh decreases their costs approach
\[
 f_q(c)=(q-1)c+\frac q{2c},\qquad
 \min_{c>0}f_q(c)=\sqrt{2q(q-1)},\qquad
 c_q=\sqrt{\frac q{2(q-1)}}.
\]

The two noncompact prime ranges require separate bounds. For $p<e^{c_-s}$,
enlarge the smoothness parameter in the positive majorant of
\eqref{arith:convolutions} to $y_0=e^{c_-s}$, keeping $x=n/p$.
Lemma~\ref{arith:upper-local} gives a total at most
\[
 n^q\exp\{-(q/(2c_-)-o(1))s\}\sum_p p^{-q}.
\]
Choose $c_-$ sufficiently small. For $p>e^{c_+s}$ the bound
$|b_p(n)|\le n/p$ gives total
$O_q(n^qe^{-(q-1)c_+s})$; choose $c_+$ sufficiently large.
For each desired error first fix these endpoints and a sufficiently fine
mesh, and only then let $n\to\infty$.
This proves the upper bound in \eqref{arith:moment-law}.
The same argument pointwise minimizes $c+1/(2c)$ to $\sqrt2$.
Its small-prime and large-prime costs are $1/(2c_-)$ and $c_+$,
respectively, proving the upper bound in \eqref{arith:maximum}
without taking a nonuniform limit $q\to\infty$.

For the lower bounds fix any $c>0$ and take primes in
$[P,2P]$, $P=e^{cs}$. Lemma~\ref{arith:bridge} and
\eqref{arith:dickman-input} give uniformly in this block
\begin{equation}\label{arith:block-law}
 |b_p(n)|=n\exp\{-(c+1/(2c)+o(1))s\}.
\end{equation}
Indeed $u=\log(n/p)/\log p$ satisfies
$u\log u=(1/(2c)+o(1))s$, and logarithms of the derivative and
$1/\log p$ prefactors are $o(s)$.
The prime number theorem supplies
$\pi(2P)-\pi(P)=\exp\{(c+o(1))s\}$ primes.
Taking $c=c_q$ yields the matching moment lower bound, while taking
$c=1/\sqrt2$ and any prime in the block yields the maximum lower bound.

Finally set $q=2$ in the bin argument. The cost is $c+1/c$,
whose unique minimum is $2$ at $c=1$.
Outside $(1-\delta,1+\delta)$ its infimum exceeds $2$.
Choose the remote endpoints so that their costs also exceed $2$,
and choose a mesh small enough to preserve half the positive gap.
The same finite-bin estimates, absorbing their number into the gap,
prove \eqref{arith:concentration}.
\end{proof}

No assertion here is uniform in a varying $q$. In particular,
\eqref{arith:moment-law} at $q=2$ gives
$E_P(n)=n^2e^{-(2+o(1))s}$, and \eqref{arith:concentration}
then gives concentration relative to the total prime energy.

\subsection{Composite multipliers and the inverse-vector norm}

Recall the full energy and the inverse-vector norm from the finite identities:
\[
 E(n)=\sum_{\substack{2\le j\le n\\\mu(j)^2=1}}
       R(n/j,P^+(j))^2,\qquad W_n=\norm{w_n}_2,
\]
\begin{equation}\label{arith:norm-recalled}
 W_n^2=E(n)+(M(n)-1)^2+n-Q(n).
\end{equation}
Here $Q(n)=\sum_{j\le n}\mu(j)^2$ is a count, not its square root.

\begin{theorem}\label{arith:energy}
The prime and full energies and the inverse-vector norm satisfy
\begin{gather}
 E_P(n)=n^2e^{-(2+o(1))s},\qquad
 E(n)=n^2e^{-(2+o(1))s},\qquad W_n=ne^{-(1+o(1))s},
 \label{arith:energy-scale}\\
 \frac{E(n)}{E_P(n)}\longrightarrow\frac{15}{\pi^2},\qquad
 \frac{W_n^2}{E_P(n)}\longrightarrow\frac{15}{\pi^2}.
 \label{arith:energy-ratio}
\end{gather}
These conclusions do not give a multiplicative equivalent for $E_P(n)$ itself.
\end{theorem}
\begin{proof}
The prime-energy assertion has just been proved. Every squarefree $j\ge2$
has a unique factorization $j=ap$ with $p=P^+(j)>P^+(a)$ and $a$
squarefree, including $a=1$. Consequently
\begin{equation}\label{arith:multiplier-decomposition}
 E(n)=\sum_{a\ge1}\mu(a)^2
       \sum_{\substack{p>P^+(a)\\ap\le n}}R(n/(ap),p)^2.
\end{equation}
Put $A=e^{3s}$ and $\mathcal C_n=\{p:e^{s/2}\le p\le e^{2s}\}$.
The contribution of $a>A$, by $|R(x,p)|\le x$ for $x\ge1$, is at most
\[
 n^2\sum_{a>A}a^{-2}\sum_p p^{-2}=O(n^2/A)=o(E_P(n)).
\]

For $a\le A$ set $N_a=\lfloor n/a\rfloor$. The exact floor identity is
\begin{equation}\label{arith:floor}
 R(n/(ap),p)=R(N_a/p,p):
 \quad pd\le n/a\ \Longleftrightarrow\ pd\le N_a
 \quad\hbox{for integers }p,d.
\end{equation}
Uniformly over these multipliers,
\[
 N_a=(n/a)(1+O(A/n)),\qquad
 \log N_a=L-\log a+O(A/n),\qquad s(N_a)/s=1+o(1),
\]
and $\min_{a\le A}N_a\to\infty$.
For all sufficiently large $n$, the complement of $\mathcal C_n$ lies
outside the band $|\log p/s(N_a)-1|\le1/3$, uniformly in $a\le A$.
Apply \eqref{arith:concentration} at the integer $N_a$ and drop the
restriction $p>P^+(a)$. For one fixed $\eta>0$, the total contribution
from $a\le A$, $p\notin\mathcal C_n$ is bounded by
\[
 \sum_{a\le A}N_a^2e^{-(2+\eta)s(N_a)}
 \le n^2e^{-(2+\eta/2)s}\sum_{a\ge1}a^{-2}=o(E_P(n)).
\]
This invocation is uniform because the same theorem holds at every
sufficiently large integer, and all the $N_a$ tend to infinity together.

It remains to compare central terms multiplicatively. For $p\in\mathcal C_n$
put $t=\log p$, $u=\log(n/p)/t$, and $v=u-\log a/t$.
Lemma~\ref{arith:bridge} applies uniformly both at $(n/p,p)$ and at
$(n/(ap),p)$ because their first logarithms equal $L+O(s)$.
Equivalently one may use \eqref{arith:floor}: replacing $N_a$ by $n/a$
changes the Dickman argument by $O(A/(nt))$, which is negligible under
the logarithmic derivative bound below. Thus
\begin{equation}\label{arith:multiplier-ratio}
 \frac{R(n/(ap),p)}{R(n/p,p)}
 =(1+o(1))\frac1a\frac{\rho'(v)}{\rho'(u)}
 \qquad(a\le A,\ p\in\mathcal C_n)
\end{equation}
with a common relative error. The denominator is nonzero by
Lemma~\ref{arith:bridge}. All arguments from $v$ to $u$ tend uniformly
to infinity, and $0\le u-v\le6$.
Equation \eqref{arith:dickman-input} gives
\[
 \left|\frac{d}{dz}\log(-\rho'(z))\right|
 \ll\log(z+1),\qquad
 \left|\frac{\rho'(v)}{\rho'(u)}\right|
 \le\exp\left\{C\frac{\log L}{s}\log a\right\}.
\]
As $\log L/s\to0$, for any fixed $0<\varepsilon<1$ we obtain
\begin{equation}\label{arith:domination}
 R(n/(ap),p)^2\le C a^{-2+\varepsilon}R(n/p,p)^2
\end{equation}
uniformly in the truncated region for all sufficiently large $n$.
For each fixed $a$, \eqref{arith:multiplier-ratio} instead gives the
uniform equivalent $a^{-2}$ after squaring; also $p>P^+(a)$ throughout
$\mathcal C_n$ eventually.

Define the normalized central contribution explicitly by
\[
 \mathcal E_{a,n}=\mathbf1_{a\le A}\frac{\mu(a)^2}{E_P(n)}
 \sum_{\substack{p\in\mathcal C_n\\p>P^+(a),\ ap\le n}}
 R(n/(ap),p)^2.
\]
It is bounded by
$C\mu(a)^2a^{-2+\varepsilon}$, a summable sequence independent of $n$.
For each fixed $a$ it tends to $\mu(a)^2/a^2$, since central prime energy
is $(1+o(1))E_P(n)$. Dominated convergence, followed by the two discarded
tail estimates, proves
\[
 \frac{E(n)}{E_P(n)}\longrightarrow
 \sum_{a\ge1}\frac{\mu(a)^2}{a^2}
 =\prod_p(1+p^{-2})=\frac{\zeta(2)}{\zeta(4)}=\frac{15}{\pi^2}.
\]
Finally \eqref{arith:mertens} implies
$(M(n)-1)^2=o(E_P(n))$, since $L^{7/12}/s\to\infty$,
and $n-Q(n)\le n=o(E_P(n))$. The exact norm identity
\eqref{arith:norm-recalled} proves the remaining assertions.
\end{proof}

\subsection{Sharp support constraints and exceptional cutoffs}

For fixed $\kappa>0$ define
\[
 F_\kappa(n)=\sup_{\substack{|a_p|\le1\\
 \#\{p\le n:a_p\ne0\}\le\lfloor e^{\kappa s}\rfloor}}
 \left|\sum_{2\le m\le n}\mu(m)a_{P^-(m)}\right|.
\]
Weights may be complex and may depend on $n$. Grouping a squarefree
integer $m$ by its least prime $p$ gives the exact identity
\begin{equation}\label{arith:least-prime}
 \sum_{2\le m\le n}\mu(m)a_{P^-(m)}
 =-\sum_{p\le n}a_p R(n/p,p)=-\sum_{p\le n}a_p b_p(n),
\end{equation}
since $m=pd$ has $P^-(d)>p$ and $\mu(pd)=-\mu(d)$.
Alladi~\cite{alladi1982} treats the unrestricted bounded-weight problem;
here the support ceiling is part of the extremal question.

\begin{theorem}\label{arith:sparse}
For every fixed $\kappa>0$,
\begin{equation}\label{arith:sparse-law}
 F_\kappa(n)=n\exp\{-(g(\kappa)+o(1))s\},\qquad
 g(\kappa)=
 \begin{cases}
 \sqrt2-\kappa,&0<\kappa\le1/\sqrt2,\\
 1/(2\kappa),&\kappa\ge1/\sqrt2.
 \end{cases}
\end{equation}
\end{theorem}
\begin{proof}
If $K\le e^{\kappa s}$ weights are nonzero, H\"older's inequality and
\eqref{arith:moment-law} imply, for each fixed $q>1$,
\[
 \left|\sum_pa_pb_p(n)\right|
 \le K^{1-1/q}\left(\sum_p|b_p(n)|^q\right)^{1/q}
 \le n\exp\{-(\sqrt{2\theta}-\kappa\theta-o(1))s\},
 \qquad\theta=1-1/q.
\]
For $\kappa>1/\sqrt2$, choose the fixed
$\theta=1/(2\kappa^2)\in(0,1)$ to obtain saving $1/(2\kappa)$.
For $\kappa\le1/\sqrt2$, use directly $K\max_p|b_p(n)|$
and \eqref{arith:maximum} to obtain saving $\sqrt2-\kappa$.
This covers the common boundary without varying $q$.

For the lower bound when $\kappa<1/\sqrt2$, take
$\lfloor e^{\kappa s}\rfloor$ primes from
$[e^{s/\sqrt2},2e^{s/\sqrt2}]$. The prime number theorem guarantees
at least this many primes for all sufficiently large $n$.
Equation \eqref{arith:block-law} gives each selected coefficient magnitude
$n e^{-(\sqrt2+o(1))s}$. Choose $a_p$ to align their contributions
in \eqref{arith:least-prime}; their sum gives the required lower bound.
For $\kappa\ge1/\sqrt2$, take every prime in
$[e^{\kappa s},2e^{\kappa s}]$. Their number is asymptotic to
$e^{\kappa s}/(\kappa s)$, hence satisfies the support ceiling, and their
total magnitude by \eqref{arith:block-law} is
$n e^{-(1/(2\kappa)+o(1))s}$. This also proves the boundary case.
In fact the coefficients in either block are eventually negative by
Lemma~\ref{arith:bridge}, so indicator weights suffice for these lower
constructions. The supremum assertion does not claim sharpness for every
prescribed support.
\end{proof}

\begin{corollary}\label{arith:exceptional}
Let $N_\lambda(n)=\#\{p\le n:|b_p(n)|\ge ne^{-\lambda s}\}$.
For fixed $\lambda>\sqrt2$,
\[
 N_\lambda(n)=\exp\{(h(\lambda)+o(1))s\},\qquad
 h(\lambda)=\frac{\lambda+\sqrt{\lambda^2-2}}2.
\]
For fixed $\lambda<\sqrt2$ the count is eventually zero.
At $\lambda=\sqrt2$ we obtain only the upper bound
$N_{\sqrt2}(n)\le\exp\{(1/\sqrt2+o(1))s\}$.
\end{corollary}
\begin{proof}
Markov's inequality and \eqref{arith:moment-law} give, for fixed $q>1$,
\[
 N_\lambda(n)\le
 \exp\{(q\lambda-\sqrt{2q(q-1)}+o(1))s\}.
\]
For $\lambda>\sqrt2$ the minimum is attained at
$q=(1+\lambda/\sqrt{\lambda^2-2})/2$, and its value is $h(\lambda)$.
For the reverse bound choose a fixed $c$ strictly between the two roots
of $c+1/(2c)=\lambda$ and arbitrarily close to the larger root
$h(\lambda)$. Every prime in $[e^{cs},2e^{cs}]$ meets the threshold
eventually by \eqref{arith:block-law}, so the count is at least
$e^{(c+o(1))s}$. Let $c$ increase to $h(\lambda)$ after taking the
limit in $n$. The assertion for $\lambda<\sqrt2$ follows from
\eqref{arith:maximum}. At equality the displayed Markov bound holds
for each fixed $q$, and its cost decreases to $1/\sqrt2$ as
$q\to\infty$. Taking the infimum after the limsup proves the stated
upper bound; no uniform moment estimate or matching lower is asserted.
\end{proof}

\subsection{A weighted harmonic series on a thin prime support}

\begin{theorem}\label{arith:thin-series}
Fix $\kappa>0$ and a fixed set $\mathcal S$ of primes such that
\[
 N_{\mathcal S}(x):=\#(\mathcal S\cap[2,x])\le e^{\kappa s(x)}
\]
for every sufficiently large $x$. For fixed complex weights $|a_p|\le1$
supported on $\mathcal S$, the series in increasing order of $m$ satisfies
\begin{equation}\label{arith:series-value}
 \sum_{m\ge2}\frac{\mu(m)a_{P^-(m)}}m=0,
 \qquad
 \left|\sum_{m>X}\frac{\mu(m)a_{P^-(m)}}m\right|
 \le\exp\{-(g(\kappa)-o(1))s(X)\}.
\end{equation}
If $\log(N_{\mathcal S}(x)+1)=o(s(x))$, the tail bound holds with
$g(\kappa)$ replaced by $\sqrt2$.
\end{theorem}
\begin{proof}
Write $c_m=\mu(m)a_{P^-(m)}$ for $m\ge2$ and
$T(x)=\sum_{2\le m\le x}c_m$.
At every sufficiently large integer endpoint the support hypothesis and
Theorem~\ref{arith:sparse} apply, and passage to real $x$ by flooring
is harmless. For each fixed $0<\varepsilon<g(\kappa)/2$,
\[
 |T(x)|\le x e^{-(g(\kappa)-\varepsilon)s(x)}
\]
eventually. For every fixed $c>0$, substitution $v=\log t$ and then
$z=\sqrt{v\log v}$ shows
\begin{align*}
 \int_X^\infty e^{-cs(t)}\frac{dt}{t}
 &=\int_{\log X}^\infty e^{-c\sqrt{v\log v}}\,dv\\
 &\le\int_{s(X)}^\infty 2z e^{-cz}\,dz
 =O_c((s(X)+1)e^{-cs(X)}),
\end{align*}
using $dv/dz=2z/(\log v+1)\le2z$ eventually.
Partial summation now proves convergence and bounds the tail by
$|T(X)|/X+\int_X^\infty|T(t)|t^{-2}\,dt$.
Absorbing the polynomial factor into an arbitrarily small exponential
loss proves the tail bound in \eqref{arith:series-value}.

Convergence alone does not identify the value. For real $\sigma>1$,
absolute convergence allows grouping by the least prime and Euler products:
\begin{equation}\label{arith:series-euler}
 \sum_{m\ge2}\frac{c_m}{m^\sigma}
 =-\frac1{\zeta(\sigma)}
   \sum_{p\in\mathcal S}\frac{a_p}{p^\sigma}
      \prod_{\substack{q\le p\\q\ \mathrm{prime}}}(1-q^{-\sigma})^{-1}.
\end{equation}
The support assumption implies $N_{\mathcal S}(x)\le x^{1/2}$
eventually. Partial summation of this estimate gives
$\sum_{p\in\mathcal S}\log(2p)/p<\infty$: the integral of
$x^{1/2}\log(2x)/x^2$ converges, and its boundary term tends to zero.
For $\sigma\ge1$, Mertens' product bound dominates the absolute value
of the $p$-summand on the right of \eqref{arith:series-euler} by
$C\log(2p)/p$. Thus the prime sum remains bounded as
$\sigma\downarrow1$, whereas $1/\zeta(\sigma)\to0$.
The Dirichlet-series limit is zero. Finally partial summation writes the
left side as
$\sigma\int_1^\infty T(t)t^{-\sigma-1}\,dt$.
The bound already proved makes $|T(t)|/t^2$ integrable, so dominated
convergence for $1<\sigma\le2$ identifies its limit with the convergent
series at $\sigma=1$. Its value is therefore zero.

Under the smaller-support assumption,
\eqref{arith:least-prime} and \eqref{arith:maximum} directly give
$|T(x)|\le x e^{-(\sqrt2-o(1))s(x)}$.
The support still satisfies the preceding polynomial counting bound,
and the identical integral and Dirichlet-series arguments yield the
last assertion.
\end{proof}

The theorem concerns the naturally ordered series, not absolute convergence
of its terms, and supplies no explicit numerical onset. The unweighted
zero-value assertion for a prime set of density zero is already contained in
the density formula of Kural--McDonald--Sah~\cite[Theorem~1.1]{kms}.
The formulation here records a quantitative tail for fixed thin supports
and bounded complex weights; it does not assert that this rate has no prior
occurrence in the literature.

\section{The singular spectrum}\label{sec:spectral}

We use the definitions of Section~\ref{sec:definitions}. All vector norms are Euclidean and all matrix norms without a subscript are induced operator norms. For a matrix or a compact operator $T$, write $s_r(T)$ for its $r$th largest singular value, counting multiplicity and adjoining zeros when appropriate. Thus $s_r(B_n)=\sigma_r(B_n)$. For a real symmetric matrix, its inertia is the triple consisting of its numbers of positive, negative, and zero eigenvalues. Transposition is denoted by a superscript $T$; for vectors $x,y$ in the real Hilbert space $\ell^2(\mathbb N)$, the notation $xy^T$ means the operator $h\mapsto x\langle y,h\rangle$.

\subsection{Exact unit singular values and a finite reducing core}

Let $\mathcal P_n$ be the set of vertices having a child and put $k_n=|\mathcal P_n|$. A subspace reduces a matrix if it and its orthogonal complement are invariant under both the matrix and its transpose. The next statement locates the entire part of the matrix that can have nonunit singular values.

\begin{theorem}\label{spec:core}
The matrix $A_n$ has exactly $k_n$ singular values greater than one, $k_n$ less than one, and $n-2k_n$ equal to one. For $n\ge4$, the corresponding counts for $B_n$ are
\[
 k_n+1,\qquad k_n,\qquad n-2k_n-1.
\]
A zero singular value is included in the count below one. There are reducing subspaces of dimensions $2k_n$ and $2k_n+1$, respectively, outside which the matrices act as the identity; their dimensions equal the numbers of nonunit singular values. For every fixed $H>0$,
\begin{equation}\label{spec:parentbound}
 k_n=O_H\!\left(\frac{n}{(\log n)^H}\right).
\end{equation}
\end{theorem}

\begin{proof}
Suppress the index $n$. Distinct nonzero columns of $S$ have disjoint supports, because each vertex has at most one parent. Every active parent $a$ has a leaf child. Indeed, take the largest prime $p\le n/a$. Since $a$ has a child, $p>P^+(a)$; and the next prime after $p$ exceeds $n/a$, so $ap$ cannot itself have a child within the truncation. Such leaf children are distinct for distinct parents.

Set $E=\spanop\{e_a:a\in\mathcal P_n\}$. In the decomposition $E\oplus E^\perp$ the matrix has the form
\[
 S=\begin{pmatrix}D&0\\ F&0\end{pmatrix}.
\]
The columns of $F$ have disjoint nonempty supports, by the leaf observation, so $F$ has full column rank $k_n$. Let $\mathcal F=\operatorname{range}F$. Use the normalized columns of $F$ as an orthonormal basis of $\mathcal F$. In $V=E\oplus\mathcal F=E+\operatorname{range}S$ we then have
\[
 S|_V=\begin{pmatrix}D&0\\ R&0\end{pmatrix},\qquad
 A|_V=\begin{pmatrix}I+D&0\\ R&I\end{pmatrix},
\]
where $R$ is diagonal with strictly positive entries: its entry for $a$ is the square root of the number of leaf children of $a$. Both $S$ and $S^T$ vanish on $U=V^\perp$. Thus $V$ reduces $A$, and $A$ is the identity on $U$.

On $V$, its Gram displacement is
\[
 A^TA-I=\begin{pmatrix}T_0&R\\ R&0\end{pmatrix},\qquad
 T_0=D+D^T+D^TD+R^2.
\]
The invertible substitution $y'=y+\tfrac12R^{-1}T_0x$ changes the quadratic form $x^TT_0x+2x^TRy$ to $2x^TRy'$. The latter has inertia $(k_n,k_n,0)$, as is seen by subsequently setting $z=Ry'$ and diagonalizing $2x^Tz$ with $x+z$ and $x-z$. The displacement is zero on $U$. Its eigenvalues are the squared singular values of $A$ minus one, proving the assertions for $A$.

The first rows of $A$ and $B$ are $e_1^T$ and $\one^T$, respectively, so
\begin{equation}\label{spec:gramdisplacement}
 B^TB-I=(A^TA-I)+\one\one^T-e_1e_1^T.
\end{equation}
For $n\ge4$, the root belongs to $E$ and $z=\operatorname{proj}_U\one$ is nonzero. Indeed, the nonsquarefree coordinate $4$ is isolated, so $e_4\in U$ and $\langle\one,e_4\rangle=1$. Put $b=\operatorname{proj}_V\one$ and $\rho=\norm{z}>0$. Relative to $V\oplus\spanop\{z/\rho\}$, the displacement in \eqref{spec:gramdisplacement} is
\[
 \begin{pmatrix}
 (A^TA-I)|_V-e_1e_1^T+bb^T&\rho b\\
 \rho b^T&\rho^2
 \end{pmatrix}.
\]
Eliminating its off-diagonal block by congruence leaves the positive scalar $\rho^2$ and the Schur complement $(A^TA-I)|_V-e_1e_1^T$. The subtraction changes only the $E\times E$ block. Its other blocks remain $R,R,0$, so the preceding quadratic-form argument gives inertia $(k_n,k_n,0)$ for this complement. The full displayed matrix therefore has inertia $(k_n+1,k_n,0)$.

The subspace $V_B=V\oplus\spanop\{z\}$ reduces $B$. To check this directly, $B-I=S+e_1(\one-e_1)^T$ and its transpose both vanish on $U\cap\one^\perp=V_B^\perp$. Thus this complement supplies exactly $n-2k_n-1$ unit singular directions. If $B$ is singular, its zero Gram eigenvalue contributes $-1$ to the displacement and causes no exception to the count.

It remains to bound the number of parents. For $y\ge2$, any parent with $P^+(a)>y$ satisfies $a<n/y$, since it has a child $ap$ with $p>P^+(a)$. Consequently
\[
 k_n\le n/y+\#\{a\le n:a\text{ squarefree},\ P^+(a)\le y\}.
\]
Take $y=(\log n)^H$ and choose a fixed $\alpha\in(0,1)$ such that $H(1-\alpha)<1$. The second term is at most
\[
 n^\alpha\prod_{p\le y}(1+p^{-\alpha})
 \le n^\alpha\exp\!\left(\sum_{m\le y}m^{-\alpha}\right)
 \le n^\alpha\exp\bigl(O_\alpha(y^{1-\alpha})\bigr)
 =n^{\alpha+o(1)}.
\]
This is $o(n/(\log n)^H)$ and proves \eqref{spec:parentbound}.
\end{proof}

\subsection{Uniform comparison with the parent degrees}

Sort the parent degrees as $d_{(1)}(n)\ge\cdots\ge d_{(n)}(n)$, including zero degrees. Write $Q(n)=\sum_{j\le n}\mu(j)^2$. The estimates below have absolute errors independent of the index and require no cancellation in M\"obius sums.

\begin{theorem}\label{spec:upper}
For $n\ge2$,
\begin{align}
 \left|\sigma_r(A_n)-\sqrt{d_{(r)}(n)}\right|&\le1
 &&(1\le r\le n),\label{spec:Adegree}\\
 \left|\sigma_{r+1}(B_n)-\sqrt{d_{(r)}(n)}\right|
 &\le1+2\sqrt{\frac{Q(n)-1}{n}}<3
 &&(1\le r\le n-1).\label{spec:Bdegree}
\end{align}
Furthermore $\sigma_1(B_n)=\sqrt n+O(1)$. In particular, for any sequence $r=r(n)\le n-1$ such that $d_{(r(n))}(n)\to\infty$,
\[
 \frac{\sigma_{r(n)+1}(B_n)}{\sqrt{d_{(r(n))}(n)}}\longrightarrow1.
\]
If $a_1=1,a_2=2,a_3=3,a_4=5,\ldots$ are the squarefree positive integers in increasing order, then for every fixed $r\ge1$,
\begin{equation}\label{spec:upperfixed}
 \sigma_r(A_n)\sim\sqrt{\frac{n}{a_r\log n}},\qquad
 \sigma_{r+1}(B_n)\sim\sqrt{\frac{n}{a_r\log n}}.
\end{equation}
The final assertion uses the classical prime number theorem.
\end{theorem}

\begin{proof}
Orthogonality of the columns gives $S_n^TS_n=\diag(d_j(n))$. Thus $s_r(S_n)=\sqrt{d_{(r)}(n)}$ and \eqref{spec:Adegree} follows from the singular-value perturbation inequality $|s_r(X)-s_r(Y)|\le\norm{X-Y}$.

For the bordered matrix put $h=\one/\sqrt n$, $P=I-hh^T$, and
\[
 \eta=\norm{S_nh}=\sqrt{(Q(n)-1)/n},\qquad
 T=\sqrt n\,e_1h^T+S_nP.
\]
The first summand acts from $\spanop\{h\}$ to $\spanop\{e_1\}$; the second acts from $h^\perp$ to $e_1^\perp$. Thus these are orthogonal domain and range decompositions. The singular values of $T$ consist of $\sqrt n$ and the $n-1$ singular values of the restriction of $S_nP$ between those complements. The root has degree $\pi(n)$ and every other degree is at most $\pi(n)$, so
\[
 \norm{S_nP}\le\norm{S_n}=\sqrt{\pi(n)}<\sqrt n.
\]
Consequently $s_1(T)=\sqrt n$ and $s_{r+1}(T)=s_r(S_nP)$ for $1\le r\le n-1$, where the full matrix $S_nP$ has one additional zero singular value. Since
\[
 B_n=\sqrt n\,e_1h^T+(I-e_1e_1^T)+S_n,
 \qquad \norm{B_n-T}\le1+\eta,
\]
we obtain $|\sigma_1(B_n)-\sqrt n|\le1+\eta<2$. Applying the same inequality at index $r+1$, followed by $\norm{S_nP-S_n}=\eta$, proves \eqref{spec:Bdegree}. Dividing that bound by a degree tending to infinity proves its growing-index consequence.

For completeness, fixed-coordinate degree asymptotics must be justified after sorting. For squarefree $a$,
\[
 d_a(n)=\max\{0,\pi(n/a)-\pi(P^+(a))\},
\]
with $P^+(1)=1$, while nonsquarefree coordinates have degree zero. The prime number theorem gives
\[
 \frac{\log n}{n}d_a(n)\longrightarrow\frac{\mu(a)^2}{a}
 \quad\text{for every fixed }a.
\]
View the corresponding diagonal matrices as operators on $\ell^2(\mathbb N)$ by extending them by zero. These converge in operator norm to the diagonal operator with entries $\mu(a)^2/a$. Indeed, for $J<a\le\sqrt n$, the classical Chebyshev bound $\pi(x)\le Cx/\log x$ gives
\[
 \frac{\log n}{n}d_a(n)\le\frac{2C}{a}\le\frac{2C}{J},
\]
and for $a>\sqrt n$ the bound $d_a(n)\le n/a$ gives an upper bound $\log n/\sqrt n$. The limiting diagonal tail is at most $1/J$. Combine these bounds with convergence on the first $J$ coordinates, then let $J$ increase. The limiting diagonal operator is compact and has positive eigenvalues $1/a_r$ in decreasing order. The min--max principle, or its direct application to diagonal matrices, therefore gives
\[
 \frac{\log n}{n}d_{(r)}(n)\longrightarrow\frac1{a_r}.
\]
The bounded errors in \eqref{spec:Adegree}--\eqref{spec:Bdegree} are negligible on these diverging scales, proving \eqref{spec:upperfixed}.
\end{proof}

\subsection{An explicit compact inverse-Gram limit}

Hilberdink~\cite{hilberdink} previously applied compact Gram limits to arithmetic matrices. Here the ancestor relation determines a different explicit kernel. Recall $c_0=1/\zeta(2)$ and, for squarefree $k$, define
\begin{equation}\label{spec:density}
 \delta_k=\frac{c_0}{k}\prod_{p\le P^+(k)}\frac{p}{p+1}.
\end{equation}
The product is empty for $k=1$, so $\delta_1=c_0$. Define the symmetric infinite matrix $K$ by zero on every nonsquarefree row or column, and on squarefree indices by
\begin{equation}\label{spec:kernel}
 K_{jk}=\begin{cases}
 \mu(k/j)\delta_k,&j\preceq k,\\
 \mu(j/k)\delta_j,&k\preceq j,\\
 0,&\text{otherwise}.
 \end{cases}
\end{equation}
The first two formulas agree when $j=k$.

An operator $T$ on $\ell^2(\mathbb N)$ is Hilbert--Schmidt if
$\norm{T}_{\mathrm{HS}}^2=\sum_{j,k}|\langle e_j,Te_k\rangle|^2<\infty$.
Such an operator is bounded, its operator norm is at most its Hilbert--Schmidt norm, and truncation to finitely many coordinates converges in operator norm; in particular it is compact. A self-adjoint operator is positive if $\langle x,Tx\rangle\ge0$ for every $x$. Positive compact operators have decreasing nonnegative eigenvalues, with their positive eigenvalues repeated according to multiplicity.

\begin{theorem}\label{spec:compact}
The kernel \eqref{spec:kernel} defines a positive Hilbert--Schmidt operator of infinite rank. With matrices extended by zero outside their first $n$ coordinates,
\begin{equation}\label{spec:Gramlimit}
 G_n:=\frac{C_n^TC_n}{n}\longrightarrow K
 \quad\text{in Hilbert--Schmidt norm}.
\end{equation}
For every fixed $r\ge1$,
\[
 \frac{s_r(C_n)}{\sqrt n}\longrightarrow\sqrt{\lambda_r(K)},
 \qquad
 \frac{\norm{C_n}}{\sqrt n}\longrightarrow\sqrt{\norm{K}},
\]
where $\lambda_r(K)$ lists the positive eigenvalues. In particular, the second limit specifies an exact positive constant without a numerical approximation.
\end{theorem}

\begin{proof}
The inverse formula in Proposition~\ref{def:inverse} says that on squarefree indices $(C_n)_{ij}=\mu(i/j)$ when $j\preceq i$, and is zero otherwise. A nonsquarefree coordinate contributes only its diagonal entry one. Two squarefree vertices have common descendants precisely when they are comparable, because the ancestors of a vertex form a single chain. If $j\preceq k$, these common descendants are exactly the descendants of $k$. For every such descendant $i$,
$\mu(i/j)\mu(i/k)=\mu(k/j)$. It follows that
\begin{equation}\label{spec:exactGram}
 (C_n^TC_n)_{jk}=\mu(k/j)D_k(n)\qquad(j\preceq k).
\end{equation}
The other entries follow by symmetry, except for the isolated nonsquarefree diagonal, which is one.

For fixed squarefree $k$, its descendants are $kd$ where $d\le n/k$ is squarefree and has no prime factor at most $P^+(k)$. We give the density argument explicitly. Let $y=P^+(k)$ and $P_y=\prod_{p\le y}p$. The identity $\mu^2(d)=\sum_{b^2\mid d}\mu(b)$ implies
\[
 \sum_{d\le x}\mu^2(d)\mathbf1_{(d,P_y)=1}
 =\sum_{\substack{b\le\sqrt x\\(b,P_y)=1}}\mu(b)
   \sum_{\substack{m\le x/b^2\\(m,P_y)=1}}1.
\]
For fixed $b$, the inner count divided by $x$ tends to
$b^{-2}\prod_{p\le y}(1-1/p)$. To interchange the limit and the sum, the absolute contribution of $b>B$, divided by $x$, is at most $\sum_{b>B}b^{-2}$, uniformly in $x$. Letting $B$ increase proves that the density is
\[
 \prod_{p\le y}(1-1/p)
 \sum_{(b,P_y)=1}\frac{\mu(b)}{b^2}
 =\prod_{p\le y}(1-1/p)\prod_{p>y}(1-p^{-2}).
\]
Replacing $x$ by $n/k$ gives $D_k(n)/n\to\delta_k$, hence entrywise convergence in \eqref{spec:Gramlimit}.

We now prove a uniform summable tail, which is essential for the operator conclusion. Always $D_k(n)\le n/k$. The ancestors of squarefree $k$ number $1+\#\{p:p\mid k,\ p\text{ prime}\}\le1+\log_2k$. Thus all squarefree entries with maximum index $k$ contribute at most
\[
 \frac{2(1+\log_2k)}{k^2}
\]
to the squared Hilbert--Schmidt norm of $G_n$. The same bound holds for $K$, since $\delta_k\le1/k$. Therefore the squared norm of the squarefree part outside the first $J\times J$ block is at most
\begin{equation}\label{spec:HStail}
 2\sum_{k>J}\frac{1+\log_2k}{k^2},
\end{equation}
uniformly in $n$. The nonsquarefree diagonal of $G_n$ contributes at most $n/n^2=1/n$. Finite-block entrywise convergence, \eqref{spec:HStail}, and the triangle inequality prove Hilbert--Schmidt convergence. The same estimate constructs $K$ as a Hilbert--Schmidt operator. Each $G_n$ is positive; operator-norm convergence shows that $K$ is positive as well.

On any finite set of distinct prime coordinates, the compression of $K$ is the positive diagonal matrix with entries $\delta_p$, because distinct primes are incomparable. Such compressions have arbitrarily large rank, so $K$ has infinite rank. Finally the min--max principle for positive compact operators gives
$|\lambda_r(G_n)-\lambda_r(K)|\le\norm{G_n-K}$ for each fixed $r$. Since $\lambda_r(G_n)=s_r(C_n)^2/n$, this proves the stated limits.
\end{proof}

\subsection{The escaping direction and two projections}

Define the positive compact operator
\begin{equation}\label{spec:Ldefinition}
 L=K-\frac{(Ke_1)(Ke_1)^T}{c_0}.
\end{equation}
We next prove both its positivity and the projection limit needed for the border. For any nonzero real vector $z$, write
\[
 \mathsf P_z=I-\frac{zz^T}{\norm{z}^2};
\]
this is the orthogonal projection onto the vectors perpendicular to $z$. For a unit vector the denominator is one. Put
\[
 D_n=\frac{C_n}{\sqrt n},\qquad
 a_n=\frac{\muvec_n}{\sqrt{Q(n)}},\qquad
 v_n=\frac{w_n}{W_n}.
\]
These are well defined for $n\ge2$. The symbols $a_n,v_n$ in this subsection denote unit vectors, distinct from the scalar enumeration $a_r$ in Theorem~\ref{spec:upper}.

\begin{lemma}\label{spec:escape}
After zero extension to $\ell^2(\mathbb N)$, $v_n$ converges weakly to zero: $\langle v_n,x\rangle\to0$ for every fixed $x\in\ell^2$. Moreover
\begin{equation}\label{spec:projectedlimit}
 (\mathsf P_{a_n}D_n\mathsf P_{v_n})^T
 (\mathsf P_{a_n}D_n\mathsf P_{v_n})\longrightarrow L
 \quad\text{in operator norm}.
\end{equation}
The operator $L$ is positive, compact, and of infinite rank.
\end{lemma}

\begin{proof}
We first justify weak escape using the arithmetic estimates rather than normalization alone. For fixed $y$, the rough M\"obius sum satisfies
\begin{equation}\label{spec:smoothconvolution}
 R(x,y)=\sum_{\substack{b\le x\\P^+(b)\le y}}M(x/b).
\end{equation}
Indeed, convolving $\mu$ with the indicator of the $y$-smooth integers cancels its Euler factor at every prime at most $y$, leaving exactly $\mu(d)\mathbf1_{P^-(d)>y}$. This coefficient identity also proves \eqref{spec:smoothconvolution} by finite summation, without convergence assumptions about Dirichlet series.

Apply the classical estimate \eqref{arith:mertens} to the terms $b\le\sqrt x$. Their absolute sum is at most
\[
 C_yx\exp(-c'(\log x)^{7/12}),
\]
because $\sum_{P^+(b)\le y}b^{-1}=\prod_{p\le y}(1-p^{-1})^{-1}<\infty$. For $b>\sqrt x$, use $|M(x/b)|\le x/b$ and
\[
 \sum_{\substack{b>\sqrt x\\P^+(b)\le y}}b^{-1}
 \le x^{-1/4}\sum_{P^+(b)\le y}b^{-1/2}=O_y(x^{-1/4}).
\]
Thus
\begin{equation}\label{spec:fixedrough}
 R(x,y)=O_y\bigl(xe^{-c'(\log x)^{7/12}}+x^{3/4}\bigr).
\end{equation}
By the exact coordinates of $w_n$ in Section~\ref{sec:definitions}, its first coordinate is $M(n)-1$, a fixed squarefree coordinate $j\ge2$ is $R(n/j,P^+(j))$, and a nonsquarefree coordinate is one. Theorem~\ref{arith:energy} gives
$W_n=n\exp(-(1+o(1))s(n))$, where $s(n)=\sqrt{\log n\log\log n}$.
Since $(\log n)^{7/12}/s(n)\to\infty$, \eqref{arith:mertens} and \eqref{spec:fixedrough} show that every fixed coordinate of $w_n$ is $o(W_n)$. Therefore each coordinate of $v_n$ tends to zero. Its norm equals one; approximating any $x\in\ell^2$ by a vector supported on finitely many coordinates proves the asserted weak convergence.

Since $\muvec_n=C_ne_1$ and $(G_n)_{11}=Q(n)/n\to c_0$, a direct multiplication gives
\begin{equation}\label{spec:Jdefinition}
 J_n:=D_n^T\mathsf P_{a_n}D_n
 =G_n-\frac{(G_ne_1)(G_ne_1)^T}{(G_n)_{11}}
 \longrightarrow L
\end{equation}
in operator norm. Here and below a finite matrix such as $J_n$ is extended by zero. This also proves positivity of $L$, since $J_n$ is positive. Compactness follows either from \eqref{spec:Ldefinition} or from compactness of $K$ and the finite rank of the subtracted term.

Compactness sends bounded weakly null sequences to norm-null sequences. In this case the fact can be checked directly: approximate $L$ in norm by finite-coordinate operators, for which coordinatewise convergence gives the assertion. Thus $\norm{Lv_n}\to0$, and \eqref{spec:Jdefinition} yields $\norm{J_nv_n}\to0$. Expanding the two factors of the projection shows
\[
 \norm{\mathsf P_{v_n}J_n\mathsf P_{v_n}-J_n}
 \le3\norm{J_nv_n}\longrightarrow0.
\]
When this identity is interpreted on $\ell^2$, the projection is $I-v_nv_n^T$ on that space; its multiplication with the zero-extended $J_n$ agrees exactly with the zero extension of the finite product. This proves \eqref{spec:projectedlimit}.

Finally, on $m$ distinct prime coordinates, $K$ compresses to a positive diagonal matrix. The subtraction in \eqref{spec:Ldefinition} has rank at most one. On the subspace perpendicular to its defining vector the diagonal quadratic form remains strictly positive, furnishing at least $m-1$ positive directions. Because $L$ is positive, this proves that it has at least $m-1$ positive eigenvalues. Let $m$ increase. In particular its positive eigenvalues $\lambda_r(L)$ are defined and strictly positive for every fixed $r$, and \eqref{spec:projectedlimit} implies
\begin{equation}\label{spec:projectedsingular}
 s_r(\mathsf P_{a_n}D_n\mathsf P_{v_n})
 \longrightarrow\sqrt{\lambda_r(L)}.
\end{equation}
\end{proof}

\subsection{Deflation and every fixed nonexceptional lower singular value}

We record the elementary singular-value lemma with its dimension-independent bounds. It explains why the projections must occur on both sides.

\begin{lemma}\label{spec:deflation}
Suppose $\norm{D_n}\le C$, $a_n,v_n$ are unit vectors, and $|t_n|\to\infty$. For each fixed $r$, in dimensions tending to infinity,
\[
 s_{r+1}(D_n-t_na_nv_n^T)
 -s_r(\mathsf P_{a_n}D_n\mathsf P_{v_n})\longrightarrow0.
\]
\end{lemma}

\begin{proof}
Choose orthonormal bases beginning with $a_n$ in the range and $v_n$ in the domain. The perturbed matrix becomes
\[
 X_n=\begin{pmatrix}\alpha_n-t_n&b_n\\c_n&E_n\end{pmatrix},
\]
where $|\alpha_n|$, $\norm{b_n}$, $\norm{c_n}$, and $\norm{E_n}$ are at most $C$. Set $z_n=\alpha_n-t_n$, which is nonzero for all sufficiently large $n$. Direct multiplication gives
\[
 \begin{pmatrix}1&0\\-c_n/z_n&I\end{pmatrix}
 X_n
 \begin{pmatrix}1&-b_n/z_n\\0&I\end{pmatrix}
 =\begin{pmatrix}z_n&0\\0&E_n-c_nb_n/z_n\end{pmatrix}.
\]
Each multiplier and its inverse has norm at most $1+C/|z_n|$. Singular values of a product with invertible left and right factors are bounded above by the product of their norms times the original singular value, and bounded below by the reciprocal product of the inverse norms times that singular value. Thus corresponding singular values before and after this multiplication differ by a factor between $(1+C/|z_n|)^{-2}$ and $(1+C/|z_n|)^2$.

The first diagonal entry has modulus tending to infinity; the lower block is bounded uniformly. All singular values after the first are therefore bounded, so these multiplicative errors tend to zero also in absolute value at each such index. Further,
$\norm{c_nb_n/z_n}\le C^2/|z_n|\to0$. The singular-value perturbation inequality now gives
$s_{r+1}(X_n)-s_r(E_n)\to0$. The matrix $\mathsf P_{a_n}D_n\mathsf P_{v_n}$ has block form $\diag(0,E_n)$ in the same bases. Its first $n-1$ singular values are those of $E_n$, with a further zero appended, which proves the claim.
\end{proof}

\begin{theorem}\label{spec:bottom}
For every fixed integer $r\ge1$,
\begin{equation}\label{spec:bottomlimit}
 \sqrt n\,\sigma_{n-r}(B_n)\longrightarrow\lambda_r(L)^{-1/2}
 \qquad(n\to\infty).
\end{equation}
The limit is along all integers, including those with $M(n)=0$. In particular,
\[
 \sigma_{n-1}(B_n)\sim\frac1{\sqrt{n\norm{L}}}.
\]
\end{theorem}

\begin{proof}
First suppose $M(n)\ne0$. The exact inverse from Proposition~\ref{def:inverse} gives
\[
 \frac{B_n^{-1}}{\sqrt n}
 =D_n-t_na_nv_n^T,\qquad
 t_n=\frac{\sqrt{Q(n)}W_n}{M(n)\sqrt n}.
\]
By Theorem~\ref{spec:compact}, $\norm{D_n}$ is bounded. Moreover $Q(n)/n\to c_0$, while Theorem~\ref{arith:energy} and \eqref{arith:mertens} imply $W_n/|M(n)|\to\infty$ along the nonzero indices. Hence $|t_n|\to\infty$. Lemma~\ref{spec:deflation}, followed by \eqref{spec:projectedsingular}, yields
\[
 \frac{s_{r+1}(B_n^{-1})}{\sqrt n}
 \longrightarrow\sqrt{\lambda_r(L)}.
\]
For an invertible matrix with descending singular values,
$s_{r+1}(B_n^{-1})=1/\sigma_{n-r}(B_n)$; this proves \eqref{spec:bottomlimit} on these indices.

Now suppose $M(n)=0$. The rank-one perturbation of the invertible $A_n$ has rank at least $n-1$, and its nonzero kernel vector $\muvec_n$ shows that its rank is exactly $n-1$. Direct use of $B_n=A_n+e_1u^T$, $u=\one-e_1$, gives
\[
 B_n\muvec_n=0,\qquad w_n^TB_n=0,\qquad
 B_nC_n=I+e_1w_n^T.
\]
For the left-kernel identity, $w_n^TA_n=u^T$ and $w_n^Te_1=M(n)-1=-1$. Thus $a_n$ and $v_n$ span the right and left kernels, respectively. It follows that
\[
 B_n(\mathsf P_{a_n}C_n\mathsf P_{v_n})=\mathsf P_{v_n}.
\]
The operator in parentheses maps into $a_n^\perp$, vanishes on $v_n$, and inverts $B_n$ from $v_n^\perp$ onto $a_n^\perp$. This is exactly the Moore--Penrose inverse: by definition, it is the inverse on those orthogonal complements and is zero on the left kernel. Consequently
\[
 \frac{B_n^\dagger}{\sqrt n}=\mathsf P_{a_n}D_n\mathsf P_{v_n}.
\]
Since $\sigma_n(B_n)=0$ and all other singular values are positive, its $r$th largest singular value is $1/\sigma_{n-r}(B_n)$. Equation~\eqref{spec:projectedsingular} proves the same limit along zero-Mertens indices. The two cases together prove the all-integer assertion.
\end{proof}

\subsection{The exceptional minimum and conditioning}

The preceding theorem omits exactly the singular direction carrying $M(n)$. For clarity we include its normalization and show how the same estimates separate it from the rest of the lower spectrum.

\begin{theorem}\label{spec:minimum}
Along the integers for which $M(n)\ne0$,
\begin{equation}\label{spec:minimumformula}
 \sigma_n(B_n)=\frac{|M(n)|}{\sqrt{Q(n)}W_n}(1+o(1)).
\end{equation}
At every zero of $M(n)$, $\sigma_n(B_n)=0$ exactly. In particular,
\[
 \sqrt n\,\sigma_n(B_n)\longrightarrow0
\]
along all integers. If $\kappa_2(B_n)=\sigma_1(B_n)/\sigma_n(B_n)$ denotes the Euclidean condition number at invertible indices, then
\[
 \kappa_2(B_n)\sim\sqrt{c_0}\,\frac{nW_n}{|M(n)|}.
\]
The stable rank, defined by $\operatorname{sr}(B_n)=\norm{B_n}_{\mathrm F}^2/\norm{B_n}^2$, satisfies
\[
 \operatorname{sr}(B_n)\longrightarrow2+c_0,
\]
where $\norm{T}_{\mathrm F}^2$ is the sum of the squared matrix entries.
\end{theorem}

\begin{proof}
At nonzero-Mertens indices, the rank-one term in the exact inverse has operator norm
\[
 \frac{\sqrt{Q(n)}W_n}{|M(n)|}.
\]
The triangle inequality and its reverse show
\[
 \left|\norm{B_n^{-1}}-\frac{\sqrt{Q(n)}W_n}{|M(n)|}\right|
 \le\norm{C_n}=O(\sqrt n).
\]
The relative error is $O(|M(n)|/W_n)=o(1)$ by the same arithmetic comparison used in Theorem~\ref{spec:bottom}. Taking reciprocals proves \eqref{spec:minimumformula}. At a zero of $M(n)$ the exact kernel already proves that the smallest singular value is zero. At the other indices \eqref{spec:minimumformula} and $Q(n)/n\to c_0$ give $\sqrt n\,\sigma_n(B_n)\sim |M(n)|/(\sqrt{c_0}W_n)\to0$.

Combining \eqref{spec:minimumformula} with $\sigma_1(B_n)\sim\sqrt n$ and $Q(n)\sim c_0n$ proves the condition-number assertion. Finally, the row of ones has $n$ entries, the other diagonal entries number $n-1$, and the parent incidences number $Q(n)-1$. These positions are disjoint, so
\[
 \norm{B_n}_{\mathrm F}^2=2n+Q(n)-2.
\]
The upper theorem gives $\norm{B_n}^2/n\to1$, proving the stable-rank limit.
\end{proof}

\begin{corollary}[Approximation of the full inverse]\label{spec:inversefrob}
Along nonzero-Mertens indices,
\[
 \left\|B_n^{-1}+\frac{\muvec_nw_n^T}{M(n)}\right\|_F
   =o\left(\frac{\sqrt{Q(n)}W_n}{|M(n)|}\right),
 \qquad \operatorname{sr}(B_n^{-1})\longrightarrow1.
\]
Up to signs, the right singular vector of $B_n$ for its smallest singular value approaches $\muvec_n/\sqrt{Q(n)}$ in Euclidean norm, and the corresponding left singular vector approaches $w_n/W_n$. The assertion is along the invertible indices and concerns normalized vectors.
\end{corollary}
\begin{proof}
Each squarefree row of $C_n$ has at most $1+\log_2 n$ entries of modulus one, and an isolated row has just one. Thus $\norm{C_n}_F^2\le n(1+\log_2 n)$. The inverse identity makes the difference in the displayed formula exactly $C_n$. Relative to the rank-one term, its Frobenius norm is
\[
 O\left(\frac{|M(n)|\sqrt{\log n}}{W_n}\right)=o(1)
\]
by the same Mertens and norm estimates. Both the operator and Frobenius norms of the inverse are therefore asymptotic to the norm of the rank-one term, proving the stable-rank assertion. After division by that norm, the inverse differs in operator norm by $o(1)$ from a rank-one matrix whose nonzero singular value is one. Its Gram matrix has the corresponding one-dimensional spectral projection plus an $o(1)$ operator error. The variational characterization of the largest eigenvalue then forces the squared inner product of its leading vector with that rank-one direction to tend to one. Apply this to both Gram matrices and interchange left and right vectors on passing from the inverse to $B_n$. This proves the directional assertions and eventual simplicity of the exceptional singular value.
\end{proof}

\section{Products, sums, and volume expansion}\label{sec:volumes}

The endpoint laws do not by themselves describe a product over a growing number of singular values. The adjugate identity supplies a different route, valid even at zero determinants. Throughout this section $n\ge4$ and singular values are those of $B_n$ in decreasing order.

\begin{theorem}[Products after removing the smallest axes]\label{vol:product}
The following finite estimate holds:
\begin{equation}\label{vol:finite}
 \left|\prod_{i=1}^{n-1}\sigma_i-\sqrt{Q(n)}W_n\right|
 \le |M(n)|\norm{C_n}_2.
\end{equation}
Consequently, along all integers $n$,
\begin{equation}\label{vol:main}
 \prod_{i=1}^{n-1}\sigma_i\sim\sqrt{Q(n)}W_n
 =n^{3/2}\exp\bigl(-(1+o(1))s(n)\bigr).
\end{equation}
If $M(n)=0$, the first product in \eqref{vol:main} equals $\sqrt{Q(n)}W_n$ exactly. For every fixed integer $r\ge1$,
\begin{equation}\label{vol:codimension}
 \prod_{i=1}^{n-r}\sigma_i
 \sim\sqrt{Q(n)}W_n\,n^{(r-1)/2}
                  \prod_{j=1}^{r-1}\sqrt{\lambda_j(L)},
\end{equation}
where an empty product equals one and $L$ is the compact operator of Theorem~\ref{spec:bottom}. Also
\begin{equation}\label{vol:middle}
 \prod_{i=2}^{n-1}\sigma_i\sim\sqrt{c_0}\,W_n.
\end{equation}
\end{theorem}
\begin{proof}
For a real $n$ by $n$ matrix $T$, the operator norm of $\adj(T)$ is the product of its largest $n-1$ singular values. For an invertible matrix this follows from $\adj(T)=\det(T)T^{-1}$ and the singular-value decomposition. For a singular matrix it follows by continuity, or directly from the induced map on alternating products of $n-1$ vectors. Equation~\eqref{def:adjugate} and the triangle and reverse triangle inequalities therefore give \eqref{vol:finite}, since
\[
 \norm{\muvec_nw_n^T}_2=\norm{\muvec_n}_2\norm{w_n}_2=\sqrt{Q(n)}W_n.
\]
Theorem~\ref{spec:compact} gives $\norm{C_n}_2=O(\sqrt n)$, while Theorem~\ref{arith:energy} and the classical Mertens bound imply
\[
 \frac{|M(n)|\norm{C_n}_2}{\sqrt{Q(n)}W_n}
 \ll\frac{|M(n)|}{W_n}\longrightarrow0.
\]
This proves \eqref{vol:main} without excluding zero-Mertens indices. At such an index the error term in \eqref{vol:finite} is zero. Divide the product in \eqref{vol:main} by the $r-1$ factors $\sigma_{n-1},\ldots,\sigma_{n-r+1}$ and apply the fixed-index limits in Theorem~\ref{spec:bottom} to obtain \eqref{vol:codimension}. This division is legitimate: the rank-one perturbation of the invertible $A_n$ has rank at least $n-1$. Finally divide \eqref{vol:main} by $\sigma_1\sim\sqrt n$ to prove \eqref{vol:middle}.
\end{proof}

The product of the largest $m$ singular values is the maximal factor by which a linear map expands an $m$-dimensional volume. More explicitly, for an orthonormal $m$-tuple $x_1,\ldots,x_m$, the squared volume of the image parallelotope is the determinant of the Gram matrix $((B_nx_i)^T(B_nx_j))_{ij}$; its maximum is $\prod_{i=1}^m\sigma_i^2$. Thus \eqref{vol:codimension} is a volume law of fixed codimension. It does not follow by multiplying a number of fixed-index equivalents that increases with $n$.

\begin{theorem}[First and second singular-value sums]\label{vol:sums}
For every $n\ge2$,
\begin{equation}\label{vol:second}
 \sum_{i=1}^n\sigma_i^2=2n+Q(n)-2.
\end{equation}
For every fixed $H>0$,
\begin{equation}\label{vol:first}
 \sum_{i=1}^n\sigma_i=n+O_H\left(\frac{n}{(\log n)^H}\right).
\end{equation}
Writing $k=k_n$, the growing upper band satisfies
\begin{equation}\label{vol:band}
 \sum_{i=2}^{k+1}\sigma_i^2\sim c_0n.
\end{equation}
Every fixed number of terms in this band contributes only $O(n/\log n)$.
\end{theorem}
\begin{proof}
The sum of squared singular values is the squared Frobenius norm, so \eqref{vol:second} is \eqref{def:frobenius}. By Theorem~\ref{spec:core}, exactly $n-r_n$ singular values equal one, where $r_n=2k_n+1$. Their complement has squared sum $n+Q(n)-2+r_n$. Cauchy--Schwarz gives
\[
 \sum_i\sigma_i\le n-r_n+
       \sqrt{r_n\bigl(n+Q(n)-2+r_n\bigr)}.
\]
The reverse bound $\sum_i\sigma_i\ge|\tr B_n|=n$ follows directly from a singular-value decomposition: if $B_n=U\Sigma V^T$, then $\tr B_n=\sum_i\sigma_i(V^TU)_{ii}$ and $|(V^TU)_{ii}|\le1$. Apply $k_n=O_{2H}(n/(\log n)^{2H})$ to obtain \eqref{vol:first}.

The $k_n$ singular values below one have total squared sum at most $k_n=o(n)$. The exact unit values contribute $n-2k_n-1=n+o(n)$, and the largest value contributes $\sigma_1^2=n+o(n)$. Subtract these contributions from \eqref{vol:second}; the remaining sum is $Q(n)+o(n)\sim c_0n$. The last assertion follows from Theorem~\ref{spec:upper} for each fixed index.
\end{proof}

In particular the empirical proportion of singular values different from one tends to zero, and their ordinary average tends to one, but their mean square tends to $2+c_0$. This is possible because a small proportion of large values carries a macroscopic squared sum. The full determinant remains
\[
 |M(n)|=\sigma_n\prod_{i=1}^{n-1}\sigma_i.
\]
Combining this identity with the smallest-value equivalent of Theorem~\ref{spec:minimum} only recovers an existing relation. An improved arithmetic estimate would require control of the exceptional thickness that is independent of the Mertens estimate already used to isolate it.

\section{Geometric residuals and prime-threshold components}\label{sec:geometry}

\subsection{Base volume and height}

Let $H_n$ be the $(n-1)$ by $n$ matrix consisting of rows two through $n$ of $A_n$, equivalently of $B_n$. It has full row rank because its minor on columns two through $n$ is unit lower triangular. Proposition~\ref{def:inverse} gives
\[
 H_n\muvec_n=0,\qquad \ker H_n=\spanop\{\muvec_n\}.
\]
The vector of signed maximal minors of $H_n$, normalized to have first coordinate one, is therefore $\muvec_n$. By Cauchy--Binet,
\begin{equation}\label{geo:base}
 \det(H_nH_n^T)=\sum_{j=1}^n\mu(j)^2=Q(n).
\end{equation}
The parallelotope generated by those rows has $(n-1)$-dimensional volume $\sqrt{Q(n)}$. The distance of the remaining row $\one_n^T$ from their span is
\begin{equation}\label{geo:height}
 h_n=\frac{|\one_n^T\muvec_n|}{\norm{\muvec_n}_2}
     =\frac{|M(n)|}{\sqrt{Q(n)}},\qquad
 \min_{z\in\mathbb R^{n-1}}\norm{\one_n-H_n^Tz}_2^2
     =\frac{M(n)^2}{Q(n)}.
\end{equation}
This is an exact height formula, not an estimate for that height.

The symmetric positive semidefinite matrix $G_n^{\mathrm{row}}=H_n^TH_n$ has a useful graph interpretation. On the squarefree coordinates it is the signless Laplacian of the tree: each edge contributes $(e_i+e_j)(e_i+e_j)^T$. Multiplication on both sides by the diagonal matrix with entries $\mu(i)$ on these coordinates changes each edge sum to an edge difference, hence gives the ordinary graph Laplacian. Each nonsquarefree coordinate contributes a one-dimensional identity block.

For $\tau>0$, spectral decomposition yields
\begin{equation}\label{geo:resolvent}
 h_n^2\le\tau\one_n^T(G_n^{\mathrm{row}}+\tau I_n)^{-1}\one_n,
 \qquad
 \lim_{\tau\downarrow0}\tau\one_n^T(G_n^{\mathrm{row}}+\tau I_n)^{-1}\one_n=h_n^2.
\end{equation}
Indeed a spectral coordinate of eigenvalue $\lambda\ge0$ is weighted by $\tau/(\lambda+\tau)$, which equals one on the kernel and tends to zero on each positive eigenvalue. Thus a regularized ellipsoid gives an upper certificate, but merely letting $\tau$ decrease to zero adds no arithmetic information. The weights of $\one_n$ in the eigenspaces matter, not just the eigenvalues.

\subsection{Removing edges by prime size}

Fix a real $y\ge2$. From $H_n$ retain every nonsquarefree diagonal row, and retain a squarefree edge row only if its added prime is greater than $y$. Call the resulting rectangular matrix $H_{n,y}$. Its squarefree components have roots
\[
 \mathcal A(n,y)=\{a\le n:\mu(a)^2=1,\ P^+(a)\le y\}
\]
and vertex sets
\[
 \mathcal C_a(n,y)=\{ad\le n:\mu(d)^2=1,\ P^-(d)>y\}.
\]
To see this, keep the primes at most $y$ in a vertex fixed and delete its larger primes in descending order. The process reaches exactly one root $a$, and reverses along retained edges. The components partition the squarefree vertices.

Set
\begin{equation}\label{geo:F}
 N_a=\Phi_{\sfq}(n/a,y),\quad S_a=\mu(a)R(n/a,y),\quad
 F(n,y)=\sum_{a\in\mathcal A(n,y)}\frac{R(n/a,y)^2}{\Phi_{\sfq}(n/a,y)}.
\end{equation}
Each denominator is a positive component size. Thus $F$ is a sum of squared signed component sums divided by component sizes. It is different from both sums in \eqref{def:energies}.

\begin{proposition}[Projection and variance]\label{geo:projection}
For $n\ge2$ and $y\ge2$,
\begin{align}
 \min_{z\in\mathbb R^{n-|\mathcal A(n,y)|}}\norm{\one_n-H_{n,y}^Tz}_2^2&=F(n,y),\label{geo:residual}\\
 F(n,y)&=\frac{M(n)^2}{Q(n)}+
  \sum_{a\in\mathcal A(n,y)}N_a
       \left(\frac{S_a}{N_a}-\frac{M(n)}{Q(n)}\right)^2.
       \label{geo:variance}
\end{align}
In particular $M(n)^2/Q(n)\le F(n,y)\le Q(n)$. For fixed $n$, $F(n,y)$ is nondecreasing in $y$. For $y\ge\sqrt n$,
\begin{equation}\label{geo:plateaubound}
 0\le Q(n)-F(n,y)\le\frac{4n}{y}.
\end{equation}
\end{proposition}
\begin{proof}
On the component $\mathcal C_a$ the kernel of the retained edge rows is spanned by the alternating vector whose entries are $\mu(ad)=\mu(a)\mu(d)$. Its squared norm is $N_a$, and its inner product with the all-ones vector is $S_a$. Orthogonal projection onto this one-dimensional kernel therefore has squared norm $S_a^2/N_a$. Distinct components have disjoint supports. Nonsquarefree diagonal rows have zero kernel on their coordinates. Summing proves \eqref{geo:residual}.

The partition gives $\sum_aN_a=Q(n)$ and $\sum_aS_a=M(n)$. Expanding the squares in \eqref{geo:variance} proves the equality, including the factor $\mu(a)$ in $S_a$. Omitting that factor would give an incorrect variance identity. The upper bound follows from $|S_a|\le N_a$.

Increasing $y$ deletes edges and splits components. If signed sums and sizes in the pieces are $S_i,N_i$, Cauchy--Schwarz gives $\sum_iS_i^2/N_i\ge(\sum_iS_i)^2/\sum_iN_i$. Thus the energy cannot decrease. If $y\ge\sqrt n$, a rough part is either one or a single prime. A component of size $N$ has alternating sum, up to sign, $2-N$ and energy $(N-2)^2/N=N-4+4/N$. Its deficit from its size is between zero and four. Let $q_1$ be the smallest prime greater than $y$. A nonsingleton root satisfies $a\le n/q_1$, so there are at most $n/q_1$ such components. Summing the deficits proves \eqref{geo:plateaubound}.
\end{proof}

\subsection{A fixed-power profile}

The Dickman function $\rho$ is the continuous function satisfying $\rho(t)=1$ on $0\le t\le1$ and $t\rho'(t)=-\rho(t-1)$ for $t>1$. The Buchstab function $\omega$ is continuous on $[1,\infty)$, equals $1/t$ on $[1,2]$, and satisfies $(t\omega(t))'=\omega(t-1)$ for $t>2$. The symbol $\omega$ in this section denotes this function, not a number of prime factors. Their delay equations determine them successively on unit intervals.

We use the following classical fixed-parameter counting inputs. As $Y\to\infty$, for each fixed $t>0$,
\begin{equation}\label{geo:smoothinput}
 \frac{\Psi(Y^t,Y)}{Y^t}\longrightarrow\rho(t).
\end{equation}
The same limit holds after multiplying $Y^t$ by any fixed positive constant; this follows by monotone bracketing and continuity of $\rho$. Uniformly when $v=\log x/\log y$ belongs to a compact subset of $(1,\infty)$,
\begin{equation}\label{geo:roughinput}
 \Phi(x,y)=\frac{x}{\log y}(\omega(v)+o(1)),\qquad y\to\infty.
\end{equation}
Smooth counting is reviewed in~\cite{hildebrandtenenbaum,mcnew}. For the rough formula on $v\ge2$ see Theorem~B of~\cite{fanpomerance}; on a compact subset of $(1,2]$ the formula follows from $\Phi(x,y)=1+\pi(x)-\pi(y)$ and the prime number theorem. We also use $\Phi(x,y)\ll_V x/\log y$ for $x\ge y$ and $\log x/\log y\le V$, including the endpoint strip. For $1\le v\le2$ this follows from Chebyshev's bound and $1\ll x/\log y$; above two it follows from the rough formula on compact intervals.

\begin{theorem}[Component profile and cutoff]\label{geo:profile}
For each fixed $u\ge1$,
\begin{equation}\label{geo:J}
 \frac{F(n,n^{1/u})}{n}\longrightarrow c_0J(u),\qquad
 J(u)=\rho(u)+\int_1^u\rho(u-v)\frac{\rho'(v)^2}{\omega(v)}\,dv.
\end{equation}
The right derivative of $\rho$ at one may be used at the immaterial endpoint. The function $J$ is positive and nonincreasing, equals one on $[1,2]$, and tends to zero as $u\to\infty$. For $2\le u\le3$,
\begin{equation}\label{geo:explicitJ}
 J(u)=1-4\int_2^u\frac{\log(v-1)}{v(1+\log(v-1))}\,dv,
 \qquad J(2+h)=1-h^2+O(h^3).
\end{equation}
For an arbitrary real sequence $y(n)\ge2$,
\begin{equation}\label{geo:cutoff}
 F(n,y(n))=o(n)\quad\Longleftrightarrow\quad
 \frac{\log y(n)}{\log n}\longrightarrow0.
\end{equation}
The last assertion follows by monotonicity from fixed-$u$ limits; no uniform formula for a growing $u$ is asserted.
\end{theorem}

\begin{proof}
We give the counting and summation details, since the singleton components contribute a term of order $n$ and cannot be dropped.

First define $\Psi_{\sfq}(z,Y)$ to count squarefree $Y$-smooth integers at most $z$. The identity for $\mu^2$ gives
\[
 \Psi_{\sfq}(z,Y)=\sum_{\substack{d\le\sqrt z\\P^+(d)\le Y}}
                         \mu(d)\Psi(z/d^2,Y).
\]
For $z=Y^t$ and each fixed $d$, the normalized summand tends to $\mu(d)\rho(t)/d^2$. The absolute tail beyond $d>D$ is at most $\sum_{d>D}d^{-2}$. Therefore
\begin{equation}\label{geo:sfsmooth}
 \Psi_{\sfq}(Y^t,Y)/Y^t\longrightarrow c_0\rho(t).
\end{equation}
Place mass $1/(a\log Y)$ at $\log a/\log Y$ for each squarefree $Y$-smooth integer $a$, and call this measure $\nu_Y$. Partial summation shows, for fixed $t>0$,
\[
 \nu_Y([0,t])=
 \frac{\Psi_{\sfq}(Y^t,Y)}{Y^t\log Y}
 +\frac1{\log Y}\int_1^{Y^t}\frac{\Psi_{\sfq}(z,Y)}{z^2}\,dz
 \longrightarrow c_0\int_0^t\rho(w)\,dw.
\]
For the integral substitute $z=Y^w$ and use \eqref{geo:sfsmooth} and domination by one. Thus on every compact interval $\nu_Y$ converges weakly to the measure with density $c_0\rho(w)$. The atom at $a=1$ vanishes because its mass is $1/\log Y$.

Second, we derive the signed counterpart of \eqref{geo:roughinput} in the fixed-parameter range:
\begin{equation}\label{geo:signedfixed}
 R(x,y)=\frac{x}{\log y}(\rho'(v)+o(1))
\end{equation}
uniformly on compact subsets of $1<v<\infty$. For $1+\eta\le v\le2$, there are only one and prime rough terms, so $R=1-\pi(x)+\pi(y)$ and the limiting coefficient is $-1/v$. For larger $v$, grouping terms by their least prime factor gives the exact identity
\[
 R(x,y)=R(x,\sqrt x)-\sum_{y<p\le\sqrt x}R(x/p,p).
\]
Proceed by induction over a bounded range of $v$. The inner parameter $w=\log(x/p)/\log p$ is at most $v-1$. In its strip $1\le w<1+\epsilon$ use $|R|\le\Phi\ll x/(p\log p)$. Chebyshev and partial summation, or the prime number theorem, bound the strip by $O(\epsilon x/\log x)+o(x/\log x)$. Off this strip the inductively uniform limit may be summed, since the total of the weights $x/(p\log p)$ is $O_V(x/\log x)$. Prime partial summation gives
\[
 \frac{x}{\log x}\int_1^{v-1}r(w)\,dw
\]
for the sum's leading term, where $r$ is the previously determined coefficient. All primes in these bounded parameter ranges tend to infinity. Since $R(x,\sqrt x)\sim-x/\log x$, the recurrence for the coefficient is
\[
 vr(v)=-1-\int_1^{v-1}r(w)\,dw.
\]
It agrees with $r(v)=\rho'(v)$ by the Dickman equation and the base interval. The recurrence determines each next unit interval, proving \eqref{geo:signedfixed}; first take $x\to\infty$, then $\epsilon\downarrow0$ to remove the strip.

Nonsquarefree rough integers contribute at most $x\sum_{p>y}p^{-2}=O(x/y)$. Hence \eqref{geo:roughinput} also holds for $\Phi_{\sfq}$, with the same positive coefficient. Put $g(v)=\rho'(v)^2/\omega(v)$. We have proved, uniformly on the same compact subsets,
\begin{equation}\label{geo:componentasymptotic}
 \frac{R(x,y)^2}{\Phi_{\sfq}(x,y)}
       =\frac{x}{\log y}(g(v)+o(1)).
\end{equation}

Now fix $u>1$ and put $y=n^{1/u}$. Roots $a>n/y$ have $n/a<y$, so their components are singletons and contribute
\[
 \Psi_{\sfq}(n,y)-\Psi_{\sfq}(n/y,y)=c_0n\rho(u)+o(n).
\]
For the other roots write $v=u-\log a/\log y$. On $v\ge1+\epsilon$, apply \eqref{geo:componentasymptotic} and the weak limit of $\nu_y$ to obtain, after division by $n$,
\[
 \frac1{\log y}\sum_{\substack{a\in\mathcal A(n,y)\\a\le n/y^{1+\epsilon}}}
       \frac{g(u-\log a/\log y)}a
 \longrightarrow c_0\int_{1+\epsilon}^u\rho(u-v)g(v)\,dv.
\]
The summed uniform remainder is $o(1)$ because $(\log y)^{-1}\sum_{a\le n}a^{-1}=O_u(1)$. Endpoint discontinuities in the cutoff cause no problem because the limiting measure has no atoms. In the omitted strip $1\le v<1+\epsilon$, the bound
\[
 R(x,y)^2/\Phi_{\sfq}(x,y)\le\Phi_{\sfq}(x,y)\le\Phi(x,y)\ll x/\log y
\]
gives $O(\epsilon n)+o(n)$ after summing over $y^{u-1-\epsilon}<a\le y^{u-1}$. Choose $0<\epsilon<u-1$; the harmonic sum of this interval is at most $\epsilon\log y+o(1)$. Let $\epsilon\downarrow0$ to prove \eqref{geo:J}. For $u=1$, every component is a singleton and the formula is immediate.

Positivity of $J$ follows from $\rho(u)>0$. Monotonicity follows from Proposition~\ref{geo:projection}, since $n^{1/u}$ decreases with $u$. For completeness the delay equation gives $u\rho(u)=\int_{u-1}^u\rho(t)\,dt$, which prevents a first zero of $\rho$. It also gives monotone decrease and limit zero: a positive limiting value would contradict integrating $-\rho'(u)\ge c/u$. The Buchstab equation gives $\omega(v)\ge1/2$ on $[1,\infty)$ by induction from $v\omega(v)=1+\int_2^v\omega(t-1)\,dt$. Thus $g(v)\le2/v^2$. Dominated convergence in the integral in \eqref{geo:J} proves $J(u)\to0$.

Repeat the component summation with the masses $\Phi_{\sfq}$ instead of squared sums divided by masses. Their total is $Q(n)$, so
\begin{equation}\label{geo:partitionlimit}
 1=\rho(u)+\int_1^u\rho(u-v)\omega(v)\,dv.
\end{equation}
For $1\le v\le2$, $g(v)=\omega(v)=1/v$, giving the plateau. For $2\le u\le3$, only $2\le v\le u$ contributes to the difference between \eqref{geo:partitionlimit} and \eqref{geo:J}, and $\rho(u-v)=1$ there. Directly from the delay equations,
\[
 \rho'(v)=\frac{\log(v-1)-1}{v},\quad
 \omega(v)=\frac{1+\log(v-1)}v,\quad
 \omega(v)-g(v)=\frac{4\log(v-1)}{v(1+\log(v-1))}.
\]
This proves \eqref{geo:explicitJ}; its integrand including the factor four is $2(v-2)+O((v-2)^2)$ at two.

Finally, if $\log y(n)/\log n\to0$, for each fixed $U\ge1$ the cutoff is eventually at most $n^{1/U}$. Monotonicity gives $\limsup F(n,y(n))/n\le c_0J(U)$, which tends to zero as $U\to\infty$. If that logarithmic ratio does not tend to zero, some subsequence satisfies $y(n)\ge n^\delta$ for a fixed $0<\delta<1$. Along it the lower bound tends to $c_0J(1/\delta)>0$. This proves both directions of \eqref{geo:cutoff}.
\end{proof}

\begin{proposition}[Fixed thresholds and elementary matching]\label{geo:limitations}
For each fixed $y\ge2$, $F(n,y)=o(n)$ is equivalent to $M(n)=o(n)$. A matching consisting of retained squarefree tree edges contains at most $\lfloor n/q_1\rfloor$ edges, where $q_1$ is the least prime greater than $y$. Consequently such a matching leaves at least $Q(n)-2\lfloor n/q_1\rfloor$ squarefree vertices uncovered.
\end{proposition}
\begin{proof}
One direction of the equivalence follows from \eqref{geo:variance} and $Q(n)\le n$. Conversely, the finite convolution
\[
 R(x,y)=\sum_{\substack{b\le x\\P^+(b)\le y}}M(x/b)
\]
and $\sum_{P^+(b)\le y}b^{-1}=\prod_{p\le y}(1-1/p)^{-1}<\infty$ imply $R(x,y)=o(x)$ by dominated convergence if $M(x)=o(x)$. For fixed $y$ there are only finitely many squarefree smooth roots. Each denominator is asymptotic to a positive constant times $n/a$, by elementary squarefree density with finitely many excluded primes. Each contribution is therefore $o(n)$ and so is their finite sum. In a matching the smaller endpoints are distinct. Each lies at most $n/q_1$, because its edge adds a prime at least $q_1$. This bounds the number of edges and proves the uncovered-vertex assertion.
\end{proof}

The matching bound limits this particular pairing certificate; it does not lower-bound $F$ itself. More generally a star has the same sparse edge-row form and a single alternating kernel vector, yet its alternating sum can have linear size. Volume, sparsity, and nested tree structure alone do not supply cancellation. The arithmetic inputs used in the component profile and in the following appendix must remain visible.

\section{Scope, applications, and remaining questions}\label{sec:discussion}

The explicit reducing core can replace the full matrix when computing nonunit singular values. The upper comparison identifies a regime of large singular values directly from prime counts; the compact operator identifies the fixed lower constants. The near-rank-one inverse describes the direction of ill-conditioning, and the fixed-codimension products describe maximal volume expansion even at zero determinants. The support-constrained arithmetic estimates apply to arbitrary thin sets of least prime factors without a distribution assumption on those sets. These are the immediate mathematical applications established here. No complexity improvement for a prime-counting algorithm, new factoring method, or practical error threshold is claimed.

\subsection{What is needed for a stronger Mertens bound}\label{disc:mertens-improvement}

The issue is not another representation of $M(n)$, but an estimate for one of its representations that uses additional cancellation. The classical comparison in this paper is \eqref{arith:mertens}: with
\[
 V(n)=\frac{(\log n)^{3/5}}{(\log\log n)^{1/5}},
\]
it gives $|M(n)|\le Cn\exp(-cV(n))$ for suitable positive constants and all sufficiently large $n$. A new proof of $M(n)=o(n)$, or of a weaker quantitative bound, would not improve this estimate. An explicit improvement must specify its constants and range; a stronger asymptotic scale must specify its rate. For example, proving $|M(n)|\le Cn\exp(-a(n))$ with $a(n)/V(n)\to\infty$ would improve the displayed scale. No such result is established here.

There is a concrete sufficient condition that does not require computing a determinant. Let $\varepsilon(n)>0$ be a proposed error rate. Construct a vector $z_n\in\mathbb R^{n-1}$ and bound its residual by
\begin{equation}\label{disc:certificate}
 \norm{\one_n-H_n^Tz_n}_2^2
 \le \frac{n^2\varepsilon(n)^2}{Q(n)}.
\end{equation}
Then $|M(n)|\le n\varepsilon(n)$. Indeed, $H_n\muvec_n=0$ implies
\[
 M(n)=\muvec_n^T(\one_n-H_n^Tz_n),
\]
and Cauchy--Schwarz gives the conclusion. Thus the missing ingredient could be an explicit choice of row coefficients together with a uniform residual estimate. Choosing the exact least-squares solution and using \eqref{geo:height} to evaluate its error merely recovers $M(n)^2/Q(n)$; it supplies no independent cancellation estimate. Since $Q(n)\sim c_0n$, the required squared residual has scale $n\varepsilon(n)^2/c_0$, with constants controlled if an explicit bound is intended.

Two other sufficient certificates are already available from the geometric identities. For a choice $y=y(n)\ge2$, it suffices to prove
\begin{equation}\label{disc:component-certificate}
 F(n,y)\le\frac{n^2\varepsilon(n)^2}{Q(n)}.
\end{equation}
Alternatively, for a choice $\tau=\tau(n)>0$, it suffices to prove
\begin{equation}\label{disc:resolvent-certificate}
 \tau\one_n^T(H_n^TH_n+\tau I_n)^{-1}\one_n
 \le\frac{n^2\varepsilon(n)^2}{Q(n)}.
\end{equation}
The first follows from \eqref{geo:variance}, the second from \eqref{geo:resolvent}. Each is a target for further work, not an estimate proved by writing the formula. In particular, the component variance identity shows that $F$ includes both the global Mertens contribution and nonnegative variation between components. A useful upper bound must control both. The fixed-power profile has $F(n,n^{1/u})\sim c_0nJ(u)$ with $J(u)>0$ for every fixed $u\ge1$; that regime cannot certify $\varepsilon(n)\to0$. Obtaining an explicit decay rate as $u$ grows requires additional quantitative control; the qualitative cutoff in Theorem~\ref{geo:profile} follows from the fixed-parameter limits and monotonicity. The displayed resolvent depends on the weights of $\one_n$ in the eigenspaces; the eigenvalue information established here alone does not control its smallness.

The singular-value approach has the same unresolved step. The formula for the exceptional smallest singular value already contains $|M(n)|$. Combining that formula with the product of the other singular values therefore reproduces $|M(n)|$, rather than bounding it independently. To gain arithmetic information, one needs an independently obtained estimate for the exceptional direction, or uniform information on a growing part of the spectrum that controls the full product without importing the desired Mertens bound. Fixed-index limits, the sum of squared singular values, and the product with the smallest factor omitted do not provide that missing estimate.

Any proposed improvement should therefore state three things separately: the new arithmetic or geometric input; the uniform argument transferring it to one of these certificates; and the resulting rate for $M(n)$ compared with \eqref{arith:mertens}. Bounds obtained by inserting the existing Mertens estimate into restricted sums are useful consequences, but cannot themselves serve as independent evidence for stronger cancellation.

\subsection{Further questions and release obligations}

Several natural refinements remain open within this program. The prime moments are determined on their leading exponential scale; a multiplicative equivalent would require a finer saddle analysis and a controlled prime-summation error. A local limit theorem would also require locating the saddle beyond its leading term before centering any Gaussian approximation. The upper-spectrum comparison is uniform, but the explicit limiting constants at both ends have been used only at fixed indices. A growing-index description that controls a logarithmic product across the intermediate band would give information absent from the endpoint laws alone.

The elementary estimate for $k_n$ proves that almost all singular values equal one, but does not give a sharp asymptotic for the size of the remaining core. The operator constants have exact definitions, but a numerical value would need a certified truncation bound and eigenvalue enclosure, not just a finite power iteration. Finally, the geometric residual could lead to an arithmetic improvement only if it were bounded by a mechanism that does not already assume comparable M\"obius cancellation.

The comparison with Alladi~\cite{alladi1982} shows that his theorems stop short of the saddle range used here. Alladi's companion paper~\cite{alladi1982b} and Tenenbaum~\cite{tenenbaum1990} could not be read for this release. Whether Lemma~\ref{arith:bridge}, or its moment consequences, occurs there or in later literature on restricted M\"obius sums remains open.

\paragraph{Provenance.}
AI assisted the research development, literature search, proof exposition, and mathematical checking under the author's direction. Process-separated AI reviews are not independent expert peer review, and no independent human expert has examined the proofs. The release repository records the inspected sources, the proof audits, and their dispositions.

\appendix
\section{An optional quantitative component bound from Hal\'asz's theorem}\label{app:halasz}

This appendix gives a complete application of a known mean-value theorem to the actual component residual $F(n,y)$. It is independent of the Mertens estimate used for the sharper spectral scales, but does not supply a new general cancellation theorem or improve known Mertens bounds.

A function $f:\mathbb N\to\mathbb C$ is multiplicative if $f(ab)=f(a)f(b)$ whenever $a,b$ are coprime, and is one-bounded if $|f(m)|\le1$. For such a function define the nonnegative quantity
\[
 \mathcal D(f,t;x)^2=\sum_{p\le x}\frac{1-\Re(f(p)p^{-it})}{p}.
\]
We use the following weakened quantitative form of Hal\'asz's theorem: there are absolute $C,c_H>0$ such that, for sufficiently large $x$ and $T>0$,
\begin{equation}\label{app:H}
 \frac1x\left|\sum_{m\le x}f(m)\right|
 \le C\left(\exp\left(-c_H\min_{|t|\le T}\mathcal D(f,t;x)^2\right)+\frac1T\right).
\end{equation}
This is the form in~\cite[Theorem~7]{taohalasz}; a sharper primary theorem is proved in~\cite{ghs}. We import \eqref{app:H}, rather than purporting to reprove it. The elementary prime estimates used below are
\[
 \sum_{p\le z}\frac1p=\log\log z+O(1),\quad
 \sum_{p\le z}\frac{\log p}{p}\ll\log z,\quad
 \pi(z)\ll\frac z{\log z},\quad
 \prod_{p\le z}(1-1/p)\asymp\frac1{\log z}.
\]
These Chebyshev--Mertens estimates do not require the prime number theorem.

\begin{proposition}\label{app:quantitative}
Let $\delta=\min(c_H/16,1/4)>0$. For every fixed $0<\alpha<1$,
\begin{equation}\label{app:rate}
 F(n,(\log n)^\alpha)
 \ll_\alpha\frac{n(\log\log n)^2}{(\log n)^{2\delta}}.
\end{equation}
Also
\begin{equation}\label{app:slowrate}
 F(n,\log\log n)\ll n\left(\frac{\log\log\log n}{\log\log n}\right)^2.
\end{equation}
The constants and onsets are not made numerically explicit.
\end{proposition}
\begin{proof}
Write $\ell=\log x$ and $\sigma=1+1/\ell$. Uniformly for real $v$, the absolutely convergent Euler product and the elementary prime estimates give
\begin{equation}\label{app:phase}
 \sum_{p\le x}\frac{\cos(v\log p)}p
          =\log|\zeta(\sigma+iv)|+O(1).
\end{equation}
Indeed, prime powers of exponent at least two contribute $O(1)$; replacing $p^{-\sigma}$ by $p^{-1}$ below $x$ costs at most $\ell^{-1}\sum_{p\le x}(\log p)/p=O(1)$; and the tail above $x$ is bounded by a constant times $\int_x^\infty t^{-1-1/\ell}(\log t)^{-1}dt=O(1)$. Each error is independent of $v$.

For $1<\sigma\le2$, comparison of the defining series with its integral gives
\[
 \zeta(z)=\frac z{z-1}-z\int_1^\infty\{t\}t^{-z-1}\,dt,
 \qquad |\zeta(\sigma+iv)|\ll1+|v|+|v|^{-1}\quad(v\ne0),
\]
where $\{t\}=t-\lfloor t\rfloor$. This calculation stays in the half-plane of absolute convergence.

Put $\mathcal E(t;x)=\sum_{p\le x}(1+\cos(t\log p))/p$. For $|t|\le\ell^{-1/2}$, retain only $p\le\exp(\sqrt\ell)$; their angles have absolute value at most one, so
\[
 \mathcal E(t;x)\ge\frac{1+\cos1}{2}\log\ell-O(1).
\]
For $\ell^{-1/2}<|t|\le T\le\sqrt\ell$, the inequality $1-\cos(2v)\le4(1+\cos v)$ and \eqref{app:phase} give
\[
 4\mathcal E(t;x)\ge\log\ell-\log|\zeta(\sigma+2it)|-O(1)
                 \ge\tfrac12\log\ell-O(1).
\]
Therefore, uniformly for $1\le T\le\sqrt{\log x}$,
\begin{equation}\label{app:distance}
 \min_{|t|\le T}\mathcal E(t;x)\ge\tfrac18\log\log x-O(1).
\end{equation}

Define $f_y(m)=\mu(m)$ when $P^-(m)>y$, and zero otherwise. It is multiplicative and one-bounded, with prime values zero below or at $y$ and minus one above it. For $x\ge y$, \eqref{app:distance} yields
\[
 \mathcal D(f_y,t;x)^2=\mathcal E(t;x)-\sum_{p\le y}\frac{\cos(t\log p)}p
 \ge\tfrac18\log\log x-\log\log y-O(1).
\]
Take $y=(\log n)^\alpha$ and let $P_y=\prod_{p\le y}p$. Every squarefree smooth root divides $P_y$. The elementary bound $P_y\le y^y=n^{o(1)}$ shows that every divisor occurs as a root and $x=n/a\ge\sqrt n$ uniformly for large $n$. Thus $\log\log x=\log\log n+O(1)$ and $\log\log y=o(\log\log x)$. Taking $T=(\log n)^{1/4}$ in \eqref{app:H} is admissible, and gives
\begin{equation}\label{app:Rbound}
 |R(x,y)|/x\ll(\log n)^{-\delta}
\end{equation}
uniformly over every component.

For the denominator, inclusion--exclusion and removal of square divisors give
\[
 \Phi_{\sfq}(x,y)\ge x\prod_{p\le y}(1-1/p)-2^{\pi(y)}-O(x/y).
\]
Since $2^{\pi(y)}\le2^y=n^{o(1)}$ and $x\ge\sqrt n$, both error terms are $o(x/\log y)$. Consequently $\Phi_{\sfq}(x,y)\gg x/\log y$. Substituting \eqref{app:Rbound} into the definition of $F$ gives
\[
 F(n,y)\ll\frac{n\log y}{(\log n)^{2\delta}}
             \sum_{a\mid P_y}\frac1a
 \ll\frac{n(\log y)^2}{(\log n)^{2\delta}},
\]
because $\prod_{p\le y}(1+1/p)\ll\log y$. This proves \eqref{app:rate}.

For $y=\log\log n$, use the same root and denominator arguments, but take $T=\log\log n$. The exponential-distance term is smaller than $1/\log\log n$ eventually, so $|R(x,y)|/x\ll1/\log\log n$. The same final summation proves \eqref{app:slowrate}.
\end{proof}

Combining \eqref{app:rate} with $M(n)^2\le Q(n)F(n,y)$ yields a weak classical cancellation bound. The saving has already entered through \eqref{app:H}; the projection does not independently generate it. In particular this appendix cannot be used to claim a purely geometric proof of the prime number theorem.

\begin{thebibliography}{99}

\bibitem{alladi1982}
K.~Alladi, \emph{Asymptotic estimates of sums involving the Moebius function},
Journal of Number Theory \textbf{14} (1982), 86--98.
\href{https://doi.org/10.1016/0022-314X(82)90060-9}{doi:10.1016/0022-314X(82)90060-9}.

\bibitem{fanpomerance}
K.~(Steve) Fan and C.~Pomerance, \emph{An inequality related to the sieve of Eratosthenes},
Journal of Number Theory \textbf{254} (2024), 169--183.
\href{https://doi.org/10.1016/j.jnt.2023.07.005}{doi:10.1016/j.jnt.2023.07.005}, Theorem~B.

\bibitem{ghs}
A.~Granville, A.~J.~Harper, and K.~Soundararajan,
\emph{A more intuitive proof of a sharp version of Hal\'asz's theorem},
\href{https://arxiv.org/abs/1706.03755}{arXiv:1706.03755} (2017).

\bibitem{hilberdink}
T.~Hilberdink, \emph{Singular values of multiplicative Toeplitz matrices},
Linear and Multilinear Algebra \textbf{65} (2017), 813--829.
\href{https://doi.org/10.1080/03081087.2016.1204978}{doi:10.1080/03081087.2016.1204978}.
\href{https://centaur.reading.ac.uk/66059/1/finitetoeplitz.pdf}{Accepted manuscript}.

\bibitem{hildebrandtenenbaum}
A.~Hildebrand and G.~Tenenbaum, \emph{Integers without large prime factors},
Journal de Th\'eorie des Nombres de Bordeaux \textbf{5} (1993), 411--484.

\bibitem{kline2019}
J.~Kline, \emph{A sparser matrix representation of the Mertens function},
Linear Algebra and its Applications \textbf{581} (2019), 354--366.
\href{https://doi.org/10.1016/j.laa.2019.07.021}{doi:10.1016/j.laa.2019.07.021}.

\bibitem{kline2020}
J.~Kline, \emph{On the eigenstructure of sparse matrices related to the prime number theorem},
Linear Algebra and its Applications \textbf{584} (2020), 409--430.
\href{https://doi.org/10.1016/j.laa.2019.09.022}{doi:10.1016/j.laa.2019.09.022}.

\bibitem{klineDraft}
J.~Kline, \emph{The smallest singular value of a sparse Mertens matrix},
working manuscript, local predecessor inspected September 27, 2026.
Project locator: \url{https://github.com/jeff-kline/sparse-mertens-singular-values}.
All identities needed from this predecessor are proved in the present paper.

\bibitem{klineExtremal}
J.~Kline, \emph{Extremal eigenvalues and singular values of sparse matrices related to the prime number theorem},
version 0.1.0 (2026).
\href{https://doi.org/10.5281/zenodo.21764114}{doi:10.5281/zenodo.21764114}.
Living source: \url{https://github.com/jeff-kline/extremal-eigenvalues}.

\bibitem{kms}
M.~Kural, V.~McDonald, and A.~Sah,
\emph{M\"obius formulas for densities of sets of prime ideals},
Archiv der Mathematik \textbf{115} (2020), 53--66.
\href{https://arxiv.org/abs/1907.02914}{arXiv:1907.02914}.

\bibitem{leeleong}
E.~S.~Lee and N.~Leong,
\emph{New explicit bounds for Mertens function and the reciprocal of the Riemann zeta function},
\href{https://arxiv.org/html/2208.06141v5}{arXiv:2208.06141v5} (September 9, 2026),
Theorem~1.1, equation~(12).

\bibitem{mcnew}
N.~McNew, \emph{The most frequent values of the largest prime divisor function},
Experimental Mathematics \textbf{26} (2017), 210--224.
\href{https://www.nathanmcnew.com/PopularPrimes.pdf}{Author manuscript},
equations~(7), (13), (14), and (17).

\bibitem{saias}
E.~Saias, \emph{Sur le nombre des entiers sans grand facteur premier},
Journal of Number Theory \textbf{32} (1989), 78--99.
\href{https://doi.org/10.1016/0022-314X(89)90099-1}{doi:10.1016/0022-314X(89)90099-1}.
The expansion used here was checked in the explicit formulation quoted by McNew.

\bibitem{taohalasz}
T.~Tao, \emph{254A, Notes 10---Mean values of nonpretentious multiplicative functions},
December 17, 2019, Theorem~7.
\href{https://terrytao.wordpress.com/2019/12/17/254a-notes-10-mean-values-of-nonpretentious-multiplicative-functions/}{Author's lecture notes}.

\end{thebibliography}

\end{document}
